%% 
%% Copyright 2007-2020 Elsevier Ltd
%% 
%% This file is part of the 'Elsarticle Bundle'.
%% ---------------------------------------------
%% 
%% It may be distributed under the conditions of the LaTeX Project Public
%% License, either version 1.2 of this license or (at your option) any
%% later version.  The latest version of this license is in
%%    http://www.latex-project.org/lppl.txt
%% and version 1.2 or later is part of all distributions of LaTeX
%% version 1999/12/01 or later.
%% 
%% The list of all files belonging to the 'Elsarticle Bundle' is
%% given in the file `manifest.txt'.
%% 

%

\documentclass[a4paper,reqno,11pt]{amsart}


\setlength{\textwidth}{\paperwidth}
\addtolength{\textwidth}{-5cm}
\calclayout

%% The amssymb package provides various useful mathematical symbols
\usepackage{amssymb}
\usepackage{amsmath}
%\usepackage{imakeidx}
\usepackage{amsthm}
%\usepackage{titleps}
%\usepackage{verbatim}
\usepackage{enumitem}
\usepackage[table]{xcolor}
\usepackage{latexsym, mathrsfs, graphicx, fancyhdr,bm}
\usepackage{multirow}
\usepackage[final]{listings}
\usepackage{longtable}
\usepackage[leftmargin=3em]{quoting}
\usepackage[hidelinks]{hyperref}

%\usepackage{showlabels}

%% The amsthm package provides extended theorem environments
%% \usepackage{amsthm}

%% The lineno packages adds line numbers. Start line numbering with
%% \begin{linenumbers}, end it with \end{linenumbers}. Or switch it on
%% for the whole article with \linenumbers.
%\usepackage[pagewise]{lineno}
%\linenumbers


\newtheorem{mainthA}{Theorem}
\renewcommand{\themainthA}{\rm{A}\arabic{mainthA}}



\newtheorem{mainthC}{Theorem}
\renewcommand{\themainthC}{\rm{C}\arabic{mainthC}}

\newtheorem*{mmainth}{Main Theorem}

\newtheorem*{mainthAA}{Theorem A}

%%%

\newtheorem{Th}{Theorem}[section]
\newtheorem{Lem}[Th]{Lemma}
\newtheorem{Prob}[Th]{Problem}
\newtheorem{Cor}[Th]{Corollary}
\newtheorem{Que}[Th]{Question}
\newtheorem{Prop}[Th]{Proposition}
\newtheorem*{T1}{Theorem~\ref{theorem}}
\newtheorem*{T2}{Theorem~\ref{theoremGU}}
\newtheorem*{T3}{Theorem~\ref{theoremSp}}
\newtheorem*{T4}{Theorem~\ref{theoremGR}}
\newtheorem*{T5}{Theorem~\ref{theoremSpGR}}

\theoremstyle{definition}
\newtheorem{Def}[Th]{Definition}
\newtheorem*{Def*}{Definition}
\newtheorem{example}[Th]{Example}
\newtheorem{xca}[Th]{Exercise}
\newtheorem{Not}[Th]{Notation}

\theoremstyle{remark}
\newtheorem{Rem}[Th]{Remark}

%\numberwithin{section}{section}
\numberwithin{equation}{section}


\newcommand{\fpr}{{\mathrm{fpr}}}
\newcommand{\Sin}{\mathrm{Sin}}
\newcommand{\Sym}{\mathrm{Sym}}
\newcommand{\Fix}{\mathrm{Fix}}
\newcommand{\Gal}{\mathrm{Gal}}
\newcommand{\diag}{\mathrm{diag}}
\newcommand{\soc}{\mathrm{soc}}
\newcommand{\Stab}{\mathrm{Stab}}
\newcommand{\sgn}{\mathrm{sgn}}
\newcommand{\per}{\mathrm{perm}}
\newcommand{\Det}{\mathrm{Det}}
\newcommand{\Reg}{\mathrm{Reg}}
\newcommand{\Aut}{\mathrm{Aut}}

\newcommand{\GaL}{{\Gamma \mathrm{L}}}
\newcommand{\PGaL}{{\mathrm{P}\Gamma \mathrm{L}}}
\newcommand{\SL}{{\mathrm{SL}}}
\newcommand{\PSL}{{\mathrm{PSL}}}
\newcommand{\GL}{{\mathrm{GL}}}
\newcommand{\PGL}{{\mathrm{PGL}}}
\newcommand{\PGU}{{\mathrm{PGU}}}
\newcommand{\Sp}{{\mathrm{Sp}}}
\newcommand{\PSp}{{\mathrm{PSp}}}


\newcommand*{\Scale}[2][4]{\scalebox{#1}{$#2$}}

\setlist[enumerate,1]{label={(\arabic*)}}




\title{Base sizes for finite linear groups with solvable stabilisers}


\author{Anton A. Baykalov}
\email{a.a.baykalov@gmail.com}

\address{Department of Mathematics and Statistics, The University of Western Australia, Perth, Australia}

\subjclass[2010]{Primary 20D06;
Secondary 20D60}


\keywords{finite groups, simple groups, solvable groups, base size}







\begin{document}


\begin{abstract}
Let $G$ be a transitive permutation group on a finite set with solvable point stabiliser. In 2010, Vdovin  conjectured that the base size of $G$ is at most 5. Burness proved this conjecture in the case of primitive $G$. The problem was reduced by Vdovin in 2012 to the case when $G$ is an almost simple group. Now the problem is further reduced to groups of Lie type through work of Baykalov and Burness. In this paper, we prove the strong form of the conjecture for all almost simple groups with socle isomorphic to $\PSL_n(q).$  
 In the upcoming second paper, the analysis is extended to the unitary and symplectic groups. The final third paper on orthogonal groups is in preparation. Together, these three paper will complete the proof of Vdovin's conjecture for all almost simple classical groups. 
\end{abstract}






\maketitle



\section{Introduction}
%\label{}

\subsection{Main problem and results}

Consider a permutation group $G$ on a finite set $\Omega.$ A { \bf base } of $G$ is a subset of $\Omega$ such that its pointwise stabiliser is trivial. The {\bf base size} of $G$ is the minimal size of a base. The study of bases and base sizes of permutation groups is an active research area with a rich history. These concepts have applications in abstract group theory and computational algebra; for a summary, see \cite{bailey} and  \cite{sfa}. In the case of a transitive group $G$ with point stabiliser $H$, we write $b_H(G)$ for the base size, noting that the action of $G$ on $\Omega$  is permutation isomorphic to the action of $G$ on right $H$-cosets by right multiplication.  


In the past few decades, significant progress has been made in determining the base sizes of groups, with a particular focus on the primitive actions of almost simple groups. Liebeck and Shalev \cite{lieb} proved the following conjecture of Cameron and Kantor \cite{Camer}:  if $G$ is an almost simple finite group and $H\le G$ is maximal and core-free, then there exists an absolute constant $c$ such that $b_H(G)$  is at most   $c$ unless $(G,H)$ lies in a prescribed list of exceptions. The exceptions arise when the action of $G$ on the set of right cosets of $H$  is {\bf standard} (see \cite[Definition 1.1]{burness}). Results specifying bounds for $b_H(G)$ include \cite{burness}, \cite{bls}, \cite{spor}, \cite{mon} and \cite{sfa}. In particular, if $G$ is an almost simple group in a primitive non-standard action, then $b_H(G) \le 7$ with equality if and only if $G=\mathrm{M}_{24}$ and $H=\mathrm{M}_{23}.$

This paper is the first in a series of three devoted to the study of  base sizes of finite transitive  permutation groups with solvable point stabilisers.  In \cite[Problem 17.41 b)]{kt}, Vdovin conjectures (in a slightly different notation) that $b_S(G) \le 5$ for a transitive permutation group $G$ with a solvable point stabiliser $S$. Burness \cite{burPS} proved this conjecture in the case of primitive $G$. It is worth mentioning that Seress \cite{seress} established $b_S(G)\le 4$ in the case of solvable primitive $G$, so the maximal subgroup $S$, of course, is solvable as well. Notice that the bound in the conjecture is the best possible since $b_S(G)=5$ if $G=\Sym(8)$ and $S=\Sym(4) \wr \Sym(2).$ This can be easily verified. In fact, there are infinitely many examples with $b_S(G)=5$, for example see \cite[Remark 8.3]{burPS}.

 Vdovin reduces the above conjecture to the case of almost simple group $G$ in \cite[Theorem 1]{vd}. In order to explain how exactly the reduction works, let us define $\Reg(G,k)$ to be the number of distinct regular orbits in the  action of a permutation group $G \le \Sym(\Omega)$ on $\Omega^k$; here $G$ acts on $\Omega^k$ by $$(\alpha_1, \ldots,\alpha_k)g = (\alpha_1g,\ldots,\alpha_kg).$$ Analogous to $b_H(G)$, we write $\Reg_H(G,k)$ for $\Reg(G,k)$ in the case when $\Omega$ is the set of right cosets of a subgroup $H$. Since a regular point in $\Omega^k$ forms a base for $G$ acting on $\Omega$, $\Reg_H(G,k)=0$ for $k < b_H(G).$    Now we return to the reduction: \cite[Theorem 1]{vd} implies that, in order to prove Vdovin's conjecture,   
it is sufficient to show $$\Reg_S(G,5) \ge 5$$ for every almost simple group $G$ and each of its maximal solvable subgroups $S$; that is, $S$ is not contained in any larger solvable subgroup of $G$. This inequality is established for almost simple groups with alternating and sporadic socle in \cite{bay} and \cite{burspor} respectively. Therefore, Vdovin's conjecture is now reduced to the case of almost simple group of Lie type. This paper is the first one to consider the general form of the conjecture for groups of Lie type. Here we  consider groups with socle isomorphic to $\PSL_n(q)$ and prove the following. 

  \begin{Th}
\label{fulltheoremPSL}
Let $\Omega$ be a finite set and let $G\le \Sym(\Omega)$ be a transitive almost simple group with socle isomorphic to $\PSL_n(q)$. If the stabiliser $S$ in $G$ of a point of $\Omega$ is solvable, then $\Reg_S(G,5)\ge 5$ and thus $b_S(G) \le 5$.
\end{Th}

\begin{Rem}
 Note that the formulation of Theorem 1  in \cite{vd} is incorrect and this is fixed in the latest arXiv version (see the link in the Bibliography). The proof in the initial paper uses the correct formulation of the theorem.
\end{Rem}


\begin{Rem}\label{remSleH}
The bound on $b_S(G)$ in Theorem \ref{fulltheoremPSL} is the best possible. It is achieved, for example, when $(2)$ of Theorem \ref{theoremGR} holds. Even though $S$ is not stated to be maximal solvable in Theorem~\ref{fulltheoremPSL}, it is sufficient for the proof to consider only maximal solvable subgroups. Indeed, if $H$ is a core-free (so $\cap_{g\in G} H^g=1$) subgroup of $G$  and $S \le H$, then $b_S(G) \le b_H(G)$.
\end{Rem}

\begin{Rem}
As already mentioned above, this paper is the first in the series of three where our goal is to prove Vdovin's conjecture for almost simple classical groups of Lie type.  In the upcoming second paper (based on \cite{thesis}) we deal with unitary and symplectic groups. The final third paper dealing with orthogonal groups is being prepared. 
\end{Rem}

The following lemma allows to omit the analysis of $\Reg_S(G,5)$ in many cases.
  \begin{Lem}[{\cite[Lemma 3]{bay}}] \label{base4}
Let $G$ be a finite permutation group with point stabiliser $H$. If  $b_H(G)\le 4$, then $\Reg_H(G,5)\ge 5.$
\end{Lem}

  By Remark \ref{remSleH}, if $G$ is almost simple, $S$ is a maximal solvable subgroup of $G$, and  $S \le H \le G$ (and $H$ does not contain the socle of $G$), then $b_S(G) \le b_H(G).$ Hence, results on the base size of primitive almost simple groups can provide useful information.
As we mention above, many results specifying bounds for base sizes of almost simple groups in primitive action are available. 
In particular, Burness \cite{burness} proves the following.
\begin{Th}
\label{bernclass}
Let $G \le \Sym(\Omega)$ be a non-standard finite almost simple  primitive permutation   group with classical socle and  point stabiliser $H$. Then either $b_H(G) \le 4$, or $G = \mathrm{U}_6 (2) \cdot 2$, $H = \mathrm{U}_4 (3) \cdot 2^2$ and $b_H(G) = 5$.
\end{Th}



Roughly speaking, Theorem \ref{bernclass} is true for maximal subgroups $H \notin \mathcal{C}_1$ (with some exceptions). Here $\mathcal{C}_i$ for $i=1, \ldots, 8$ are Aschbacher's classes introduced in \cite{asch}   and described in \cite[\S 2.1]{maxlow}  and \cite[Chapter 4]{kleidlieb}.  If $H \in \mathcal{C}_1,$ then it stabilises a subspace (or a pair of subspaces) of the natural module of $G$. Tables 2 and 3 in \cite{burness} contain detailed information on $b_H(G)$ for $n\le 5$, where $n$ is the dimension of the natural module.

If $G$ is classical of Lie type and $H \in \mathcal{C}_1$, then $b_H(G)$ can be arbitrarily large since the order
of $G$ is not always bounded by a fixed polynomial function of the degree of the action.  In particular, by Lemma \ref{logd}, $b_H(G)$ is not bounded by a constant. For example, if $G=\PGL_n(q)$ acting on the set of $1$-subspaces of the natural module, then the point stabiliser $H$ is a parabolic subgroup lying in  $\mathcal{C}_1$ and $b_H(G)=n+1.$ 

 Therefore,  if a maximal solvable subgroup $S$ lies only in a $\mathcal{C}_1$-subgroup of $G,$ then one cannot find a bound for $b_S(G)$ simply by studying the corresponding problem for maximal subgroups.
 
 \medskip
 
 Before we describe techniques and ideas we use to prove  Theorem \ref{fulltheoremPSL},  %In order to comment on the methods we use to establish our results and to understand the cases that arise in the proof 
 let us shortly state the structure of an almost simple group $G$ with socle $\PSL_n(q)$ and state more detailed results which also illustrate the cases that arise in the proof.
 
 Let $\GaL_n(q)$ be the group of invertible $\mathbb{F}_q$-semilinear transformations of $V=\mathbb{F}_q^n$, let  $\iota$ be the inverse-transpose map of $\GL_n(q)$, and let $Z$ be the group of all scalar matrices in $\GL_n(q)$ (see Section~\ref{secnot} for details). If $\PSL_n(q)$ is simple, so $n \ge 2$ and $(n,q)$ is neither $(2,2)$ nor $(2,3)$, then $$\Aut(\PSL_n(q)) \cong 
 \begin{cases}
 \PGaL_n(q) \rtimes \langle \iota \rangle =  ( \GaL_n(q) \rtimes \langle \iota \rangle)/Z & \text{ if } n \ge 3 \\
  \PGaL_n(q) =  \GaL_n(q) /Z & \text{ if } n = 2.
 \end{cases}$$
 
See \cite{Dieu} for details.  We achieve Theorem \ref{fulltheoremPSL} by proving the following two statements. 
\begin{mainthA}
\label{theorem}
Let $n \ge 2$ and $(n,q)$ is neither $(2,2)$ nor $(2,3).$ If $S$ is a maximal solvable subgroup of $\PGaL_n(q)$, 
 then $\Reg_S(S \cdot \PSL_n(q),5)\ge 5$. In particular $b_S(S \cdot \PSL_n(q)) \le 5.$
\end{mainthA}

\begin{mainthA}
\label{theoremGR}
Let $n\ge 3$. If $S$ is a maximal solvable subgroup of $\PGaL_n(q) \rtimes \langle \iota \rangle$ not contained in $\PGaL_n(q),$ then one of the following holds:
\begin{enumerate}
\item[$(1)$] $b_S(S \cdot \PSL_n(q))\le 4$;
\item[$(2)$] $n=4,$ $q\in \{2,3\}$, $S=P_2$, $b_S(S \cdot \PSL_4(q))=5$ and $\Reg_S(S \cdot \PSL_4(q),5)\ge 5.$
\end{enumerate}
\end{mainthA}

Here $P_2$ is the stabiliser in $\PGaL_n(q) \rtimes \langle \iota \rangle$ of a $2$-dimensional subspace of the natural module for $\PSL_n(q).$ The combination of the two theorems above gives a stronger statement than Theorem \ref{fulltheoremPSL}: if $\PSL_n(q)$ is simple and $S$ is a maximal solvable subgroup of $\Aut(\PSL_n(q))$, then $\Reg_S(S \cdot \PSL_n(q)) \ge 5.$ Indeed, Theorem \ref{fulltheoremPSL} follows from this statement via Lemma \ref{SdotG0}. 

\subsection{Methods and ideas.}\label{secmeth}

 It is more convenient for us to work with groups of matrices which makes the action  on the set of right cosets of a subgroup not faithful in most cases.  It is possible to extend the notion of $b_H(G)$ and $\Reg_H(G,k)$ to an arbitrary finite group $G$ by 
  $$b_H(G):=b(G/H_G) \text{ and } \Reg_H(G,k):=\Reg(G/H_G,k).$$
  Here $H_G = \cap_{g \in G} H^g$ and the action of $G/H_G$ on the set $\Omega$ of right cosets of $H$ in $G$ is induced from the natural action of $G$. Another straightforward but important notice is that the statement $b_H(G) \le k$ is equivalent to the existence of $k$  subgroups of $G$ conjugate to $H$ in $G$ such that their intersection is equal to $H_G.$ Indeed, $H^{g_1} \cap \ldots \cap H^{g_k}$ is the pointwise stabiliser of $(Hg_1, \ldots, Hg_k) \in~\Omega.$   In particular, in the proofs of Theorems \ref{theorem} and \ref{theoremGR}, we work with $\GaL_n(q)$ and $ \GaL_n(q) \rtimes \langle \iota \rangle $ rather than $\PGaL_n(q)$ and $\PGaL_n(q) \rtimes \langle \iota \rangle $.

  
 Now we are ready to outline main approaches and ideas of this paper.  We use combination of  probabilistic, constructive and computational methods to establish our results. Probabilistic method is described in details in Section \ref{fprsec}. It uses the ratio of points fixed by an element of the group $G \le \Sym(\Omega)$ to obtain an upper bound on the probability $Q(G,c)$ that a randomly chosen $c$-tuple of points in $\Omega$ is not a base. By showing $Q(G,c)<1$, one shows that there exists a base of size $c$. This method is used in most works related to base sizes of primitive permutation groups including  \cite{burness, sfa, bls, spor, lieb, mon}. We use this method to obtain bounds for $b_S(S \cdot \SL_n(q))$ for irreducible maximal solvable subgroups of $\GL_n(q).$  Our results  are refinements of Theorem \ref{bernclass} in the sense that they provide better estimates for $b_H(G)$  for solvable $H$ not lying in a $\mathcal{C}_1$-subgroup of $G \le \PGL_n(q)$.  In particular, with an explicit list of exceptions, we obtain $b_S(S \cdot \SL_n(q)) =2$ for an irreducible maximal solvable subgroup $S$ of $\GL_n(q)$  (see Theorem \ref{irred}).
 
  Probabilistic method does not work for us in the general case since the fixed points ratio is much harder to estimate when $S$ is reducible. Our reduction of the general case to the case of an irreducible subgroup of $\GL_n(q)$ is of constructive manner. We illustrate this for  a reducible maximal solvable subgroup $S$ of $\GL_n(q)$; the analysis for $S$ not contained in $\GL_n(q)$ is much more technical and splits into number of smaller cases but inherits the same general idea. So, let $S$ be a reducible maximal solvable subgroup of $\GL_n(q)$. By Lemma \ref{supreduce}, in a suitable basis, matrices in $S$ have blocks on their diagonal and all the entries below these blocks are zero. The projection on each block forms an irreducible solvable matrix group of smaller dimension for which we can apply Theorem \ref{irred}. As a result, we obtain that the intersection of two (in most cases) conjugates of $S$ consists of upper triangular matrices. Further, using this ``anti-symmetry'' we explicitly construct elements of $\SL_n(q)$ as in Lemmas \ref{irrtog}  and \ref{diag} that, as a result of using them as conjugating elements, give us, at first, a subgroup of diagonal matrices in the intersection of four conjugates of $S$, and then, the subgroup of scalar matrices in the intersection of five conjugates of $S$, so $b_S(S \cdot \SL_n(q))\le 5$.  The five distinct regular orbits then obtained by an explicit altering of these conjugating elements (see the proof of Proposition \ref{case5prop}).
  
  Both our theoretical tools, probabilistic method and the constructive method, turn out to be powerful yet technical and demanding to details techniques. They require knowledge of the structure of solvable subgroups of $\Aut(\PSL_n(q)).$  Here we heavily rely on the classical work of Suprunenko \cite{sup} on solvable matrix groups. Many useful details also can be found in \cite{short}.  
  
  We use the  computer algebra systems {\sf GAP} \cite{GAP4} and {\sc Magma} \cite{magma} to find and verify  $b_S(S \cdot \SL_n(q))$ for some small values of $n$ and $q$ when theoretical approach does not provide a sufficient result. Mostly, we turn to computations when an irreducible maximal solvable subgroup $S$ of $\GL_n(q)$ is ``too large'' and it is much easier to obtain results this way rather than establish a suitable estimate for $Q(S \cdot \SL_n(q), c).$ For details see Section \ref{gapsec}.


\subsection{Further comments and implications}

 

Vdovin's conjecture is originally stated in terms of intersection of subgroups in \cite[Problem 17.41 b)]{kt}:

\begin{Prob}\label{prob}
Let $S$ be a solvable subgroup of a finite group $G$ with trivial solvable radical.
 Do there always exist five conjugates of $S$ whose intersection is trivial?
\end{Prob}

As we mentioned above, since the stabiliser of $Sg$ in $G$ is $S^g$ for $g \in G$, the existence of five conjugates of $S$ whose intersection is trivial is equivalent to the statement that $b_S(G)\le 5.$  The number of conjugates of a subgroup with certain property one needs to obtain a normal subgroup in the intersection often provides useful information. In particular, as we discuss below, it gives an estimate on the index of such a normal subgroup. In a pioneering work \cite{passman} in 1966, Passman proves that for a Sylow $p$-subgroup $P$ of a $p$-solvable finite group $G$ there exist $x,y \in G$ such that $P \cap P^x \cap P^y = O_p(G)$ where $O_p(G)$ is the unique maximal normal $p$-subgroup of $G$. In 1966, Zenkov \cite{ZenP} extended this result to an arbitrary finite group $G.$  Similar results are achieved by Zenkov \cite{Zen93}, Vdovin \cite{vdovinReg} and Dolfi \cite{dolfi} for solvable Hall $\pi$-subgroups of a $\pi$-solvable group $G$.  We use the following well known result of Zenkov in our proofs on multiple occasions.   

\begin{Th}[\cite{zen}]
\label{zenab}
If $A$ and $B$ are abelian subgroups of a finite group $G$, then there exists $x \in G$ such that $A \cap B^x \le {\bf F}(G).$
\end{Th}
Here ${\bf F}(G)$ is the Fitting subgroup of $G$ (the maximal normal nilpotent subgroup).
In a recent paper \cite{zen21}, Zenkov shows that for nilpotent subgroups $A,B,C$ of a finite group $G$ there exist $x,y \in G$ such that $A \cap B^x \cap C^y \le {\bf F}(G)$.


%A positive solution of Problem \ref{prob} would lead to an answer

The study of intersections of subgroups is connected directly to a more general and abstract question of interest. Consider some property $\Psi$ of a finite group inherited by all its subgroups.  Important examples of  such a property are  the following: cyclicity, commutativity, nilpotence, solvability.
A natural question arises: how large is a normal $\Psi$-subgroup in an arbitrary finite group $G$? A more precise formulation of this question is the following:

\begin{Que}\label{que1}
Given a finite group $G$ with $\Psi$-subgroup $H$ of index $n$, is it true that $G$ has a normal $\Psi$-subgroup whose index is bounded by some function $f(n)?$
\end{Que}

Since $H_G$, which is a normal $\Psi$-subgroup, is the kernel of the action of $G$ on the set of right cosets of $H$, and such an action provides a homomorphism to the symmetric group $\Sym(n),$   it always suffices to take $f(n)=n!$ for every such $\Psi$. We are interested in  stronger bounds, in particular those of shape $f(n)=n^c$ for some constant $c.$



Babai, Goodman and Pyber \cite{bab} prove some related results and state several conjectures.
In particular, they prove that
if a finite group $G$ has a cyclic subgroup ${C}$ of index $n$, then $ \cap_{g \in G} {{C}}^g$ has index at most $n^7.$
They also conjectured that the bound $n^2-n$ holds  and showed that it is best possible. Lucchini \cite{luccini} and, independently,  Kazarin and  Strunkov \cite{kaz} proved this. It follows from 
results  by Chermak and  Delgado \cite{cher} that the bound $n^2$ suffices for commutativity, and from results  by Zenkov \cite{zen21} that the bound $n^3$ suffices for nilpotence. We are interested in the case when $\Psi$ is solvability.



Babai, Goodman and Pyber \cite{bab}   prove the following statement.
\begin{Th}\label{indexBab}
There is an absolute constant $c$ such that, if a finite group $G$ has a
solvable subgroup $S$ of index $n$, then $G$ has a solvable normal subgroup of index
at most $n^c$.
\end{Th}

\noindent Although their proof does not yield an explicit value, they  conjectured  that $c \le 7$. It is straightforward to obtain a bound for $c$ from $b_S(G)$ for a maximal solvable subgroup $S$.
 Indeed, let  $G$ act on the set $\Omega$ of right $S$-cosets via right multiplication, so $|\Omega|=|G:S|.$  If $(\beta_1 , \ldots , \beta_k )\in \Omega^k$ is a base for the natural action of $G/S_G$ on $\Omega$, then $$|(\beta_1 , \ldots , \beta_k )^G| \le |\Omega| \cdot (|\Omega|-1) \ldots (|\Omega|-k+1)<|\Omega|^k=|G:S|^k.$$
Therefore,
$$|G:S_G|<|G:S|^k,$$
and if Problem \ref{prob} has a positive answer, then $c \le 5$ in Theorem \ref{indexBab} (we can assume that $S$ is maximal solvable, so $S_G$ is the solvable radical of $G$ and $G/S_G$ satisfies the condition of Problem \ref{prob}).

\bigskip







   
    
    The paper is organised as follows. In Section 2 we state necessary definitions and preliminary results. In Section 3 we obtain bounds for $b_S(S \cdot \SL_n(q))$ for irreducible maximal solvable subgroups of $\GL_n(q).$   Finally, we  use these results to prove Theorems \ref{theorem} and \ref{theoremGR} in Sections \ref{sec411} and \ref{sec412} respectively.  

\section{Definitions and preliminaries}\label{secnot}

All group actions are right actions. For example, the action of a linear transformation $g$ of a vector space $V$ on $v \in V$ is $(v)g \in V.$

 Let $p$ be a prime and $q=p^f$,  $f \in \mathbb{N}$. Denote a finite field of size $q$ by $\mathbb{F}_q$, its algebraic closure by $K$ and the multiplicative group of a field $\mathbb{F}$ by $\mathbb{F}^*.$ 
Throughout, unless stated otherwise, $V=\mathbb{F}_{q}^n$ denotes a vector space of dimension $n$ over $\mathbb{F}_{q}$.

For finite classical groups we follow the  notation of \cite{kleidlieb}. For algebraic groups our standard references are \cite[Chapter 1]{carter}, \cite[Chapter 1]{gorlyo} and \cite{hump}.
An $\mathbb{F}$-{\bf semilinear} transformation of a vector space $V$ over a field $\mathbb{F}$ is a map $g: V \to V$ such that there exists a field automorphism $\sigma(g) \in \Aut(\mathbb{F})$ with 
$$(v+u)g=(v)g+ (u)g \text{ and } (\lambda v)g= \lambda^{\sigma(g)}(v)g$$
for $v,u \in V$ and $\lambda \in \mathbb{F}.$
An $\mathbb{F}$-semilinear transformation $g$ is {\bf non-singular} if $\{v \in V \mid (v)g=0\}=\{0\}.$ We define the group $\GaL(V,\mathbb{F})$ to be the set of all non-singular $\mathbb{F}$-semilinear transformations of $V$ with composition as the binary operation.   
We reserve the letter $\beta$ for a {\bf basis} of $V$. A basis is an ordered set. 
If $\alpha \in \Aut(\mathbb{F})$, then $\phi_{\beta}(\alpha)$ denotes the unique  $g \in \GaL(V, \mathbb{F})$ such that
\begin{equation}
\label{defphibet}
\left(\sum_{i=1}^n \lambda_i v_i \right) \phi_{\beta}(\alpha)= \sum_{i=1}^n\lambda_i^{\alpha} v_i
\end{equation}
 where $\beta =\{v_1, \ldots, v_n\}.$ If $\mathbb{F}= \mathbb{F}_{q}$ and $\alpha \in \Aut(\mathbb{F})$ is such that $\lambda^{\alpha}=\lambda^p$ for all $\lambda \in \mathbb{F}$,  then we denote $\phi_{\beta}(\alpha)$ by $\phi_{\beta}$ \index{$\phi_{\beta}$} or simply $\phi$ when $\beta$ is understood. It is routine to check that  $$\GaL(V, \mathbb{F}_q)= \GL(V, \mathbb{F}_q) \rtimes \langle \phi \rangle \cong \GL_n(q) \rtimes \langle \phi \rangle.$$


\medskip

\begin{Not}\label{not} We fix the following notation.
\begin{longtable}{p{1em} p{7em} p{27em} }
1. & $p'$ & set of all primes except $p$\\
2. & $(a,b)$ & greatest common divisor of integers $a$ and $b$\\
3. & $g^G$ & conjugacy class of $g \in G$\\
4. & $Z(G)$ & {\bf center} of a group $G$\\
5. & ${\bf F}(G)$ & {\bf Fitting subgroup} \index{Fitting subgroup} of a finite group $G$ (unique maximal \\ & & normal  nilpotent subgroup)\\
6. & $O_{\pi}(G)$ & unique maximal normal $\pi$-subgroup for a set of primes $\pi$\\
7. & $A \rtimes B$ & semidirect product of groups $A$ and $B$ with $A$ normal\\
8. &  $\Sym(n)$ &  symmetric group of degree $n$ \\
9. & $\sgn(\pi)$ &  {\bf sign} of a permutation $\pi$\\
10. & $M_n(\mathbb{F})$& algebra of all $n \times n$ matrices over $\mathbb{F}$\\
11. & $D(G)$ & subgroup of all diagonal matrices of a matrix group $G$\\
12. & $RT(G)$ & subgroup of all upper-triangular %(right-triangular)
 matrices of a matrix group $G$\\
13. & $\Det(H)$ & $\{\det(h) \mid h \in H \}$ for $H \le \GL(V)$\\ 
14. & $\diag(\alpha_1, \ldots, \alpha_n)$ &  {\bf diagonal} matrix with entries $\alpha_1, \ldots, \alpha_n$ on its diagonal  \\  
15. &  $\diag[g_1, \ldots, g_k]$ &  {\bf block-diagonal} matrix with blocks $g_1, \ldots, g_k$ on its diagonal  \\
16. & $\per(\sigma)$ & {\bf permutation matrix} corresponding to $\sigma \in \Sym(n)$\\
17. & $g^{\top}$ & {\bf transpose} of a matrix $g$\\
18. & $g \otimes h$ &  {\bf Kronecker product} \index{Kronecker product}   
 $\begin{pmatrix}
g \cdot h_{1,1}     &  \ldots & g \cdot h_{1,m}  \\
\ldots           &   \ldots   & \ldots   \\
g \cdot h_{m,1}     & \ldots     & g \cdot h_{m,m}      
\end{pmatrix} \in \GL_{nm}(q)$ \\ & & for $g \in \GL_n(q)$ and $h \in \GL_m(q)$ \\
\end{longtable}
\end{Not}

 An irreducible subgroup $G$ of $\GL(V)$ is {\bf imprimitive} if there exist a decomposition 
 $$V=V_1 \oplus \ldots \oplus V_k; \text{ } k>1$$  stabilised by $G$. So for each $i \in \{1, \ldots, k\}$ and $g \in G$ there exist   $j \in \{1, \ldots, k\}$ such that $(V_i)g=V_j.$ If such a decomposition does not exist, then $G$ is {\bf primitive}.

It is convenient to view  the symmetric group 
as a group of permutation matrices. We define the {\bf wreath product} \index{wreath product}  of  $X \le \GL_n(q)$ and a  group of permutation matrices $Y \le \GL_m(q)$ as the matrix group $X \wr Y \le \GL_{nm}(q)$ obtained
by replacing the entries 1 and 0  in every matrix in $Y$ by
arbitrary matrices in $X$ and by zero $(n \times n)$ matrices respectively. 

Let $A$ be an $(nm\times nm)$ matrix. We can view $A$ as the matrix
$$\begin{pmatrix}
A_{11}      &  \ldots & A_{1m}  \\
\ldots           &   \ldots   & \ldots   \\
A_{m1}     & \ldots     & A_{mm}      
\end{pmatrix}$$
where the $A_{ij}$ are $(n \times n)$ matrices. The vector $(A_{i1}, \ldots, A_{im})$ is the $i$-th $(n \times n)$-{\bf row} \index{$(n \times n)$-row} of $A.$



\subsection{General results and solvable linear groups}
\label{missec}

We begin by proving general technical results that are mentioned in Introduction.

\begin{Lem}\label{logd}
If $G\le \Sym(\Omega)$ and $d=|\Omega|,$ then $b(G) \ge \log_d|G|$.
\end{Lem}
\begin{proof}
Let $B \in \Omega^{b(G)}$ be a base. Every element of $G$ is uniquely determined by its action on $B$. Indeed, if $Bx=By$ for $x,y \in G$, then $Bxy^{-1}=B$ and $x=y$ since $B$ is a regular point. Hence $|G|\le d^{b(G)}.$ 
\end{proof}

\begin{Lem}\label{SdotG0}
Let $G_0$ be a finite simple nonabelian group. Let $G_0 \le G \le \Aut (G_0)$ and let $S\le G$ be solvable. If  $\Reg_H(H \cdot G_0,5)\ge 5,$ for every maximal solvable subgroups $H$ of $\Aut(G_0),$ then $\Reg_S(G,5)\ge 5,$ for every solvable subgroup $S$ of $G$.
\end{Lem}
\begin{proof} Let $G_0 \le G \le \Aut (G_0)$ and let $S\le G$ be solvable. Let $H$ be a maximal solvable subgroup of $  \Aut(G_0)$ containing $S$. Since 
$$\Reg_H(H \cdot G_0,5)\ge 5,$$
 there exist $x_{(i,1)},x_{(i,2)},x_{(i,3)},x_{(i,4)},x_{(i,5)} \in G_0$ for $i \in\{1, \ldots, 5\}$ such that
$$\omega_i=(Hx_{(i,1)},Hx_{(i,2)},Hx_{(i,3)},Hx_{(i,4)},Hx_{(i,5)})$$ 
are $H \cdot G_0$-regular points   in $\Omega^5=\{Hx \mid x \in H \cdot G_0 \}^5$, and the $\omega_i$ lie in distinct orbits. 
We claim that $$\omega_i'=(Sx_{(i,1)},Sx_{(i,2)},Sx_{(i,3)},Sx_{(i,4)},Sx_{(i,4)}) \text{ for } i \in \{1, \ldots, 5\}$$ lie in distinct $G$-regular orbits in $(\Omega')^5=\{Sx \mid x \in G\}^5.$ Indeed, $$S^{x_{(i,1)}} \cap S^{x_{(i,2)}} \cap S^{x_{(i,3)}} \cap S^{x_{(i,4)}} \cap S^{x_{(i,5)}} \le H^{x_{(i,1)}} \cap H^{x_{(i,2)}} \cap H^{x_{(i,3)}} \cap H^{x_{(i,4)}} \cap H^{x_{(i,5)}}=1,$$ so $\omega_i'$ are regular. Assume that $\omega_1' g = \omega_2'$ for some $g \in G.$ Therefore,
$$(Sx_{(1,i)})g=Sx_{(2,i)} \text{ for } i \in \{1, \ldots, 5\}$$
and 
$$g \in \bigcap_{i=1}^5 (x_{(1,i)}^{-1}S x_{(2,i)}) \subseteq \bigcap_{i=1}^5 (x_{(1,i)}^{-1}H x_{(2,i)})= \emptyset$$
where the last equality holds since $\omega_1$ and $\omega_2$ lie in distinct $H \cdot G_0$-orbits. Hence $\omega_1'$ and $\omega_2'$ lie in distinct $G$-orbits. The same argument shows that all of the $\omega_i'$ lie in distinct $G$-orbits, so $\Reg_S(G,5)\ge 5$.
\end{proof}


  Further we state two classical results on the structure of solvable linear groups.
\begin{Lem}[{\cite[\S 18, Theorem 5]{sup}}] \label{supirr}
An irreducible solvable subgroup of $\GL_n(q)$ is either primitive, or conjugate in $\GL_n(q)$ to a subgroup of the wreath product
$S \wr \Gamma$
where $S$ is a primitive solvable subgroup of $\GL_m(q)$ and $\Gamma$ is a transitive solvable
subgroup of the symmetric group $\Sym(k)$ and $km = n$. In particular, an irreducible maximal solvable  subgroup of 
$\GL_n(q)$ is either primitive, or conjugate in $\GL_n(q)$
to  $S \wr \Gamma$, where $S$ is a  primitive maximal solvable subgroup of $\GL_m(q)$
and $\Gamma$ is a  transitive maximal solvable subgroup of $\Sym(k)$.
\end{Lem}

\begin{Lem}[{\cite[\S 18, Theorem 3]{sup}}]
\label{supreduce}
Let  $H$ be a
subgroup of $\GL_n(q)$. In a suitable basis $\beta$ of $V$, the matrices $g\in H$ have the
shape
\begin{gather}\label{stup}
\begin{pmatrix}
g_k     & g_{k,(k-1)} & \ldots & g_{k,1}  \\
0          &    g_{k-1}  & \ldots & g_{(k-1),1}  \\
\ldots     & \ldots     & \ldots & \ldots     \\
      0    &   0        & \ldots & g_1 
\end{pmatrix}
\end{gather}
where the mapping $\gamma_i: H \to \GL_{n_i}(q)$, $g \mapsto g_i$ is an irreducible 
representation of $H$ of degree $n_i$, and $g_{i,j}$ is an $(n_i \times n_j)$ matrix over $\mathbb{F}_q$, and $n _1 + \ldots + n _k = n$.
 The group $H$  is solvable if and only if all the
groups
$$H_i ={\mathrm{Im} } (\gamma_i) \text{ for } i = 1, \ldots, k,$$
are solvable.
\end{Lem}

 Now we state three technical lemmas about solvable linear groups. We use them and ideas from their proofs many times throughout  this paper. 

\begin{Lem}\label{supSL}
Let $H \le G \le \GL_n(q).$ Assume $\Det(H)=\Det(G).$ Then the following hold:
\begin{enumerate}[font=\normalfont]
\item $H \cdot (\SL_n(q) \cap G)=G$.
\item If $g \in G$, then there exists $g_1 \in \SL_n(q) \cap G$ such that $H^g=H^{g_1}.$ 
\end{enumerate}
\end{Lem}
\begin{proof}
Let $g \in G$. Since $\Det(H)=\Det(G)$, there exists $h \in H$ such that $\det(h)=\det(g).$ Therefore, 
$h^{-1}g \in \SL_n(q) \cap G$, so $g=h \cdot (h^{-1}g)\in H \cdot (\SL_n(q) \cap G).$

Let $g_1=h^{-1}g,$ so $H^{g_1}=g^{-1}hHh^{-1}g=g^{-1}Hg=H^g.$
\end{proof}

\begin{Lem}\label{irrtog}
Let $H$ be a subgroup of $\GL_n(q)$ of shape \eqref{stup}. If for every $H_i$ there exists $x_i \in \GL_{n_i}(q)$ (respectively $\SL_{n_i}(q)$) such that the intersection $H_i \cap H_i^{x_i}$ consists of upper triangular matrices, then there exist $x,y \in \GL_n(q)$ (respectively $\SL_n(q)$) such that 
$$(H\cap H^x) \cap (H\cap H^x)^y \le D(\GL_n(q)).$$
\end{Lem}
\begin{proof}
Let $x=\diag(x_k, \ldots, x_1)$ and let $y$ be $$\diag(\sgn(\sigma),1 \ldots,1) \cdot \per(\sigma)$$ where  $$\sigma =(1,n)(2, n-1) \ldots ([n/2], [n/2+3/2]).$$ Here $[r]$ is the integer part of a positive number $r$. Since $\det(\per(\sigma))=\sgn(\sigma),$ $\det(y)=1.$ Since
$H\cap H^x$ consists of upper triangular matrices and $(H\cap H^x)^y$ consists of lower triangular matrices, 
\begin{equation*}
(H\cap H^x) \cap (H\cap H^x)^y \le D(\GL_n(q)). %\qedhere
\end{equation*} 
Finally, note that if $x_i \in \SL_{n_i}(q)$, then $x \in \SL_n(q).$ % Notice that
\end{proof}




\begin{Def}
Let $\{v_1, v_2, \ldots,  v_n \}$ be a basis of $V$. The equality 
$$v=\alpha_1v_1+ \alpha_2v_2+  \ldots+ \alpha_nv_n $$ 
is the {\bf decomposition} \index{decomposition} of $v \in V$ with respect to $\{v_1, v_2, \ldots,  v_n \}$.
 \end{Def}


\begin{Lem}\label{diag}
Assume that matrices in $H\le \GL_n(q)$ have shape \eqref{stup} with respect to the  basis $$\{v_1, v_2, \ldots,  v_n \}$$ and $n_1 < n,$ so $H$ stabilises $$U=\langle v_{n-n_1+1}, \ldots, v_n \rangle.$$  Then there exists $z \in \SL_n(q)$ such that $D(\GL_n(q)) \cap H^z \le Z(\GL_n(q)).$
\end{Lem} 
\begin{proof}
Let $m=n_1.$ Define vectors 
\begin{equation*}
\begin{split}
u_1 & = v_1 + \ldots + v_{n-m} + v_{n-m+1}  \\
u_2 & = v_1 + \ldots + v_{n-m} + v_{n-m +2}\\
u_3 & = v_1 + \ldots + v_{n-m}  + v_{n-m +3}\\
\vdots \\
u_m & = v_1 + \ldots + v_{n-m} + v_{n}.
\end{split}
\end{equation*}
Let $z \in \GL_n(q)$ be such that 
\begin{equation}\label{ui}
(v_{n-m+i})z=u_i, \text{ } i=1, \ldots, m.
\end{equation}
 Such  $z$ exists since $u_1, \ldots, u_m$ are linearly independent and we can assume $z\in \SL_n(q)$ since $m<n.$ So \eqref{ui} implies that 
$$Uz = \langle u_1, \ldots, u_m \rangle $$ is $H^z$-invariant.

Let $g \in D(\GL_n(q)) \cap H^z.$ So $g$ is $\text{diag}(\alpha_1, \ldots , \alpha_n)$ with respect to the basis $\{v_1, v_2, \ldots, v_{n}\}$. Thus,
$$(u_j)g = \alpha_1 v_1 + \ldots + \alpha_{n-m} v_{n-m} + \alpha_{n-m+j} v_{n-m+j}. $$
Also, $(u_j)g$ must lie in $Uz,$ since $g \in H^z,$ so 
$$(u_j)g = \beta_1 u_1 + \beta_2 u_2 + \ldots + \beta_m u_m.$$
But the decomposition of $(u_j)g$ with respect to    $\{v_1, v_2, \ldots, v_{n}\}$ does not contain $v_{n-m +i}$ for $i \ne j$, so 
$$\beta_1 = \beta_2 = \ldots =\beta_{j-1} =\beta_{j+1}= \ldots = \beta_m =0$$ and $(u_j)g= \beta_j u_j$. Thus, 
$$\alpha_1 = \alpha_2 = \ldots =  \alpha_{n-m}=\alpha_{n-m+j}$$
for $0<j\le m.$
Therefore, $g$ is scalar and lies in $Z(\GL_n(q)).$
\end{proof}

The following lemma plays an important role in our proof of Theorems \ref{theorem} and \ref{theoremGR}. 

\begin{Lem}
\label{scfield}
Let  $\phi=\phi_{\beta}$ where $\beta$ is a basis of $V,$ and assume $n\ge 2$ with $(n,q) \ne (2,2),(2,3).$ Let $H \le \GaL_n(q) $ and assume  $H \cap \GL_n(q) \le Z(\GL_n(q))$. Then there exists $b \in \GL_n(q)$ such that every element of $H^b$ has the form $\phi^i g$ for some $i \in \{1, \ldots, f\}$ and $g \in Z(\GL_n(q)).$   
\end{Lem}
\begin{proof}
 Let $Z=Z(\GL_n(q))$ and $\Gamma=\GaL_n(q).$ Notice that $\Gamma/Z$ is almost simple. Let $G_0$ and $\hat{G}$ be the socle of $\Gamma/Z$ and the group of inner-diagonal automorphisms of $G_0$ respectively, so $G_0=\PSL_n(q)$ and $\hat{G}= \PGL_n(q).$ Without loss of generality, we may assume $Z \le H.$ Observe $H\cap \GL_n(q) = Z,$ so $H/Z$ is cyclic and consists of field automorphisms of $G_0.$ Let $\varphi \in H$ be such that $\langle Z \varphi \rangle=H/Z.$   Since 
$$\Gamma=\GL_n(q) \rtimes \langle \phi \rangle,$$
we obtain that $\varphi \in \phi^i  \GL_n(q)$ for some  $i \in \{1, \ldots, f\}$ and $Z\varphi \in (Z\phi^i) \hat{G}$.

 Since $Z \varphi$ and $Z \phi^i$ can be viewed as field automorphisms of $G_0$ of the same order, they are conjugate in $\hat{G}$ by \cite[(7-2)]{conjaut}. So   there exists  $Zb \in \PGL_n(q)$ such that $(Z\varphi)^{Zb}=Z\phi^i$ for some $i \in \{1, \ldots,  f\}.$  Therefore, $H^b = Z\langle \varphi^b \rangle = Z\langle \phi^i \rangle = \langle\phi^i \rangle Z.$
\end{proof}





\subsection{Singer cycles}\label{singsec}

\begin{Def}\label{sindef} A  {\bf Singer cycle} \index{Singer cycle} of $\GL_n(q)$ is a cyclic subgroup of order $q^n - 1$.
\end{Def}

 A Singer cycle always exists by {\cite{singer38}}, and all Singer cycles are conjugate in $\GL_n(q)$ as we prove in Lemma \ref{singconj}. Sometimes, it is convenient to identify a Singer cycle with the multiplicative group of a filed of order $q^n$ as follows. A field $\mathbb{F}_{q^n}$ can be considered as an $n$-dimensional vector space $V=\mathbb{F}_q^n$ over $\mathbb{F}_q.$  Right multiplication by a generator of  $\mathbb{F}_{q^n}^*$ determines a bijective linear map from $V$ to itself of order $q^n-1$. So  $\mathbb{F}_{q^n}^*$  is isomorphic to a cyclic subgroup of $\GL_n(q)$ of order $q^n-1$. It is well known, and is easy to deduce from the observation above, that such a subgroup $T$ is irreducible, and $T \cup \{0\}$ is a field under the usual matrix addition and multiplication. 



\begin{Lem}\label{singconj} 
All Singer cycles are conjugate in $\GL_n(q).$ 
\end{Lem}
\begin{proof}
 By \cite[Satz 3.10]{hupp}, an irreducible cyclic subgroup of $\GL_n(q)$ is conjugate to a subgroup of the Singer cycle constructed below Definition \ref{sindef}. It remains to show that a Singer cycle is irreducible.

Let $T$ be a Singer cycle. By Maschke's Theorem, $T$ is completely reducible, so 
$V = V_1 \oplus \ldots \oplus V_k$ where $T$ stabilises the $V_i$ and each $V_i$ is $\mathbb{F}_q[T]$-irreducible.  Notice that $q^n-1$ is the largest possible order of an element of $\GL_n(q).$ Indeed,  let $\chi_x(t)$ be the characteristic polynomial of $x \in \GL_n(q)$. By the Cayley--Hamilton Theorem $\chi_x(x)=0$ and the dimension of the subalgebra $\mathbb{F}_q[\langle x \rangle]$ of $M_n(\mathbb{F}_q)$, generated by $x$, is at most 
$\deg \chi_x(t) \le n.$ So $|\mathbb{F}_q[\langle x \rangle]|\le|\mathbb{F}_q|^n=q^n$ and $|x|\le q^n-1.$ Therefore,  $|T| \le \prod_i^k (q^{n_i}-1)$ which is  possible only if $k=1$, so $T$ is irreducible. 
\end{proof} 

\begin{Th}[{\cite[Chapter II, \S 7]{hupp}}]\label{sin1} 
If $T$ is a Singer cycle of $\GL_n(q)$, then $$C_{\GL_n(q)}(T) = T$$
and $N_{\GL_n(q)} (T) /T$ is cyclic of order $n$. Moreover, $N_{\GL_n(q)} (T)=T \rtimes \langle \varphi \rangle$, where $t^{\varphi}=t^q$ for $t \in T.$ 
\end{Th}




Let $t \in \GL_n(q)$ be a generator of a Singer cycle. By the identification of $\mathbb{F}_q^n$ and $\mathbb{F}_{q^n}$ above, the charateristic polinomial of $t$ is the characteristic polynomial $\prod_{i=0}^{n-1}(x-\lambda^{q^i})$ of a primitive element $\lambda \in \mathbb{F}_{q^n}$ over $\mathbb{F}_q.$ So $t$ is conjugate to $\diag(\lambda, \lambda^q, \ldots, \lambda^{q^{n-1}})$ in $\GL_n(K)$ where $K$ is the algebraic closure of $\mathbb{F}_q.$ In particular, $\det(t)$ generates $\mathbb{F}_q^*,$ so $\langle t \rangle \SL_n(q)=\GL_n(q)$ by Lemma \ref{supSL}.

\subsection{Fixed point ratios and elements of prime order}
\label{fprsec}

\begin{Def}
If a group $G$ acts on a set $\Omega$, then   $C_{\Omega}(x)$ is the set of points in $\Omega$  fixed
by   $x \in G$. If $G$ and $\Omega$ are finite, then  the {\bf fixed point ratio} \index{fixed point ratio} of $x$, 
denoted  by $\fpr(x)$, is the proportion of points in $\Omega$ fixed by $x$, i.e. $\fpr(x) = |C_{\Omega}(x)|/|\Omega|.$
\end{Def}




If  $G$ acts transitively on a set $\Omega$  and $H$ is a point stabiliser,  then it is easy to see that
\begin{equation}\label{fpr}
\fpr(x)= \frac{|x^{G} \cap H|}{|x^G|}
\end{equation}


 In \cite{fpr, fpr2, fpr3, fpr4} Burness studies   fixed point ratios for primitive actions of almost simple classical groups and we recall some  observations from \cite{fpr}.
 
Let a group $G$ act faithfully   on the set $\Omega$ of right cosets of a subgroup $H$ of $G.$  
Let $Q(G, c)$ be the probability that
a randomly chosen $c$-tuple of points in $\Omega$ is not a base for $G$, so $G$ admits a base of size $c$ if and
only if $Q(G, c) < 1$. Of course, a $c$-tuple is not a  base if and only if it is fixed by some element  $x \in G$ of prime order, and the probability that a random $c$-tuple is fixed
by $x$ is equal to $\fpr(x)^c$. Let $\mathscr{P}$ be the set of elements of prime order in ${G}$, and let ${x}_1, \ldots, {x}_k$
be representatives for the ${G}$-classes of elements in $\mathscr{P}$. Since  fixed point
ratios are constant on conjugacy classes by  \eqref{fpr}, 
\begin{equation}\label{ver}
Q(G,c) \le \sum_{{x} \in \mathscr{P}}\fpr({x})^c = \sum_{i=1}^{k}|{x_i}^{{G}}|\cdot\fpr({x}_i)^c=:\widehat{Q}(G,c).
\end{equation}

The following observation is useful for bounding $\widehat{Q}(G,c)$ in \eqref{ver}.
 
\begin{Lem}[{\cite[Lemma 2.1]{burness}}]
\label{fprAB}
Let $G$ act faithfully and transitively on $\Omega$ and let $H$ be a point stabiliser. If  $x_1,\ldots,x_k$ represent distinct $G$-classes such that $\sum_{i=1}^k |x_i^G \cap  H| \le A$ and $|x_i^G| \ge B$ for all $i \in \{1, \ldots, k\},$ then
$$\sum_{i=1}^k |x_i^G| \cdot \fpr (x_i)^c \le B \cdot (A/B)^c.$$ 
for all $c \in \mathbb{N}.$
\end{Lem}


Notice that if there exists $\xi \in \mathbb{R}$ such that $\fpr({x})\le |{x}^{{G}}|^{-\xi}$ for every  ${x} \in \mathscr{P}$,  then  
$$\widehat{Q}(G,c)\le \sum_{i=1}^{k}|{x}_i^{{G}}|^{1-c\xi}.$$



\begin{Def}
Let $\mathscr{C}$ be the set of conjugacy classes of prime order elements in ${G}$.
For $t \in \mathbb{R},$ 
$$\eta_G(t):= \sum_{C \in \mathscr{C}}|C|^{-t}.$$
If $Z(G)=1,$ then there exists $T_G \in \mathbb{R}$ such that $\eta_G(T_G) = 1.$
\end{Def}





\begin{Lem}\label{11}
	If $G$ acts faithfully and transitively on $\Omega$ and  $\fpr({x})\le|{x}^{{G}}|^{-\xi}$ for all  ${x} \in \mathscr{P}$ and $T_G < c\xi - 1$, then $b(G) \le c$.
\end{Lem}
\begin{proof}
We follow the proof of \cite[Proposition 2.1]{burness}. Let ${x}_1, \ldots, {x}_k$ be representatives of the ${G}$-classes
of prime order elements in ${G}$. %, and let $Q(G, c)$ be the probability that a randomly chosen
%$c$-tuple of points in $\Omega$ is not a base for $G$. Evidently, $G$ admits a base of size $c$ if and only if
%$Q(G, c) < 1.$ 
By \eqref{ver}, 
$$Q(G,c) \le  \sum_{i=1}^{k}|{x_i}^{{G}}|\cdot \fpr({x}_i)^c \le \eta_G(c\xi - 1).$$
The result follows since $\eta_G(t) < 1$ for all $t>T_G.$
\end{proof}

%\begin{Def} \label{algdef}
We fix the following notation for the rest of the section. Let $\overline{G}$ be an adjoint simple algebraic group of classical type over the algebraic closure $K$ of ${\mathbb{F}_p}$. Let $\overline{G}_{\sigma}=\{g \in \overline{G} \mid g^{\sigma}=g\}$ where $\sigma$ is a Frobenius morphism of $\overline{G}.$ Let $\overline{G}$ be such that $G_0=O^{p'}(\overline{G}_{\sigma})'$ is a finite simple group of Lie type over $\mathbb{F}_q$ for some $p$-power $q$. Here  $O^{p'}(G)$ is the subgroup of a finite group $G$ generated by all $p$-elements of $G$. In particular, if $\overline{G}=\PSL_n(K),$ then $\overline{G}_{\sigma}=\PGL_n^{\varepsilon}(q)$ and $G_0=\PSL_n^{\varepsilon}(q)$ with $\varepsilon=\pm$ (here $\varepsilon =+$ corresponds to linear groups and $\varepsilon =+$ corresponds to unitary groups, so $\PGL_n^+(q)=\PGL_n(q)$ and $\PGL_n^-(q)=\PGU_n(q)$). Let  $G$ be a finite almost simple  group with  socle $G_0.$  
%\end{Def}

As proved in \cite[Proposition 2.2]{burness},  if  $n \ge 6$, then $T_G$ exists and $T_G<1/3$. 
Thus, if for such $G$  
\begin{equation}\label{0} 
\fpr(x)<|x^G|^{-\frac{4}{3c}}
\end{equation}
 for all $x \in \mathscr{P} $, then $\xi \ge {4}/{(3c)}$ and $c\xi -1 \ge 1/3 >T_G$ and $G$ has a base of size $c$.

Therefore, Lemma \ref{11} allows us to estimate the base size by calculating bounds for $|x^G|$ and $|x^G \cap H|$ for elements $x$ of prime order. 





\begin{Def}
\label{nudef}
Let $x \in \PGL(V)=\PGL_n(q)$. Let $K$ be the algebraic closure of $\mathbb{F}_q$, and let $\overline{V}=K \otimes V.$ Let $\hat{x}$ be a preimage of $x$ in $\GL_n(q).$ Define
$$ \nu_{V, K}(x):= \min\{\dim [\overline{V}, \lambda \hat{x}] : \lambda \in K^* \}. $$ Here $[V,g]$ for a vector space $V$ and $g \in \GL(V)$ is the commutator in $V \rtimes \GL(V)$.  Therefore, $\nu_{V, K}(x)$ is  the minimal codimension of an eigenspace of $\hat{x}$  on $\overline{V}$. In particular, $\nu_{V, K}(x)$ does not depend on the choice of the preimage $\hat{x}.$  Sometimes we denote this number by $\nu(x)$ and $\nu_{V, K}(\hat{x})$. 
\end{Def}

\begin{Lem}[{\cite[Lemma 3.11]{fpr2}}]\label{prost}  %\label{prost}
Let $x \in \PGL_n(q)$ have  prime order $r.$ One of the following holds: 
\begin{enumerate}[font=\normalfont]
\item $x$ lifts to  $\hat{x} \in \GL_n(q)$ of order $r$ such that $|x^{\PGL_n(q)}|=|\hat{x}^{\GL_n(q)}|;$ \label{prost1}
\item $r$ divides both $q-1$ and $n$, and $x$ is $\PGL_n(K)$-conjugate to the image of $$\diag[I_{n/r}, \omega I_{n/r}, \ldots , \omega^{r-1} I_{n/r} ],$$ where $\omega \in K$ is a primitive $r$-th root of unity. \label{prost2}
\end{enumerate} 
\end{Lem}

\begin{Rem}
Lemma 3.11 from \cite{fpr2} is formulated for   all classical groups, but only for $r \ne 2$. It is easy to see  from its proof  that the condition $r\ne 2$ is necessary only for orthogonal and symplectic cases; if $x \in \PGL_n(q)$ then the statement is true for all primes $|x|.$  
\end{Rem}



%In the following lemmas we use notation from Definition \ref{algdef}.
\begin{Lem}\label{xGoGs}
Let $x \in \overline{G}_{\sigma}$ have prime order.
\begin{enumerate}[font=\normalfont]
\item If $x$ is semisimple, then $x^{\overline{G}_{\sigma}}=x^{G_0}.$
\item If $x$ is unipotent and $G_0=\PSL_n(q)$, then $|x^{\overline{G}_{\sigma}}| \le n|x^{G_0}|.$
\end{enumerate} 
\end{Lem}
\begin{proof}
See \cite[4.2.2(j)]{gorlyo} for the proof of $(1)$ and \cite[Lemma 3.20]{fpr2} for $(2)$.
\end{proof}

\begin{Lem}\label{lemxGbound}
Let ${G_0}=\PSL_n(q)$, $x \in G$ have prime order $r$ and $s:=\nu(x).$ Then the following holds.
\begin{equation}\label{5uni}
 |x^G| > 
\begin{cases}
  \frac{1}{2n}  q^{ns} \ge \frac{1}{2n} q^{n^2/2}  & \text{ for } s \ge n/2;\\
  \frac{1}{2n} q^{2s(n-s)} \ge  
\frac{1}{2n}   q^{(3/8)n^2}  & \text{ for }  n/4 \le s < n/2.
\end{cases}
\end{equation}
\end{Lem}
\begin{proof}
The statement follows by Lemma \ref{xGoGs} and \cite[Propositions 3.22 and 3.36, Lemmas 3.34 and 3.38]{fpr2}.
\end{proof}

\subsection{Computations using {\sf GAP} and {\sc Magma}}
\label{gapsec}

We use the  computer algebra systems {\sf GAP} \cite{GAP4} and {\sc Magma} \cite{magma}  to check the inequality $b_S(S \cdot \SL_n(q))\le c$ for  particular cases where $S$ is a maximal solvable subgroup of $\GL_n(q)$ and $q$ and $n$ are small. In this section we discuss how the statement $b_S(G)\le c$ can be checked for given  $S \le \GL_n(q)$ and integer $c$, and how to construct desired $S$  using {\sf GAP} and {\sc Magma}.

Since a nontrivial transitive permutation group on $\Omega$ has base size $c>0$ if and only if the stabiliser of  $\omega \in \Omega$ acting on $\Omega \backslash \{\omega\}$ has base size $c-1$, it is enough to check that the action of the stabiliser of $S$ in $G/S_G$ on $(\Omega \backslash \{S\})^{c-1}$ has a regular point. Here $\Omega$ is the set of right cosets of $S$ in $G$. The following code in {\sc Magma} checks the existence of such an orbit  and returns {\sf true} if $b_S(G) \le c$: 
\lstset{
basicstyle=\ttfamily}  
{\small
\begin{lstlisting}
hom, R:=CosetAction(G,S);
St:=Stabiliser(R,1);
X:=Set(GSet(R));
Exclude(~X,1);
Y:=Set(CartesianPower(X,c-1));
Z:=GSet(St,Y);
O:=Orbits(St,Z);
OrbReg:=[x: x in O| #x eq #St];
#OrbReg gt 0;
\end{lstlisting}}

 This procedure works well when $|G:S|^{c-1}$ is relatively small. For example, if $|G:S|^{c-1}$ is at most $4000$, then this code executes in {\sc Magma} V2.28-4 in at most 12 seconds on a machine with a 2.6 GHz processor. We use this method to check that the base size is larger than $c$. Since for the most cases we verify computationally $b_S(G)=2,$ we only use this code in a couple of cases for very small $|G:S|^c$.
 % (for example, when $S$ is the Singer cycle of $G=\GL_3(2)$ with $|G:S|=8$).
    The way to establish  $b_S(G)\le c$, that we use in most situations, is to find $a_1, \ldots, a_{c-1}$ such that 
\begin{equation}
\label{gapint}
S \cap S^{a_1} \cap \ldots \cap S^{a_{c-1}} =S_G=\cap_{g \in G} S^g.
\end{equation}
Randomised search allows us to find such $a_i$ in all cases we consider. 
% The function {\sf Random}$(G)$, in both {\sf GAP} and {\sc Magma}, allows us to find such $a_i$ in most cases we consider. 

A more difficult  task is to construct all desired (up to conjugacy) maximal solvable subgroups of $G$.  Generating sets for primitive maximal solvable subgroups of $\GL_n(q)$ for small $n$ (and for subgroups containing primitive maximal solvable subgroups of $\GL_n(q)$ in the general case) can be found in \cite{short} and \cite[\S 21]{sup}. The function \texttt{IrreducibleSolvableGroupMS}$(n,p,i)$ in the {\sf GAP} package \texttt{PrimGrp} \cite{PrimGrp} realises the results of \cite{short}. It returns a representative of the $i$-th conjugacy class of irreducible solvable subgroups of $\GL_n(p)$, where $n>1$, $p$ is a prime, and $p^n < 256$.   While constructing a specific subgroup can be difficult, it is usually not necessary:  for $M>S$ if
\begin{equation}
\label{gapintM}
M \cap M^{a_1} \cap \ldots \cap M^{a_{c-1}} =S_G,
\end{equation}
then \eqref{gapint} holds, so it is enough to construct an overgroup $M$ of $S$ such that \eqref{gapintM} holds.   The function \texttt{ClassicalMaximals} in  {\sc Magma} realises the results of \cite{maxlow}. It  constructs all maximal subgroups of a classical group ($\GL_n(q)$ for example) up to conjugation for $n \le 17$. Since the generators of these subgroups are written in the code explicitly in terms of the field elements, there are no practical restrictions on the value of $q$.  If a maximal subgroup $M$ is too big and \eqref{gapintM} does not hold, then we use the function \texttt{MaximalSubgroups} to obtain all maximal subgroups of $M$ up to conjugation and check if \eqref{gapintM} holds for them. In practice, at most three iterations are needed to obtain \eqref{gapintM} for some overgroup of the solvable subgroup under investigation.



We often write ``$b_S(G)<c$ is verified by computation''\index{computation} to imply that the statement is verified using one of the procedures described above.   
Sometimes we carry out similar computation to show  that the intersection of a specific number of conjugates of $S$ satisfies particular order bounds.  We often write ``computation shows'' to summarise such routine calculations.

\section{Irreducible solvable subgroups}
\label{ch2}

Recall from the introduction that $b_S(G)$ is the minimal number such that there exist $$x_1, \ldots, x_{b_S(G)} \in G  \text{ with } S^{x_1} \cap \ldots \cap S^{x_{b_S(G)}}=S_G$$ where $S_G= \cap _{g \in G}S^g$. Let $S$ be an irreducible maximal solvable subgroup of $\GL_n(q)$.  The goal of this section is to obtain upper bounds for $b_S(S \cdot \SL_n(q))$. These bounds play an important role in the proof of Theorems \ref{theorem}, \ref{theoremGR}.  While, with some exceptions,  $b_S(S \cdot \SL_n(q)) \le 4$ follows by Theorem \ref{bernclass}, it is not sufficient for our purposes. In this section we prove that  $b_S(S \cdot \SL_n(q)) \le 2$ with a short list of exceptions.

\subsection{Primitive subgroups} \label{sec31}


    We start our study with a special case:  $S$ is a primitive maximal solvable subgroup  of $\GL_n(q)$ (recall the definition of a primitive linear group after Notation \ref{not}). In Section \ref{impsec} we use these results to obtain bounds for $b_S(S \cdot \SL_n(q))$ where $S$ is irreducible.  

To prove results about $b_S(S \cdot \SL_n(q))$, we need  information about  primitive solvable groups, upper bounds for $|S|$ and lower bound for $\nu(x)$ (the codimension of a largest eigenspace of $x$, see Definition \ref{nudef}), where $x$ is a prime order element of the image of $S$ in $\PGL_n(q).$  This information is needed to apply the probabilistic method described in Section \ref{fprsec}.

The following theorem collects properties of primitive maximal solvable subgroups from \cite[\S \S 19 -- 20]{sup} and \cite[\S 2.5]{short}.  For $A \le \GL_m(q)$ and $B \le \GL_n(q)$, $A \otimes B = \{a \otimes b \mid a \in A, b \in B\}$ where $a \otimes b$ is the Kronecker product defined  after Notation \ref{not}.  

\begin{Th}\label{suplem}
Let $S \le \GL_n(q)$ be a primitive maximal solvable subgroup. Then $S$
admits a unique chain of subgroups 
\begin{equation}\label{ryad}
S \trianglerighteq C \trianglerighteq F \trianglerighteq A
\end{equation}
 where $A$ is the unique maximal abelian normal subgroup of $S$,  $C =C_S(A)$, and  $F$ is the full preimage in $C$ of the maximal  abelian normal subgroup of $S/A$ contained in $C/A.$ The following hold: 
\begin{enumerate}[font=\normalfont]
\item[$(a)$] $A$ is the multiplicative group of a field extension $K$ of the field of scalar matrices $ \Delta = \{\alpha I_n: \alpha \in \mathbb{F}_q\}$ and $m:=|K : \Delta|$ divides  $n$.
\item[$(b)$]  $S/C$ is isomorphic to a subgroup of the Galois group of the extension $K:\Delta,$ so 
$$S=C.\sigma,$$
where $\sigma$ is cyclic of order dividing $m$. 
\item[$(c)$] $C \le \GL_e(K)$ where $e=n/m$.
\item[$(d)$] If
\begin{equation}\label{r}
 e=p_1^{l_1} \cdot \ldots \cdot p_t^{l_t},
\end{equation}
where $t \in \mathbb{N},$ and the $p_i$ are distinct primes, then each $p_i$ divides $|A|=q^m-1.$
\item[$(e)$] $|F/A|=e^2$ and $F/A$ is the direct product of elementary abelian groups. 
\item[$(f)$] If $Q_i$ is the full preimage in $F$ of the Sylow $p_i$-subgroup of $F/A,$ then
$$Q_i= \langle u_1 \rangle \langle v_1 \rangle \ldots \langle u_{l_i} \rangle \langle v_{l_i} \rangle A,$$
where 
\begin{equation}\label{quv}
[u_j,v_j]=\eta_j, \text{ } \eta_j^{p_i}=1, \text{ } \eta_j \ne 1, \text{ } \eta_j \in A; \text{ } u_j^{p_i},v_j^{p_i} \in A
\end{equation}
and elements from distinct pairs $(u_j,v_j)$ commute.
\item[$(g)$] In a suitable $K$-basis of $K^e$, 
$$F=\tilde{Q_1} \otimes \ldots \otimes \tilde{Q_t},$$
where $\tilde{Q_i}$ is an absolutely irreducible subgroup of $\GL_{p_i^{l_i}}(K)$ isomorphic to $Q_i.$
\item[$(h)$] If $N=N_{\GL_e(K)}(F)$ and $N_i=N_{G_i}(\tilde{Q_i}),$ then
 $$N=N_1 \otimes \ldots \otimes N_t.$$
\item[$(i)$] $N_i/\tilde{Q_i}$ is isomorphic to a completely reducible subgroup of $\Sp_{2l_i}(p_i).$
\item[$(j)$] $F/A=C_{C/A}(F/A)$.

\end{enumerate}
\end{Th}

The following result is due to Gluck and Manz {\cite[Theorem 3.5]{manz}}.

\begin{Th}  \label{lowboundsol}
If $S$ is a completely reducible solvable  subgroup of $\GL_n(q)$, then $|S| < q^{9n/4}/2.8.$
\end{Th}


Let $S$ be a primitive maximal solvable  subgroup of $\GL_n(q)$. If $A,$ $F$ and $C$ are as in  \eqref{ryad}, then by Theorem \ref{suplem}
\begin{equation*}
\begin{split}
|A|&=q^m-1;\\
|F:A|&=(n/m)^2=e^2;\\
|S:C|&\le m.\\
\end{split}
\end{equation*}
By  Theorem \ref{suplem} (i) and Theorem \ref{lowboundsol}
\begin{equation}
\label{CFbound}
|C:F| \le 
\begin{cases}
\prod_{i=1}^t |\Sp_{2l_i}(p_i)|;\\
\prod_{i=1}^t ((p_i^{2l_i})^{9/4}/2.8)<e^{9/2},
\end{cases}
\end{equation}
where $e=n/m$, $p_i$ and $l_i$ are as in \eqref{r}. Notice that 
$$|\Sp_{2l_i}(p_i)|=p_i^{l_i^2} \prod_{j=1}^{l_i}(p_i^{2j}-1)\le p_i^{l_i^2} \prod_{j=1}^{l_i}p_i^{2j}=p_i^{l_i(2l_i+1)}. $$
Denote $\log_2(e)$ by $l$, so $l_i \le l$ for $i=1, \ldots, t.$ Therefore,
$$|C:F| \le \prod_{i=1}^t p_i^{l_i(2l+1)}=e^{2l+1} $$ 
and
\begin{equation*}
|S|\le (q^m-1)\frac{n^2}{m^2}m \cdot \min \{ e^{2l+1}, e^{9/2}\}=(q^m-1)m \cdot \min \{e^{2l+1}, e^{13/2}\},
\end{equation*}
so
\begin{equation}\label{H}
|S/Z(\GL_n(q))|\le \left( \frac{q^m-1}{q-1} \right) m\cdot \min \{e^{2l+1}, e^{13/2}\}.
\end{equation}

\begin{Lem}\label{6}
 Let $S$  be a primitive maximal  solvable subgroup of $\GL_n(q)$ and let $H=S/Z(\GL_n(q)) \le \PGL_n(q).$ Let $x \in H$ and let $\hat{x}$ be a preimage of $x$ in $S$. If $x$ has prime order $k$, so  $\hat{x}$ lies in $S \backslash Z(\GL_n(q))$, then %$\nu(x) \ge n/4$. Furthermore, if $|\hat{x}|$ is prime, then 
the following hold:
\begin{enumerate}[font=\normalfont]
\item if $\hat{x} \in A$, then $\nu(x) \ge n/2; $ \label{6it1}
\item if $\hat{x}  \in F \backslash A$, then $\nu(x) \ge n/2; $ \label{6it2}
\item if $\hat{x}  \in C \backslash F$, then $\nu(x) \ge n/4; $ \label{6it3}
\item if $\hat{x}  \in S \backslash C$, then $\nu(x)=n-n/k \ge n/2. $ \label{6it4}
\end{enumerate}
\end{Lem}
\begin{proof}
The proof follows the beginning of the proof of \cite[Proposition 4]{gluck}. Let $\hat{x}$ be a preimage of $x$ in $\GL_n(q),$ so $\hat{x} \in S$ since $Z(\GL_n(q)) \le S$. Observe that $\nu(x)=\nu(\hat{x}).$ Let $\overline{\mathbb{F}_q}$ be the algebraic closure of $\mathbb{F}_q.$

If $\hat{x} \in A$, then, since $A$ is cyclic and $(|A|,q)=1$, $\hat{x}$ is conjugate in $\GL_n(\overline{\mathbb{F}_q})$ to 
$$\diag(\lambda, \lambda^q, \ldots, \lambda^{q^{m_1-1}}, \ldots, \lambda, \lambda^q, \ldots, \lambda^{q^{m_1-1}}); \text{ } \lambda \in \overline{\mathbb{F}_q}$$
by \cite[Lemma 1.3]{buturl}. Here $m_1$ is the smallest possible integer such that  $\lambda^{q^{m_1}}=\lambda.$ Therefore, $\nu(\hat{x})=n-n/m_1\ge n/2.$ Moreover, if $z \in A$ is nontrivial, then $C_V(z)=\{0\}.$

Let $\lambda \in \overline{\mathbb{F}_q}^*.$ If $\hat{x} \in S \backslash C,$ then 
$$[\lambda \hat{x}, z]=[\hat{x},z] \in A \backslash \{1\}$$ for some $z \in A.$ Notice 
$$C_V((\lambda \hat{x})^{-1}) \cap C_V(z^{-1} \lambda \hat{x} z) \subseteq C_V([\lambda \hat{x}, z])=\{0\}$$
and $\dim(C_V((\lambda \hat{x})^{-1}))= \dim(C_V((\lambda \hat{x})^{z})),$ so $ \dim (C_V((\lambda \hat{x}))) \le n/2$ for every $\lambda \in \overline{\mathbb{F}_q}^*$. Hence $\nu(\hat{x})\ge n/2.$
Moreover, if  \ref{prost2} holds in Lemma \ref{prost}, then $\nu(x)=n-n/k.$ 
Suppose \ref{prost1} holds in Lemma \ref{prost}, so $\hat{x} \in S$ has order $k.$
 If $\hat{x}  \in S \backslash C$ then, by $(b)$ of Theorem \ref{suplem} and \cite[7-2]{conjaut}, $\hat{x}$ is conjugate to a field automorphism $\sigma_1 \le \langle \sigma \rangle$ of $\GL_e(K)$ of order $k$ where $K$ is as in Theorem \ref{suplem}$(d)$.
Such an element acts as a permutation with $(n/k)$ $k$-cycles on a suitable basis of $V$, so  $\nu(x)=n-n/k$.



If $\hat{x} \in F \backslash A,$ then, by $(f)$ of Theorem \ref{suplem}, there exists $h \in F \backslash A$ such that $$[ \lambda \hat{x}, h]= [\hat{x},h] \in Z \backslash \{1\}$$ and $\nu(\hat{x}) \ge n/2$ as above.



If $\hat{x} \in C \backslash F$, then $[\lambda \hat{x}, h] \in F \backslash A$ for some $h \in F$ by $(j)$ of Theorem \ref{suplem}. Therefore, 
$$\dim (C_V([\lambda \hat{x}, h])) \le n/2 \text{ and } \dim (C_V(\lambda \hat{x})) \le 3n/4$$
for every $\lambda \in K^*$,
so $\nu(\hat{x}) \ge n/4.$
\end{proof}

\begin{Lem}\label{sinbase}
If $S$ is the normaliser of a Singer cycle of $\GL_n(q)$ with $n\ge 3$, then $b_{S}(S \cdot \SL_n(q)) \le 3$ with equality if and only if $(n,q)=(3,2).$
\end{Lem}
\begin{proof}
We follow \cite[Lemma 6.4]{burPS} which proves the statement of the lemma in the case of prime $n$. Let $\hat{G}$ be $S \cdot \SL_n(q),$ let  $G$ be $\hat{G}/Z(\hat{G})\le \PGL_n(q)$ and let $H$ be $S/Z(\hat{G})\le G.$   Obviously, 
$$b_S(\hat{G})=b_H(G).$$

Let $x\in H$ be an element of prime order. Since $S$ is the normaliser of a Singer cycle $A$, $A=F=C$ in the notation of Theorem \ref{suplem}, so $\nu(x) \ge n/2$ for all such $x \in H$ by Lemma \ref{6}.  Hence
$$|x^G| \ge \frac{1}{2n}q^{n^2/2}=b$$ 
for all $x \in H$ of prime order  by Lemma \ref{lemxGbound}. Let $$a:=|H|=n(q^n-1)/(n-1).$$
Recall $\hat{Q}(G,c)$ from \eqref{ver} and that $b_H(G)\le c$ if  $\hat{Q}(G,c)<1.$ By Lemma \ref{fprAB}, 
$$\hat{Q}(G,2) \le  a^2/b.$$
 It is routine to check that, if $n \ge 4,$ then this upper bound is less than $1$ unless  either $(n,q)=(6,2),(5,3), (5,2)$, or $n=4$ and  $q<13.$ For these cases the lemma is verified by computations. If $n=3$, then the lemma follows directly from \cite[Lemma 6.4]{burPS}.
\end{proof}


\begin{Lem}\label{sch} 
 Let $n \ge 6.$ If $S$ is a primitive maximal solvable  subgroup of $\GL_n(q)$, then $$b_S(S \cdot \SL_n(q)) = 2.$$ 
\end{Lem}
\begin{proof}

Let   $G$ be $(S \cdot \SL_n(q))/Z(\GL_n(q))\le \PGL_n(q)$ and let $H$ be $S/Z(\GL_n(q))\le G$, so 
$$b_S(S \cdot \SL_n(q))=b_H(G).$$

If $n \ge 6$ and for all $x \in G$ of prime order   
$$|x^G \cap H| < |x^G|^{(3c-4)/(3c)},$$
then $b_H(G) \le c$ by  \eqref{fpr}  and \eqref{0}. Therefore,  it suffices to show this inequality for $c=2$. 


Let $s:=\nu(x).$ We use bounds  \eqref{5uni} for $|x^G|.$
In most cases the bound $|x^G \cap H| \le |H|$ is sufficient. 

For $n \ge 16$, the inequality $|x^G \cap H| < |x^G|^{1/3}$ follows for all $x \in H$ or prime order by combining  Lemma \ref{6} and the bounds   in \eqref{5uni} for $|x^G|$ and \eqref{H} for $|x^G \cap H|$ respectively. And there is a finite list of cases with $6 <n \le 15$ where these bounds  are insufficient. By  Theorem \ref{suplem} $(d)$,   $(q^m-1)$ must be divisible by $p_i$ for all $p_i,$ $i=1, \ldots, t.$ Using this statement and the bound for $|C:F|$ obtained in \eqref{CFbound} by using the precise orders of the $\Sp_{2l_i}(p_i)$, we reduce this list to  cases 1--6 in Table~\ref{tab}. Here $e$ is as in Theorem \ref{suplem}. %For $n=6$ there exists an  infinite sequence of such cases.
The lemma follows by Lemma \ref{sinbase} in case $1$ and it is verified by computation  in cases 2--6. 
\begin{table}[h]
\centering
\caption{Exceptional cases in proof of Lemma \ref{sch}}
\label{tab}
\begin{tabular}{|l|l|l|l|}
\hline
\textbf{Case} & $n$ & $e$ & $q$   \\ \hline
\textbf{1}  & 6   & 1   & any   \\ \hline
\textbf{2}  & 6   & 2   & 3     \\ \hline
\textbf{3}  & 6   & 3   & 2,4 \\ \hline
\textbf{4}  & 7   & 1   & 2,3,4 \\ \hline
\textbf{5}  & 8   & 1   & 2     \\ \hline
\textbf{6}  & 8   & 8   & 3,5    \\ \hline
\end{tabular}
\end{table} 
\end{proof}

Before the next lemma, we remind the reader the following definitions.
A subgroup of $\GL_n(q)$ is {\bf absolutely irreducible} if it is irreducible as a subgroup of $\GL_n(K)$, where $K$ is the algebraic closure of $\mathbb{F}_q.$ For a prime $r$, a finite $r$-group $R$ is of {\bf symplectic-type} if every characteristic abelian subgroup of $R$ is cyclic. Symplectic-type groups are closely related to extra-special groups, see \cite[\S 4.6]{kleidlieb} and \cite[\S 2.4]{short} for a summary and details.

\begin{Lem}
\label{c6small}
 Let $S$ be a primitive maximal solvable subgroup of $\GL_n(q).$  Recall that $q=p^f.$ If $e=n=r^l$ for some integer $l$ and prime $r$ in Theorem \ref{suplem}, then  $S=S_1 \cdot Z(\GL_n(q))$ where $S_1=S \cap \GL_n(p^t)$,   $t$ divides $f$, and $S_1$ lies in the normaliser $M$ in $\GL_n(p^t)$ of an absolutely irreducible symplectic-type subgroup of $\GL_n(p^t)$. %So $M$ is a maximal group of $\GL_n(p^t).$  
\end{Lem}
\begin{proof}
Let $F$  be as in Theorem \ref{suplem}. Let $W \le V$ be an irreducible $\mathbb{F}_{q}[F]$-submodule. By  Theorem \ref{suplem} $(f)$, $F=A \cdot F_1$ where $F_1$ is extra-special of order $r^{2l+1},$ so $W$ is a faithful irreducible $\mathbb{F}_{q}[F_1]$-module. Therefore, by \cite[Propositions 4.6.2 and 4.6.3]{kleidlieb}, $\dim W=r^{l}$, and $F_1$ is an absolutely irreducible symplectic-type subgroup of $ \GL_n(p^t)$  where $\mathbb{F}_{p^t}$ is the smallest field (of characteristic $p$) over which such a representation of $F_1$ can be realised. In particular, $t$ is the smallest positive integer for which $p^t \equiv 1 \mod |Z(F_1)|.$ Moreover, by \cite[Corollary 2.4.12]{short}, $|Z(F_1)|$ is either $r$ or $4$.  By  \cite[Theorem 2.4.12]{short}, $S=S_1 \cdot Z(\GL_n(q))$ where $S_1 \le N_{\GL_n(p^t)}(F_1)=M.$     %so $S_1$ lies in a maximal subgroup $M \in \mathcal{C}_6$ of $\hat{G} \cap \GL_n(p^t)$ (see  \cite[\S 4.6]{kleidlieb} for the definition of $\mathcal{C}_6$). 
\end{proof}

Let $S$ be a  primitive maximal solvable subgroup of $\GL_n(q)$ of  degree $n \le 5.$ The equation  $b_S(S \cdot \SL_n(q))=2$ does not always hold for such $n$. However, in view of Lemmas %\ref{irrtog}, \ref{prtoirr} and \ref{diag} for every maximal primitive solvable subgroup $S$ it suffices  to find $x \in \GL(n,q)$  such that
\ref{irrtog}, \ref{diag}   and \ref{prtoirr}, for every primitive maximal  solvable subgroup $S$ it suffices  to find $x \in \SL_n(q)$  such that
\begin{equation*}%\label{uptr}
 S \cap S^x \le RT(\GL_n(q))
\end{equation*}
 to prove that $b_M(M \cdot \SL_n(q)) \le 5$ for every maximal solvable subgroup $M$.  Recall that $D(\GL_n(q))$ and $RT(\GL_n(q))$ denote the subgroups of all diagonal and all upper-triangular matrices in $\GL_n(q)$ respectively. Therefore, if $b_S(S \cdot \SL_n(q))>2$, then we decide if there exists $x \in \SL_n(q)$ such that  $S \cap S^x$ lies in $D(\GL_n(q))$ or $RT(\GL_n(q)).$ In particular, we prove in this section that if $q>7$, then such $x$ always exists, so  $$b_M(M \cdot \SL_n(q)) \le 5$$
for every maximal solvable subgroup $M$ of $\GL_n(q)$, for every $n \ge 2$; see Corollary \ref{corir} and Theorem \ref{irred}. We also use this observation in the proof of Theorem \ref{theorem} in Section \ref{sec411}.


\medskip

The following theorem is the main result of Section \ref{sec31}.

\begin{Th}
\label{sec311lem}
 Let $S$ be a primitive maximal solvable subgroup of $\GL_n(q)$ where $n \ge 2$ and  $(n,q) \ne (2,2), (2,3).$
Then either $b_S(S \cdot \SL_n(q))=2$, or one of the following holds:
\begin{enumerate}[font=\normalfont]
\item $n =2$,  $q>3$ is odd,  $S$ is the normaliser of a Singer cycle and  $b_S(S \cdot \SL_2(q))=3$. If $q>5$, then there exists $x\in \SL_2(q)$ such that  $S \cap S^x \le D(\GL_2(q));$ \label{thirr2odd}
\item $n =2$,  $q\ge 4$ is even,  $S$ is the normaliser of a Singer cycle and  $b_S(S \cdot \SL_2(q))=3$. In this case there exists $x\in \SL_2(q)$ such that  $S \cap S^x \le RT(\GL_2(q));$ \label{thirr2even}
\item $n=2,$ $q=9,$ $S$ is generated by all matrices with entries in $\mathbb{F}_3$ and scalar matrices, so $S=\GL_2(3) \cdot Z(\GL_2(9)),$ and  $b_S(S \cdot \SL_2(9))=3$.  In this case there exists $x\in \SL_2(9)$ such that  $S \cap S^x \le RT(\GL_2(9));$ \label{thirr29}
\item $n=2$, $q \in \{5,7\}$, and $S$  is an absolutely irreducible subgroup such that $S/Z(\GL_2(q))$ is isomorphic to $2^2.\Sp_2(2)$. Here  $b_S(S \cdot \SL_2(q))$ is $4$ and $3$ for $q$ equal to $5$ and $7$ respectively;
\item $n=3$, $q=2$,   $S$ is the normaliser of a Singer cycle and  $b_S(S \cdot \SL_3(2))=3$. \label{thirr32} 
\end{enumerate}  
\end{Th}
\begin{proof}
For $n \ge 6$ the statement holds by Lemma \ref{sch}. Assume $2 \le n \le 5$. First, we prove the statements for the base sizes and then we prove additional assertions of \ref{thirr2odd} and \ref{thirr2even}; part \ref{thirr29} follows by computations.

Let $m$ and $e$ be as in Theorem \ref{suplem}. Notice that, for $n=2,3,5$, either $m=n$ or $m=1,$ since $n$ is prime. If $n=4,$ then there is additional  case $m=e=2.$

\medskip

{\bf Case 1:} $m=n.$

 If $m=n$, then $S$ is the normaliser of a Singer cycle and the statement of the lemma on $b_S(S \cdot \SL_n(q))$ follows by Lemma \ref{sinbase} for $n \ge 3$ and by \cite[Proposition 4.1]{burPS} for $n=2$.

\medskip

{\bf Case 2:} $m=1.$ 

Here  Lemma \ref{c6small} holds, since $n$ is either prime or $4=2^2$. Hence  $S= S_1 \cdot Z(\GL_2(q))$ where $S_1$ lies in the normaliser  $M$ of an absolutely irreducible symplectic-type subgroup of $\GL_n(p^t)$ where $t$ is minimal such that $p^t \equiv 1 \mod |Z(F_1)|$ and $F_1$ is as in the proof of Lemma \ref{c6small}. If  $t=1$, then $M$ (up to scalars) is a maximal subgroup of either $\GL_n(p)$ or $\Sp_{n}(q)$ (if $n$ is even) from the Achbacher's class $\mathit{C}_6$ by \cite[Table 4.6.B]{kleidlieb}. Bounds for base size for primitive actions of  classical groups with $n \le 5$ are listed in \cite[Tables 2 and 3]{burness}. In particular, it follows that there exists $g \in \SL_n(p)$ such that $M \cap M^g \le Z(\GL_n(q))$ unless $n=2,4$ and $p \le 7.$ Hence  $S \cap S^g \le Z(\GL_n(q))$ and  $b_S(S \cdot \SL_n(q))=2$.  The results for $n=2,4$ and $q \le 7$ are verified by computation.

For $t >1$ we prove the statement directly using the probabilistic method. Let   $G$  be $(S \cdot \SL_n(q))/Z(\GL_n(q))\le \PGL_n(q)$ and let $H$ be $S/Z(\GL_n(q))\le G$, so 
$$b_S(S \cdot \SL_n(q))=b_H(G).$$ 
 By Theorem \ref{suplem}, $H$ lies in a subgroup of $\PGL_n(q)$ isomorphic to $e^{2}.\Sp_{2l}(p_1^l)$, where $e=p_1^l$ for a prime $p_1$. The number of elements of order equal to each  prime dividing this subgroup can be computed directly. Further,
  using bounds for $|x^G|$ for $x$ of prime order from Lemma \ref{lemxGbound}, we obtain an upper bound for $\widehat{Q}(G,2)$ via Lemma \ref{fprAB}. For example, if $n=5$, then
  there are   $25$ elements of order $2$, $500$ elements of order $3$ and $624$ elements of order $5$, so 
$$\widehat{Q}(G,2)\le \frac{(25^2+500^2+624^2)}{(1/10)q^{12}}.$$ 
 This upper bounds are less than $1$ for $q$ greater than $17,$ $17,$ $19$ and $7$ for $n=2,3,4$ and $5$ respectively. For smaller $q$, the value of $b_H(G)$ is verified using computations. In particular, if $q=9$ then \ref{thirr29} occurs.   
 
\medskip 
 
{\bf Case 3:} $n=4,m=2.$

 Since $e=n/m=2$, $q$ is odd by Theorem \ref{suplem}$(d)$. Recall from Theorem \ref{suplem} that $S$ has a chain of normal subgroups:
$$S \ge C \ge F \ge A,$$
where $C \le \GL_2(q^2),$ so $S \le \GL_2(q^2).2,$ and $C/A$ is isomorphic to a subgroup of $2^2.\Sp_2(2).$ 
Therefore, $|S| $ divides $ 48(q^2-1)$ and $|H|$ divides $a=48(q+1).$ 
 
 For all elements of prime order in $H$,
 $$|x^G| \ge (1/8)q^6=b$$ by Lemma \ref{lemxGbound}. Hence, by Lemma \ref{fprAB},
 $\hat{Q}(G,2) \le  a^2/b$.  This upper bound is less than $1$ unless  $q \ge 13.$
If $q \le 11$, then  $b_S(S \cdot \SL_4(q))=2$  is verified  by computation.

This completes the proof of the values of $b_S(S \cdot \SL_n(q)).$



\medskip

%{\bf Case 1:} $n=2$. 

%  Let $m$ and $e$ be as in Theorem \ref{suplem}. So, either $m=2$ and $S$ is the normaliser of a Singer cycle, or $m=2,$ so $e=n=2$ and $S$  is as in Lemma \ref{c6small}. In the  second case, which arises only if $q$ is odd by Theorem \ref{suplem}$(d)$, $S= S_1 \cdot Z(\GL_2(q))$ where $S_1$ lies in a maximal subgroup $M$(see the third line of \cite[Table 4.6.B]{kleidlieb}) of $\GL_2(p)$ and $S/Z(\GL_n(q))$ is isomorphic to $2^2.\Sp_2(2).$ By \cite[Lemma 4.4]{burness}, there exists $g \in \SL_2(p)$ such that $M \cap M^g \le Z(\GL_2(p))$ for $p \ge 7.$ Hence  $S \cap S^g \le Z(\GL_2(q))$ and  $b_S(S \cdot \SL_2(q))=2$ if $q > 7$. The results for $q \in \{5,7\}$ are verified by computation. 


 Let us prove the remaining claims in parts $(1)$ and $(2)$.
 First, assume $q$ is odd and fix a non-square $a \in \mathbb{F}_q$ (there are $(q-1)/2$ such elements). Consider 
\begin{equation}\label{thesin}
 S_a= \left\{
\alpha \begin{pmatrix}
1    & 0  \\
      0    &   1       
\end{pmatrix} + 
\beta \begin{pmatrix}
0    & 1  \\
a    &  0       
\end{pmatrix}
\right\} \backslash \{0\}.
\end{equation}
 Notice that for an element of $S_a$
$$\det \begin{pmatrix}
\alpha    & \beta  \\
    a  \beta    & \alpha       
\end{pmatrix} =0$$ if, and only if, $a=(\alpha/ \beta)^2,$ so all matrices in $S_a$ are invertible. 
 Calculations show that $S_a$ is an abelian subgroup of $\GL_2(q)$ of order $q^2-1$ and $S_a \cup \{0\}$ under usual matrix addition and multiplication is a field, so $S_a$ is a Singer cycle. Notice that 
$$\varphi = \begin{pmatrix}
-1    & 0  \\
0    &  1       
\end{pmatrix}$$
normalises $S_a.$ Therefore, we can view $S$ as   $N_{\GL_2(q)}(S_a)=S_a \rtimes \langle \varphi \rangle$. 

 It is easy to see that if there exist $a, b \in \mathbb{F}_q$ such that $a \ne b$
and neither $a$ nor $b$ has square roots in $\mathbb{F}_q$, then $$S_a \cap S_b \le Z(\GL_2(q)).$$ 
If, in addition, $a \ne -b$, then the direct comparison of the matrices show that 
\begin{equation}
\label{2sindiagodd}
 N_{\GL_2(q)}(S_a) \cap N_{\GL_2(q)}(S_b) =\langle \varphi \rangle Z(\GL_2(q)) \le D(\GL_2(q)).
\end{equation} 
It is possible to find such $a$ and $b$ if $q>5,$ so in this case there exists $x \in \SL_2(q)$ such that $$S \cap S^x \le D(\GL_2(q)),$$ since all Singer cycles are conjugate in $\GL_2(q)$ and $\Det(S_a)=\Det(\GL_2(q)).$

Let $q$ be even and let $a \in \mathbb{F}_q $ be such that there are no roots of $x^2+x+a$ in $\mathbb{F}_q$ (there are $q/2$ such elements in $\mathbb{F}_q$). 
 Consider  
\begin{equation} \label{thesineven}
S_a= \left\{
\alpha \begin{pmatrix}
1    & 0  \\
      0    &   1       
\end{pmatrix} + 
\beta \begin{pmatrix}
0    & 1  \\
a    &  1       
\end{pmatrix}
\right\} \backslash \{0\}.
\end{equation} Notice that  for an element of $S_a$
$$\det \begin{pmatrix}
\alpha    & \beta  \\
    a  \beta    & \alpha +\beta      
\end{pmatrix} =0$$ if, and only if, $a=(\alpha/ \beta)^2 +(\alpha/\beta),$ so all matrices in $S_a$ are invertible. 
 Calculations show that $S_a$ is an abelian subgroup of $\GL_2(q)$ of order $q^2-1$ and $S_a \cup \{0\}$ under usual matrix addition and multiplication is a field, so $S_a$ is a Singer cycle. Notice that 
$$\varphi = \begin{pmatrix}
1    & 0  \\
1    &  1       
\end{pmatrix}$$
normalises $S_a.$  Therefore, we can view $S$ as   $N_{\GL_2(q)}(S_a)=S_a \rtimes \langle \varphi \rangle$.

 It is easy to see that if there exist $a, b \in \mathbb{F}_q$ such that $a \ne b$
and neither $x^2+x+a$ nor $x^2+x+b$ has roots in $\mathbb{F}_q$, then $$S_a \cap S_b \le Z(\GL_2(q)).$$
Calculations show that  
\begin{equation}
\label{2sindiageven}
N_{\GL_2(q)}(S_a) \cap N_{\GL_2(q)}(S_b)=\langle \varphi \rangle Z(\GL_2(q))
\end{equation} consists of  lower triangular matrices. It is possible to find such $a$ and $b$ if $q \ge 4,$ so in this case there exists $x \in \SL_2(q)$ such that all matrices in $S \cap S^x$ are  lower triangular,   since all Singer cycles are conjugate in $\GL_2(q)$ and $\Det(S_a)=\Det(\GL_2(q)).$
\end{proof}


\subsection{Imprimitive irreducible subgroups}\label{impsec}

We commence by obtaining  results about the  subgroups of wreath products in $\GL_n(q).$ We combine these results with those of the previous section to obtain 
an upper bound to $b_S(S \cdot \SL_n(q))$ for imprimitive maximal solvable subgroups $S$ of $\GL_n(q).$   

\medskip

Let $A(n)$  be the following $n \times n$ matrix:
\begin{equation}\label{igrek}
\begin{aligned}
A(n) & = \begin{pmatrix}
1      & -1     & 1      & -1     & \ldots & (-1)^{n+1}         \\
0      &   1    & -1     & 1      & \ldots &  (-1)^{n}         \\
0      &    0   & 1      & -1     & \ldots &  (-1)^{n-1}        \\
       &        &        & \ddots &\ddots  &          \\
       &        &        &        & 1      & -1        \\
       &        &        &        & 0      &  1      
\end{pmatrix}, \text{ so } \\
A(n)^{-1} & =
\begin{pmatrix}
1      & 1      & 0      & 0      &        &          \\
0      &   1    & 1      & 0      &        &           \\
0      &    0   & 1      & 1      &        &          \\
       &        &        & \ddots &\ddots  &          \\
       &        &        &        & 1      & 1        \\
       &        &        &        & 0      &  1      
\end{pmatrix}.
\end{aligned}
\end{equation}


\begin{Lem}\label{igrekl}
Let $H \le X \wr Y,$ where $X \le \GL_m(q),$ $Y \le \Sym(k).$  
Let $A(k)=(y_{ij})$ be as in \eqref{igrek} and let $x_i$ for $i =1, \ldots, k$ be arbitrary elements of $X$.
Define $x \in \GL_{mk}(q)$ to be
\begin{gather*}
\diag(x_1, \ldots, x_k) (I_m \otimes A(k))=
\begin{pmatrix}
y_{11}x_1      & y_{12}x_1      & \ldots     & y_{1k}x_1            \\
y_{21}x_2      &  y_{22}x_2     & \ldots     & y_{2k}x_2             \\
\vdots      &             &            & \vdots               \\
y_{k 1}x_{k}   & y_{k 2}x_{k}   &  \ldots    & y_{k k}x_{k}               
\end{pmatrix},
\end{gather*} 
 Let $h=\diag[D_1, \ldots, D_k] \cdot s  \in H$, where $D_i \in X$ and $s \in Y,$ so $h$ is obtained from the permutation matrix $s$ by replacing $1$ in the $j$-th row by the $(m \times m)$ matrix $D_j$ for $j =1, \ldots, k$ and replacing each zero by an $(m\times m)$ zero matrix. 
If $h^x \in H$, then $s$ is trivial and $D_j^{x_j}=D_{j+1}^{x_{j+1}}$ for $j=1, \ldots, k-1.$ 
\end{Lem}
\begin{proof}
 %We say that rows with numbers $(j-1)m+1,  \ldots, jm$ of a matrix form its $j$th $m \times m$-row.
   If $s$ does not stabilise the point $k$, then there is more than one non-zero $(m \times m)$-block in the last $(m \times m)$-row of $h^{x}$ and, thus, $h^{x}$ does not lie in   $ X \wr Y.$
Assume that $s$ stabilises the last ${k-j}$ points and $(j)s=i \le j.$ The  $j$-th $(m \times m)$-row of $h^{x}$ is 
\begin{equation*}	
\begin{split}
\bigl( \overbrace{0, \ldots, 0}^{i-1},  \overbrace{{x_j}^{-1}(D_{j}){x_i},-{x_j}^{-1}(D_{j}){x_i},  \ldots,(-1)^{j-i} {x_j}^{-1}(D_{j}){x_i}}^{j+1-i}, \ldots \bigr).
\end{split}
\end{equation*}
Therefore, $h^x$ does not lie in $X \wr Y$ if $i \ne j,$ since the $j$-th  $(m \times m)$-row contains more than one non-zero $(m \times m)$-block in that case. So $i=j$ and the  $j$-th $(m \times m)$-row of $h^{x}$ is
\begin{equation*}
\begin{split}
\bigl( \overbrace{0, \ldots, 0}^{i-1},  (D_{j})^{x_j},   
( ( D_{j+1})^{x_{j+1}} -(D_{j})^{x_j}  ) ,  \ldots, (-1)^{k-j+1}((D_{j+1})^{x_{j+1}} - (D_{j})^{x_j}  ) \bigr).
\end{split}
\end{equation*}
%with $(i-1)$ zeros in the beginning.
 So, if $h^x \in H$, then $s$ stabilises $j$ and $D_{j}^{{x_j}}=D_{j+1}^{x_{j+1}}.$
\end{proof}

%If $x \in GL(n_1, q)$, $y \in GL(n_2,q)$ then 
%\begin{gather}
%x \otimes y:= 
%\begin{pmatrix}
%y_{11}x      & y_{12}x      & \ldots     & y_{1n_2}x            \\
%y_{21}x      &  y_{22}x     & \ldots     & y_{2n_2}x             \\
%\vdots      &             &            & \vdots               \\
%y_{n_2 1}x   & y_{n_2 2}x   &  \ldots    & y_{n_2 n_2}x               
%\end{pmatrix}.
%\end{gather} 
 

\begin{Lem}\label{prtoirr}
Let $T(\GL_n(q))$ be one of the following subgroups: $Z(\GL_n(q))$, $D(\GL_n(q))$ or $RT(\GL_n(q))$. Let $H$ be a subgroup of $\GL_n(q)$ such that
$$H \le H_1 \wr \Sym(k), $$
where $n=mk$, $H_1 \le \GL_m(q)$. If there exist $g_1, \ldots, g_b \in \GL_m(q)$ ({respectively } $\SL_m(q))$ such that 
$$H_1 \cap H_1^{g_1} \cap \ldots \cap H_1^{g_b} \le T(\GL_m(q)),$$
then there exist   $x_1, \ldots, x_b \in \GL_n(q)$ ({respectively} $\SL_n(q))$ such that 
\begin{equation}\label{irreq}
H \cap H^{x_1} \cap \ldots \cap H^{x_b} \le T(\GL_n(q)).
\end{equation}
\end{Lem}
\begin{proof}
Define $x_i$ to be 
\begin{gather*}
\begin{pmatrix}
y_{11}g_{i}      & y_{1 2}g_{i }      & \ldots     & y_{1 k}g_{i }            \\
y_{21}g_{i}      &  y_{22}g_{i}     & \ldots     & y_{2k}g_{i}             \\
\vdots      &             &            & \vdots               \\
y_{k 1}g_{i }   & y_{k 2}g_{i}   &  \ldots    & y_{kk}g_{i}               
\end{pmatrix},
\end{gather*} 
 where $y=(y_{ij})=A(n).$ Let us show that \eqref{irreq} holds for such $x_i.$  Let  $h=\diag[D_1, \ldots, D_k] \cdot s  \in H$, where $D_i \in H_1$ and $s \in \Sym(k).$
If $$h \in H \cap H^{x_1} \cap \ldots \cap H^{x_b}$$  then, by  Lemma \ref{igrekl}, $s$ is trivial and $D_i=D_j$ for all $1 \le i, j \le {k}$. Thus, $$D_i \in H_1 \cap H_1^{g_1} \cap \ldots \cap H_1^{g_b} \le T(\GL_m(q))$$
and $h \in  T(\GL_n(q)).$ 

Notice that $\det(y)=1$ and if $\det(g_i)=1,$ then 
\begin{equation*}
\det(x_i)=\det(g_i \otimes y)=\det(g_i)^k \cdot \det({y})^m=1. \qedhere
\end{equation*} 
\end{proof}

\begin{Cor}\label{corir}
Let $S$ be a maximal solvable subgroup of $\GL_n(q)$ and assume that matrices in $S$ have shape \eqref{stup}; so $S_i =\gamma_i(S)= P_{i} \wr \Gamma_i$,
where $P_{i}$ is a primitive solvable subgroup of $\GL_{m_i}(q)$, $\Gamma_i$ is a transitive solvable
subgroup of the symmetric group $\Sym(k_i)$, and $k_im_i = n_i$. If, for every $P_{i}$, there exist $x_i \in \SL_{m_i}(q)$ such that
\begin{equation*}%\label{cor29}
P_{i} \cap P_{i}^{x_i}\le RT(\GL_{m_i}(q)),
\end{equation*}
then $b_S(S \cdot \SL_n(q))\le5.$
\end{Cor}
\begin{proof}
The statement follows  from Lemmas \ref{irrtog},  \ref{prtoirr} and  \ref{diag}.
\end{proof}

%\begin{Rem}
%Since $RT(GL(m_i,q))$ is not normal in $GL(m_i,q)$ when $m_i \ge 2$, it may be  that $P_{i} \cap P_{i}^{x_i}$ does not lie in $RT(GL(m_i,q))$ but is conjugate to a subgroup of $RT(GL(m_i,q))$. Since the conclusion of the corollary holds for every conjugate of such subgroup, we may assume  % So $S$ is conjugate to the subgroup which satisfies the condition in the corollary and the conclusion is true for $S$. Therefore, we  assume in such cases 
%that $S$ satisfies \eqref{cor29}.
%\end{Rem}

\begin{Th}\label{irred}
Let $S$ be an irreducible maximal solvable  subgroup of $\GL_n(q)$ with $n \ge 2$, and $(n,q)$ is neither  $(2,2)$ nor $(2,3)$. Then either $b_S(S \cdot \SL_n(q))=2$ or one of the following holds: 
\begin{enumerate}[font=\normalfont]
\item $n =2$,  $q>3$ is odd, $S$ is the normaliser of a Singer cycle and $b_S(S \cdot \SL_2(q))=3$. If $q>5$, then there exists $x\in \SL_2(q)$ such that  $S \cap S^x \le D(\GL_2(q));$ \label{irred11}
\item $n =2$,  $q\ge 4$ is even,  $S$ is the normaliser of a Singer cycle and $b_S(S \cdot \SL_2(q))=3$. In this case there exists $x\in \SL_2(q)$ such that  $S \cap S^x \le RT(\GL_2(q));$ \label{irred12}
\item $n=2,$ $q=9,$ $S$ (up to conjugacy) is generated by all matrices with entries in $\mathbb{F}_3$ and scalar matrices, so $S=\GL_2(3) \cdot Z(\GL_2(9)),$ and $b_S(S \cdot \SL_2(9))=3$. In this case there exists $x \in \SL_2(9)$ such that $S \cap S^x \le RT(\GL_2(9));$ \label{irred135}
\item $n=2$, $q \in \{5,7\}$, and $S$  is an absolutely irreducible subgroup such that $S/Z(\GL_2(q))$ is isomorphic to $2^2.\Sp_2(2).$  Here  $b_S(S \cdot \SL_2(q))$ is $4$ and $3$ for $q$ equal to $5$ and $7$ respectively;\label{irred13}
\item $n=3$, $q=2$,  $S$ is the normaliser of a Singer cycle and $b_S(S \cdot \SL_3(2))=3$; \label{irred14} 
\item $n=4$, $q=3$, $S=\GL_2(3) \wr \Sym(2)$ and $b_S(S \cdot \SL_4(3))=3$. In this case there exists $x \in \SL_4(3)$ such that $S \cap S^x \le RT(\GL_4(3)).$ \label{irred15} %(v) $n=4$, $q \ge 9$ is odd, and $S$ is conjugate to a subgroup of $GL(2,q^2).2.$ In this case there exists $x\in GL(n,q)$ such that  $S \cap S^x \le D(GL(n,q)).$\\ %(vi)  $n=4$, $q\ge 5$ is odd, and $S=S_1 \wr Sym(2)$, where $S_1$ is the normalizer of a Singer cycle of $GL(2,q)$. In this case there exists $x\in GL(n,q)$ such that  $S \cap S^x \le D(GL(n,q)).$\\
\end{enumerate} 
\end{Th}
\begin{proof}
Let $S$ be an  irreducible maximal solvable subgroup of $\GL_n(q).$ The statement  follows by Lemmas \ref{supirr} and \ref{prtoirr} and Theorem \ref{sec311lem}
%there exists $x$ such that 
%\begin{equation}\label{ut}
%S \cap S^x \le RT(GL(n,q))
%\end{equation}
 for all cases except the case when $S$ is  conjugate to $S_1 \wr \Gamma$, where 
$n=kl$, $\Gamma$ is a transitive maximal  solvable subgroup of $\Sym(l)$, and $S_1$ is  one of the following groups:
\begin{enumerate}[label=\alph*)]
\item $k=3, q=2$ and $S_1$ is the normaliser of a Singer cycle of $\GL_3(2);$ \label{irred1}
\item $k=2$, $q=2,3$ and $S_1=\GL_2(q);$ 
\item $k=2$, $q=3$ and $S_1=\GL_2(3)$; 
\item $k=2$, $q \in \{5,7\}$ and $S_1/Z(\GL_2(q))$ is isomorphic to $2^2.\Sp_2(2);$
\item $k=2$, $q=9$ and $S_1$ is $\GL_2(3) \cdot Z(\GL_2(9));$ \label{irred4}
\item $k=2$, $q>3$ is odd and $S_1$ is the normaliser of a Singer cycle of $\GL_2(q);$ \label{irred6}
\item $k=2$, $q>2$ is even and $S_1$ is the normaliser of a Singer cycle of $\GL_2(q).$ \label{irred5}

\end{enumerate}
If $l \ge 3$, then $b_S(S \cdot \SL_n(q))=2$ by \cite[Theorem 3.1]{james}. % Also by the remark after \cite[Theorem 1.4]{james}, if $l=2$ and $q$ is even then $b_S(\GL_n(q))=2$ for all such $S,$ in particular for case \ref{irred5}.

We verify by computation that $b_S(S \cdot \SL_n(q))=2$ for $l=2$ for all cases \ref{irred1} -- \ref{irred4} except  $S_1=\GL_2(3)$. If $S_1=\GL_2(3)$, then there exists $x \in \SL_4(3)$ such that $S \cap S^x \le RT(\GL_4(3))$. 
%Therefore, the only cases when there is no  $x$ such that \eqref{ut} holds are cases 1) -- 4) with $l=1$, so $S=S_1$ is primitive. 
%\medskip

Consider  case \ref{irred6}: so we assume $S=S_1 \wr \Sym(2)\ \le \GL_4(q)$, where $S_1$ is the normaliser of a Singer cycle $S_a$ in $\GL_2(q)$ and $q>3$ is odd, as in \eqref{thesin}. Hence $S \cdot \SL_n(q)=\GL_n(q)$ by Lemma \ref{supSL} since the determinant of a generator of a Singer cycle generates $\mathbb{F}_q^*$. 

Let $$s=
\begin{pmatrix}
0 & 0& 1 & 0\\
0 & 0& 0 & 1\\
1 & 0& 0 & 0\\
0 & 1& 0 & 0
\end{pmatrix},
$$  
so  $g \in S$ has shape 
$$
\begin{pmatrix}
\alpha_1' & \beta_1& 0 & 0\\
a \beta_1' &  \alpha_1 & 0 & 0 \\
0 & 0& \alpha_2' & \beta_2\\
0 & 0& a \beta_2' & \alpha_2
\end{pmatrix} \cdot s^i,
$$
where $i =0,1,$ $\alpha_j' = \pm \alpha_j$, $\beta_j' = \pm \beta_j$ and $\alpha_j'\beta_j' = \alpha_j \beta_j.$ Consider $g^y$ where $y=A(n)$ is as in \eqref{igrek}. First let $i=0,$ so 
$$g^y = \Scale[0.95]{
\begin{pmatrix}
\alpha_1' +a \beta_1' &\alpha_1+ \beta_1 - \alpha_1' -a\beta_1' & \alpha_1' +a\beta_1'- \alpha_1 - \beta_1 & \alpha_1 + \beta_1 -\alpha_1' -a\beta_1'\\
a \beta_1' &  \alpha_1 -a\beta_1' & a\beta_1' - \alpha_1 + \alpha_2' & \alpha_1 + \beta_2 -a\beta_1' - \alpha_2' \\
0 & 0& \alpha_2' +a\beta_2' & \alpha_2+\beta_2-\alpha_2' -a\beta_2'\\
0 & 0& a \beta_2' & \alpha_2 - a\beta_2'
\end{pmatrix}}.
$$ 
Assume that $g^y \in S,$ so ${g^y}_{1,3}={g^y}_{1,4}={g^y}_{2,3}={g^y}_{2,4}=0.$ Thus, $0={g^y}_{2,3}+{g^y}_{2,4}=\beta_2.$ Also, ${g^y}_{1,2}= - {g^y}_{1,3} = 0$, but, since the left upper $(2 \times 2)$ block must lie in $N_{\GL_2(q)}(S_a)$, 
$$({g^y}_{1,2})a = \pm {g^y}_{2,1},$$
so $0 ={g^y}_{2,1} = a\beta_1'$ and $\beta_1=0.$ Therefore, ${g^y}_{1,2}= \alpha_1 -\alpha_1',$ so
$$\alpha_1 =\alpha_1'.$$ Now, since $\beta_2=0,$ ${g^y}_{3,4} = \pm ({g^y}_{4,3}) a^{-1}=0,$  
$$\alpha_2=\alpha_2'$$
and since ${g^y}_{2,4}=0$ we obtain $\alpha_1=\alpha_2,$ so $g^y$ is a scalar.  

Now let $i=1$, so 
$$g^y = \Scale[0.95]{
\begin{pmatrix}
0 & 0 & \alpha_1' + a\beta_1'& \alpha_1 +\beta_1 - \alpha_1' - a\beta_1'\\
\alpha_2' &  \beta_2'-\alpha_2' & a\beta_1' - \beta_2' + \alpha_2' & \alpha_1 - a\beta_1' +a\beta_2 - \alpha_2' \\
a \beta_2' + \alpha_2' & \alpha_2 + \beta_2 -a\beta_2' -\alpha_2'& a\beta_2' + \alpha_2 - \alpha_2 -\beta_2 & \alpha_2+ \beta_2 - a\beta_2' \alpha_2'\\
a\beta_2' & \alpha_2 -a \beta_2'& a\beta_2' - \alpha_2 & \alpha_2 - a\beta_2'
\end{pmatrix}}.
$$ 
If $g^y \in  S$, then ${g^y}_{2,1}={g^y}_{2,2}=0.$ So $\alpha_2 = \beta_2= 0$ which  contradicts the invertibility of $g.$
Therefore, $S \cap S^y  \le Z(\GL_4(q)).$ 

Consider \ref{irred5},  so we  assume $S=S_1 \wr \Sym(2)\ \le \GL_4(q)$, where $S_1$ is the normaliser of a Singer cycle $S_a$ in $\GL_2(q)$ as in \eqref{thesineven}. Since $q>2$, we can choose $a\ne 1$ and arguments similar to those in case \ref{irred6} show that $$S \cap S^y \le Z(\GL_n(q)).$$ 
\end{proof}

We conclude the section with a technical lemma about the normalisers in $\GaL_n(q)$ of some specific solvable subgroups of $\GL_n(q)$ that is used in the following sections.

\begin{Lem} \label{lemRTdiagfield} \phantom{sdfs}
\begin{enumerate}[font=\normalfont]
\item Let $q=p^f$ with $f>1$ and let $S$ be the normaliser of the Singer cycle $S_a$ defined in \eqref{thesin} and \eqref{thesineven} for odd and even $q$ respectively. Then there exists $x \in \SL_2(q)$ such that $$S \cap S^x \le RT(\GL_2(q)).$$ \label{lrtdf1}
\item Let $S$ be $\GL_2(3) \cdot Z(\GL_2(9)) \rtimes \langle \phi \rangle$ where $\phi$ is the standard field automorphism as in \eqref{defphibet}. Then there exist $x, y \in \SL_2(9)$ such that 
$$S \cap S^x \le RT(\GL_2(9))\rtimes \langle \phi \rangle \text{ and } (S \cap S^x) \cap (S \cap S^x)^y \le Z(\GL_2(9)) \rtimes \langle \phi \rangle.$$ \label{lrtdf2}
\item Let $S$ be the normaliser of the Singer cycle of $\GL_3(2)$ generated by the matrix 
$$ \begin{pmatrix}
0    & 0 & 1  \\
1 &0    & 0  \\
  0&    1    &   1      
\end{pmatrix}.$$ Then there exists $x \in \SL_2(q)$ such that 
$$ S\cap S^x = \left\langle
 \begin{pmatrix}
1    & 1 & 1  \\
0 &1    & 0  \\
  1&    0    &   0      
\end{pmatrix} 
\right\rangle. $$ \label{lrtdf3}
\end{enumerate}
\end{Lem}
\begin{proof}
\ref{lrtdf1} Let $q$ be odd. It is routine to check that 
$$N_{\GaL_2(q)}(S_a) = S_a \rtimes \left\langle \phi \begin{pmatrix}
a^{-(p-1)/2}    & 0  \\
0    &  1       
\end{pmatrix} \right\rangle$$ 
where $\phi : \lambda v_i \mapsto \lambda^p v_i$ for $\lambda \in \mathbb{F}_q$ and $\{v_1, v_2\}$ is the basis of $V$ with respect to which matrices from $S_a$ have shape \eqref{thesin}. Since $f>1$, $q>5$ and there exist non-squares $a, b \in \mathbb{F}_q$ such that $a \ne \pm b$. For such $a,b$, the direct comparison of the elements shows that  
$$N_{\GaL_2(q)}(S_a) \cap N_{\GaL_2(q)}(S_b)=\langle \varphi \rangle Z(\GL_2(q)) \le D(\GL_2(q)).$$

\medskip

Let $q$ be even.  It is routine to check that  $N_{\GaL_2(q)}(S_a)$ is solvable and 
$$N_{\GaL_2(q)}(S_a)= (S_a) \rtimes \left\langle \phi \begin{pmatrix}
a^{-1}    & 0  \\
1    &  a^1       
\end{pmatrix} \right\rangle$$ 
where $\phi : \lambda v_i \mapsto \lambda^p v_i$ for $\lambda \in \mathbb{F}_q$ and $\{v_1, v_2\}$ is the basis of $V$ with respect  to which matrices from $S_a$ have shape \eqref{thesineven}. Since $q \ge 4,$ there exist $a, b \in \mathbb{F}_q$ such that $a \ne b$
and neither $x^2+x+a$ nor $x^2+x+b$ has roots in $\mathbb{F}_q$. 
For such $a,b$, the direct comparison of the elements shows that   $$N_{\GaL_2(q)}(S_a) \cap N_{\GaL_2(q)}(S_b)=\langle \varphi \rangle Z(\GL_2(q))\le RT(\GL_2(q)).$$

\medskip

\ref{lrtdf2}
Let $\omega$ be a primitive element of $\mathbb{F}_9$ such that $\omega^2-\omega-1=0.$ Let
$$x = \begin{pmatrix}
\omega^{-1} & \omega^2 \\
0 & \omega
\end{pmatrix} \text{ and } y=\begin{pmatrix}
0 & -1 \\
1 & 0
\end{pmatrix}.$$ It is routine to check the statement of the lemma directly by hand or computationally. 

\medskip

\ref{lrtdf3} 
Computation shows that 
$$ S\cap S^x = \left\langle
 \begin{pmatrix}
1    & 1 & 1  \\
0 &1    & 0  \\
  1&    0    &   0      
\end{pmatrix} 
\right\rangle $$
has order 3 where $$
x= \begin{pmatrix}
1    & 0 & 1  \\
1 &1    & 1  \\
  0&    0    &   1      
\end{pmatrix}. $$ 
\end{proof}

\section{Solvable subgroups of $\GaL_n(q)$}
\label{sec411}

We prove Theorems \ref{theorem} and \ref{theoremGR} in Sections \ref{sec411} and \ref{sec412}, respectively.
In this section $S$ is a maximal solvable subgroup of $\GaL_n(q),$ $G=S \cdot \SL_n(q)$, $H=S/Z(\GL_n(q))$ and $\overline{G}=G/Z(\GL_n(q)).$ Our goal is to prove the following theorem.

\begin{T1}
Let  $n \ge 2$ and $(n,q)$ is neither $(2,2)$ nor $(2,3).$ If $S$ is a maximal solvable subgroup of $\GaL_n(q)$, 
 then $\Reg_S(S \cdot \SL_n(q),5)\ge 5$. In particular $b_S(S \cdot \SL_n(q)) \le 5.$
\end{T1}



Before we start the proof, let us discuss the structure of a maximal solvable subgroup $S \le \GaL_n(q)$ and fix some notation. 

Let $0< V_1 \le V$ be such that $(V_1)S=V_1$ and $V_1$ has no non-zero proper $S$-invariant subspace. It is easy to see that $S$ acts semilinearly on $V_1$. Let $$\gamma_1:S \to \GaL_{n_1}(q),$$ where $n_1=\dim V_1,$ be the homomorphism defined by $\gamma_1:g \mapsto g|_{_{V_1}}.$  Since $V_1$ is $S$-invariant, $S$ acts (semilinearly) on $V/V_1.$ Let $V_1 < V_2 \le V$ be such that $(V_2/V_1)S=V_2/V_1$ and $V_2/V_1$ has no non-zero proper $S$-invariant subspace.  Observe that $S$ acts semilinearly on $V_2/V_1$. Let $\gamma_2:S \to \GaL_{n_2}(q)$, where $n_2=\dim (V_2/V_1),$ be the homomorphism defined by $\gamma_2:g \mapsto g|_{_{(V_2/V_1)}}.$ 
Continuing this procedure we obtain the chain of subspaces 
\begin{equation}
\label{Gchain}
0=V_0<V_1< \ldots < V_k=V
\end{equation}
 and a sequence of homomorphisms $\gamma_1, \ldots, \gamma_k$ such that $V_{i}/V_{i-1}$ is $S$-invariant and has no  non-zero proper $S$-invariant  subspaces and $\gamma_i(S) \le \GaL_{n_i}(q)$ is the restriction of $S$ to $V_{i}/V_{i-1}$ for $i \in \{1, \ldots, k\}.$ 
 
 
 \begin{Lem}
 \label{GammairGL}
 If $k=1$, then either $S \cap \GL_n(q)$   is an irreducible solvable subgroup of $\GL_n(q)$ or there exists $x \in \SL_n(q)$ such that $S \cap S^x \cap \GL_n(q) \le Z(\GL_n(q)).$
 \end{Lem}
 \begin{proof}
 Since $k=1,$  $V$ has no  non-zero proper $S$-invariant subspace. Let $M$ be $S \cap \GL_n(q).$ Assume that $M$ is reducible, so there exists $0<U_1 <V$ of dimension $m$ such that $(U_1)M=U_1$ and $U_1$ is $\mathbb{F}_q[M]$-irreducible.
 
 Since $M \trianglelefteq S,$ $U_1 \varphi$ is $M$-invariant and $\mathbb{F}_q[M]$-irreducible subspace of dimension $m$ for all $\varphi \in S.$    Let $\varphi \in S$ be such that $M \varphi$ is a generator of $S/M$, so $|\varphi|=|S:M|=r.$ Let $U_i=U_1 \varphi^{i-1}$ for $i \in \{1, \ldots, r\},$ so $V=\sum_{i=1}^r U_i$ since $k=1$ and there is no non-zero proper $S$-invariant subspaces in $V$.

We claim that $$V=U_1 \oplus \ldots \oplus U_{n/m},$$
so $M$ is completely reducible. The proof of this fact is similar to a proof of Clifford's Theorem. Let $W_i=\sum_{j=1}^i U_i.$ If  $W_{i+1}=W_i$ for some $i$, then $W_{i}$ is $S$-invariant. Indeed,
$$(W_i)\varphi=(U_1 \varphi + \ldots + U_i \varphi)=(U_2 + \ldots + U_{i+1})\subseteq W_{i+1}=W_i$$
 and $(W_i)\varphi=W_i$ since $\dim (W_i)\varphi = \dim W_i.$ On the other hand, if $W_{i+1}>W_i$, then  $W_i \cap U_{i+1}$ is a proper $\mathbb{F}_q[M]$-submodule of $U_{i+1},$ so it must be zero. Hence  $W_{i+1}=W_i\oplus U_{i+1}$. Since $k=1,$ $W_{i+1}>W_i$ for $1 \le i \le n/m -1.$ So, by induction, 
 $$V=W_{n/m}=U_1 \oplus \ldots \oplus U_{n/m}.$$
 
  Let $M_i$ be the restriction of $M$ on $U_i$, so $M_i$ is an irreducible solvable subgroup of $\GL_m(q)$ for all $i \in \{1, \ldots, n/m\}.$ Notice that $q$ is not a prime since $k=1$ and $m<n$. Therefore, by Theorem \ref{irred}, one of the following holds:
  \begin{enumerate}
  \item  there exists $x_i \in \SL_m(q)$ such that $M_i \cap M_i^{x_i} \le Z(\GL_m(q))$; \label{clif1}
  \item  $m=2$ and $M_i$ lies in the normaliser $N_i$ of a Singer cycle of $\GL_2(q);$ \label{clif2}
  \item $m=2$, $q=9$ and $M_i$ lies in a subgroup $N_i$ of $\GL_2(9)$ that is conjugate in $\GL_2(9)$ to  $\GL_2(3) \cdot Z(\GL_2(9)).$ \label{clif3} 
\end{enumerate}  
In  case \ref{clif1}, define $x \in \SL_n(q)$ to be $\diag(x_1, \ldots, x_{n/m}) (I_m \otimes A(n/m))$ as in Lemma  \ref{igrekl}. It follows directely from Lemma \ref{igrekl} that $S \cap S^x \cap \GL_n(q) \le Z(\GL_n(q)).$

 In cases \ref{clif2} and \ref{clif3}, all $N_i$ are conjugate in $\GL_2(q)$ by  Lemma \ref{singconj} and computations respectively. Hence, in a suitable basis, $M$ is a subgroup of the group
$$\{\diag[g_1, \ldots, g_{n/m}] \mid g_i \in N_1\}.$$ Hence $M \le N_1 \wr P$ for any solvable transitive subgroup $P$ of $\Sym(n/m).$ By \cite[Lemma 4,\S15]{sup}, $N_1 \wr P$ is irreducible. Now there exists  $x \in \SL_n(q)$ such that $S \cap S^x \cap \GL_n(q)=M \cap M^x \le Z(\GL_n(q))$ by Theorem \ref{irred} since $n\ge 4$ and $q$ is not prime.  
 \end{proof}
 
 We start the proof of Theorem \ref{theorem} with the case $k=1.$ 
 
 
\begin{Th}
Theorem {\rm \ref{theorem}} holds for $k=1.$ 
\end{Th}
\begin{proof}
Let $M = S \cap \GL_n(q)$. By Lemma \ref{GammairGL}, if $M$ is not a subgroup of one of the groups listed in  \ref{irred11} -- \ref{irred15} of Theorem \ref{irred}, then there exists $x \in \SL_n(q)$ such that $M \cap M^x \le Z(\GL_n(q)).$ So $H \cap H^{\overline{x}}$ is a cyclic subgroup of $\overline{G}$ and by Theorem \ref{zenab} there exists $\overline{y} \in \overline{G}$ such that 
$$(H \cap H^{\overline{x}}) \cap (H \cap H^{\overline{x}})^{\overline{y}}=1.$$ Hence $b_S(G) \le 4$ and $\Reg_S(G,5) \ge 5$ by Lemma \ref{base4}.
If  $M$ is a subgroup of one of the groups in \ref{irred13} -- \ref{irred15} of Theorem \ref{irred}, then $S=M$ and $b_S(G) \le 4.$ If  $M$ is a subgroup of one of the groups in \ref{irred11} -- \ref{irred135}, then $n=2$ and $b_S(G) \le 3$ by \cite[Table 3]{burness}  for $q>4$ and by computation for $q=4$.  So $\Reg_S(G,5) \ge 5$ by Lemma \ref{base4}.
\end{proof}


 For the rest of the section we assume that $k>1.$ Our proof for $q=2$ naturally splits into two cases.  If $q=2$ and $S_0$ is the normaliser of a Singer cycle of $\GL_3(2)$, then there is no $x, y \in \GL_3(2)$ such that $S_0^{x} \cap S_0^{y}$ is contained in $RT(\GL_3(2))$ (see Theorem \ref{irred}). So, in Theorem \ref{lem41}, we assume that if $q=2$, then there is no  $i \in \{1, \ldots, k\}$ such that $\gamma_i(S)$ is the normaliser of a Singer cycle of $\GL_3(2).$ In Theorem \ref{lem42} we address the case where there exists such an $i$.


\begin{Th}
\label{lem41}
Let $k>1.$ Theorem {\rm \ref{theorem}} holds if
 $q \ge 3$ or
if  $q=2$ and there is no  $i \in \{1, \ldots, k\}$ such that $\gamma_i(S)$ is the normaliser of a Singer cycle of $\GL_3(2).$

\end{Th} 
\begin{proof}

Since $k>1,$ there exists a nontrivial $S$-invariant subspace $U<V$ of dimension $m<n$. Let $M$ be $S \cap \GL_n(q).$ We fix  
\begin{equation}
\label{basGLG}
\beta=\{v_1, v_2, \ldots, v_{n-m+1}, v_{n-m+2}, \ldots, v_n \} 
\end{equation}
such that the last $ \left( \sum_{j=1}^i n_j \right)$  
 vectors in $\beta$ form a basis of $V_i$. So  $g \in M$ has  shape 
\begin{equation}
\label{stupG}
\begin{pmatrix}
\gamma_k(g)     & * & \ldots & * & * \\
0          &    \gamma_{k-1}(g)  & \ldots & * & * \\
     &     & \ddots & \ddots  &    \\
   0    &    \ldots& 0 & \gamma_{2}(g) & * \\
      0    &   \ldots        & \ldots& 0 & \gamma_{1}(g) 
\end{pmatrix}
\end{equation}
 with respect to $\beta.$ Recall that $q=p^f$, so if $f=1,$ then $\GaL_n(q)=\GL_n(q)$ and $S=M.$ Let  $\phi \in \GaL_n(q)$ be such that \begin{equation}
\label{phidef}
\phi : \lambda v_i \mapsto \lambda^p v_i \text{ for all } i \in \{1, \ldots, n\} \text{ and } \lambda \in \mathbb{F}_q.
\end{equation}
% There are cases:\\
%{\bf A.} $n_i \ne 2$ for all $i =1, \ldots , k$ in \eqref{stup};\\
%{\bf B.} $n_i=2$ at least for two distinct $i$;\\
%{\bf C.} $n_i=2$ exactly for one $i$.

Assume that $f>1.$  Let $x= \diag[x_k, \ldots, x_1]$, where $x_i\in \SL_{n_i}(q)$ is such that $\gamma_i(M) \cap (\gamma_i(M))^{x_i} \le RT(\GL_{n_i}(q))$. Such $x_i$ exist by Lemma \ref{GammairGL} and Theorem \ref{irred}. Let $y$ be as in the proof of Lemma \ref{irrtog}. Recall that  we defined $\gamma_i$ only on $S$, but it is easy to see that $\gamma_i$ can be extended to $\Stab_{\GaL_n(q)}(V_i, V_{i-1})$ since $\gamma_i$ is the restriction on $V_i/V_{i-1}.$

 
  We present the following piece of the proof as a proposition for easy reference. 
\begin{Prop} 
\label{clm44}
Let $f>1.$ There exists $\beta$ as in \eqref{basGLG} such that for every $\varphi \in (S \cap S^x) \cap (S \cap S^x)^y$, 
$$\varphi=\phi^j g \text{ for } j \in \{0,1, \ldots,f-1\}$$
where $g=\diag[g_k, \ldots, g_1] \in \GL_n(q)$ is diagonal with $g_i=\gamma_i(g) \in D(\GL_{n_i}(q))$.
 Moreover, either $g_i \in Z(\GL_{n_i}(q))$ or $n_i=2$,  $\gamma_i(S)$ normalisers a Singer cycle of $\GL_2(q)$ and $(S \cap S^x) \cap (S \cap S^x)^y \le D(\GL_n(q)).$ %if at least one of the following holds:
%\begin{enumerate}[font=\normalfont]
% \item if  $\gamma_i(M) \cap \gamma_i(M)^{x_i} \le Z(\GL_{n_i}(q));$
% \item if $j \ne 0.$ 
%\end{enumerate} 
\end{Prop}
\begin{proof}
Consider $\varphi \in S \cap S^x.$
Notice that, by Theorem \ref{irred}, either $\gamma_i(M) \cap \gamma_i(M)^{x_i} \le Z(\GL_{n_i}(q))$ or $n_i=2$ and either $\gamma_i(M)$ lies  in the normaliser $N$ of a Singer cycle $T$ in $\GL_2(q)$ or $q=9$ and $\gamma_i(M)$ lies in $\GL_2(3) \cdot Z(\GL_2(9)).$ Let us show that we can assume that $S$ lies in the normaliser of $T$ in $\GaL_2(q)$ in the first case, and in the normaliser of $\GL_2(3)$ in $\GaL_2(9)$ in the second case. In the first case, if $\gamma_i(M)$ is imprimitive, so lies in a conjugate of $\GL_1(q)\wr \Sym(2),$ then  $\gamma_i(M) \cap \gamma_i(M)^{x_i} \le Z(\GL_{n_i}(q))$ by Theorem \ref{irred}. If $\gamma_i(M)$ is primitive, then $R=\gamma_i(M) \cap T$ is the unique maximal normal irreducible abelian subgroup, so $S$ normalises $C_{\GL_2(q)}(R)=T$. In the second case, the claim is verified using computations.    

 If $n_i=2$ and $\gamma_i(S)$ lies in the normaliser of a Singer cycle of $\GL_2(q)$ in $\GaL_2(q)$, then, by Lemma \ref{lemRTdiagfield}, we can choose $x_i \in \SL_2(q)$ such that $\gamma_i(S) \cap \gamma_i(S)^{x_i}\le RT(\GL_2(q))$.  If $\varphi$ acts linearly on $V_i/V_{i-1},$ then $\varphi$ acts linearly on $V$, so $S \cap S^x \le \GL_n(q)$.
 
  If $n_i=2,$ $q=9$ and $\gamma_i(S)$ lies in the normaliser of $\GL_2(3)$ in $\GaL_2(9),$ then
  $$\gamma_i((S \cap S^x) \cap (S \cap S^x)^y) \le Z(\GL_2(9))\rtimes \langle \phi \rangle$$  by  Lemma \ref{lemRTdiagfield}.

     Otherwise $\gamma_i(M) \cap \gamma_i(M)^{x_i} \le Z(\GL_{n_i}(q))$ and, by Lemma \ref{scfield}, there exists a basis of $V_i/V_{i-1}$ such that $\gamma_i(\varphi)= \phi^{j_i} g_i$ with scalar $g_i$ and $j_i \in \{0,1, \ldots , f-1 \}.$ Hence we can choose $\beta$ in \eqref{basGLG} such that if $\varphi \in S \cap S^x,$ then $\gamma_i(\varphi)= \phi^{j_i} g_i$ with scalar $g_i.$ Since $(\phi^{j_1})^{-1} \varphi$ acts on $V_1\ne 0$ linearly,   $(\phi^{j_1})^{-1} \varphi \in \GL_n(q),$ so $j_i=j_l$ for all $i,l \in \{1, \ldots, k\}$
and 
\begin{equation}
\label{phitri}
\varphi = \phi^j 
\begin{pmatrix}
g_k     & * & \ldots & * & * \\
0          &    g_{k-1}  & \ldots & * & * \\
     &     & \ddots & \ddots  &    \\
   0    &    \ldots& 0 & g_{2} & * \\
      0    &   \ldots        & \ldots& 0 & g_{1}
\end{pmatrix}
\end{equation}
 with $g_i \in Z(\GL_{n_i}(q))$ unless $\gamma_i(S)$ normalisers a Singer cycle of $\GL_2(q)$  where $g_i \in RT(\GL_2(q))$ and $j=0$.

%Let $i \in \{1, \ldots, k\}$ be such that $V_i=\langle v_{j_1},v_{j_1+1}, v_{j_1+2}, \ldots, v_n \rangle$. Since $\varphi \in S \cap S^x$,  $\varphi$ stabilises $V_i$, $(V_i)x$ and $\varphi =\varphi_1^x$ for some $\varphi \in S.$ So
%$(v_{j_1+1})\varphi \in V_1.$ On the other hand 
%$$(v_{j_1+1})\varphi=(v_{j_1+1})x^{-1}\varphi_1 x=(v_{j_t+1})\varphi_1 x,
%$$ so
%$$(v_{j_1+1})\varphi \in \langle v_{j_t},v_{j_t+1}, v_{j_t+2}, \ldots, v_n \rangle x =\langle v_{j_1+1},v_{j_t+1}, v_{j_t+2}, \ldots, v_n \rangle.$$ 
%So
%\begin{equation}
%\label{2diag2}
%(\varphi) \gamma_i= \phi^j 
%\begin{pmatrix}
%\lambda_1 & \lambda_2\\
%0 & \lambda_3
%\end{pmatrix}
%\end{equation}
% for some $j \in \{0, 1, \ldots, f-1\}$ and $\alpha_1, \ldots, \alpha_3 \in \mathbb{F}_q$ since $v_{j_1} \notin \langle v_{j_1+1},v_{j_t+1}, v_{j_t+2}, \ldots, v_n \rangle.$ Similar arguments show that \eqref{2diag2} holds for all $i$ such that $V_i=\langle v_{j_l},v_{j_l+1}, v_{j_l+2}, \ldots, v_n \rangle$ with $l \in \{2, \ldots, t\}.$
%Therefore, \eqref{phitri} holds with  with $g_i \in Z(\GL_{n_i}(q))$ if $n_i \ne 2$ and $g_i \in RT(\GL_2(q))$ if $n_i=2.$ 


Since $\varphi \in (S \cap S^x)^y,$
$$
\varphi = \phi^j 
\begin{pmatrix}
g_1'     & 0 & \ldots & 0 & 0 \\
*          &    g_{k-1}'  & \ldots & 0 & 0 \\
     &     & \ddots & \ddots  &    \\
   *    &    \ldots& * & g_{2}' & 0 \\
      *    &   \ldots        & \ldots& * & g_{1}'
\end{pmatrix}
$$ with $g_i$ either scalar or lower-triangular. So 
\begin{equation}
\label{phidiagG}
\varphi =\phi^j g \text{ where } g=\diag[g_k, \ldots, g_1]
\end{equation}
 with $g_i$ scalar unless $\gamma_i(S)$ normalisers a Singer cycle of $\GL_2(q)$  where $g_i \in D(\GL_2(q))$ and $j=0$.  
\end{proof}

\medskip

We now resume our proof of Theorem \ref{lem41}.
Let $\Omega$ be the set of right $S$-cosets in $G$ and let $\overline{\Omega}$ be the set of right $H$-cosets in $\overline{G}$, where the action is given by  right multiplication.  Since $S$ is maximal solvable and $$S_{G}:=\cap_{g \in G} S^g=Z(\GL_n(q)),$$ $\Reg_S(G,5)$ is the number of $\overline{G}$-regular orbits on  $\overline{\Omega}^5$.   Therefore, $$\overline{\omega}=(H\overline{g_1},H\overline{g_2},H\overline{g_3},H\overline{g_4},H\overline{g_5}) \in \overline{\Omega}^5$$ is regular under the action of $\overline{G}$ if and only if the stabiliser of $$ \omega =(Sg_1, Sg_2, Sg_3,Sg_4, Sg_5) \in \Omega^5$$ under the action of $G$ is equal to $Z(\GL_n(q))$, where $g_i$ is a preimage in $\GaL_n(q)$ of $\overline{g_i}.$  
%In addition, we can induce the action of $\PGL_n(q)$ on $\Omega$ by 
%$$(Sg) \overline{h}= Sgh,$$
%so, since $S \ge Z(\GL_n(q))$, $\omega \in \Omega^5$ is regular under the action of $\PGL_n(q)$ if and only if the stabiliser of $\omega$ under the action of $\GL_n(q)$ is $Z(\GL_n(q)).$ 
In our proof we say that   $\omega \in \Omega^5$ is regular if it is stabilised only by elements from $Z(\GL_n(q))$, so $\overline{\omega}$ is regular in terms of the standard definition that is the stabiliser of $\overline{\omega}$ in $\overline{G}$ is trivial.  

The proof of Theorem \ref{lem41} splits into five cases. To show $\Reg_S(G,5) \ge 5$, in each case we find five regular orbits in $\Omega^5$ or show that $b_S(G)\le 4,$ so $\Reg_S(G,5) \ge 5$ by Lemma \ref{base4}. Recall that $n_i = \dim V_i/V_{i-1}$ for $i \in \{1, \ldots, k\}.$  Different cases arise according to the number of $n_i$ for $i \in \{1, \ldots, k\}$ which equal 2  and in what rows of  $g \in M=S\cap \GL_n(q)$ the submatrix $\gamma_i(g)$ is located for $n_i=2.$ We now specify the cases.
\begin{description}[before={\renewcommand\makelabel[1]{\bfseries ##1}}]
\item[{\bf Case 1.}] Either $f>1$, or $f=1$ and the number of $i \in \{1, \ldots, k\}$ with $n_i=2$ is not one;
\item[{\bf Case 2.}] $f=1, $ $n$ is even, there exist exactly one $i \in \{1, \ldots, k\}$ with $n_i=2$ and $\gamma_i(g)$ appears in rows  $\{n/2,n/2+1\}$ of $g \in M$ for such $i$;
\item[{\bf Case 3.}] $f=1, $ $n$ is odd, there exist exactly one $i \in \{1, \ldots, k\}$ with $n_i=2$ and $\gamma_i(g)$ appears in rows  $\{(n+1)/2,(n+1)/2+1\}$ of $g \in M$ for such $i$;
\item[{\bf Case 4.}] $f=1, $ $n$ is odd, there exist exactly one $i \in \{1, \ldots, k\}$ with $n_i=2$ and $\gamma_i(g)$ appears in rows  $\{(n-1)/2,(n-1)/2+1\}$ of $g \in M$ for such $i$;
\item[{\bf Case 5.}] $f=1,$ there exist exactly one $i \in \{1, \ldots, k\}$ with $n_i=2$ and none of {\bf Cases 2 -- 4} holds. 
\end{description}

%In {\bf Case 1} we deal with the situation when either $f>1$, or $f=1$ and the number of $n_i$ equal to $2$ is not one. {\bf Cases 2--5} consider the situation when $f=1$ and there is exactly one $n_i$ such that $n_i=2.$

Our proof of Theorem \ref{lem41} in {\bf Cases 2 -- 5} is analogous and repetitive, so we only demonstrate a proof for {\bf Cases 1} and {\bf 5}. A detailed proof for {\bf Cases 2 -- 4} can be found in \cite{thesis}.

\medskip

Before we proceed with the proof of Theorem \ref{lem41}, let us resolve two situations which both arise often in our analysis of these cases and can be readily settled.

\begin{Prop}
\label{ni1}
Let $S$ be a maximal solvable subgroup of $\GaL_n(q),$ let $\beta$ be as in \eqref{basGLG},  $n \ge 2$ and $(n,q)$ is neither $(2,2)$ nor $(2,3).$ If  all $V_i$ in \eqref{Gchain} are such that $n_i=1,$ then there exist $x,y,z \in \SL_n(q)$ such that $S \cap S^x \cap S^y \cap S^z \le Z(\GL_n(q)).$ 
\end{Prop}
\begin{proof}
Let $\sigma =(1, n)(2, n - 1) \ldots ([n/2], [n/2 + 3/2]) \in \Sym(n).$ We define $x,y,z \in \SL_n(q)$ as follows:
\begin{itemize}
 \item $x = \diag(\sgn(\sigma), 1\ldots, 1) \cdot \per (\sigma);$ 
\item $(v_i)y=v_i$ for $i \in \{1, \ldots, n-1\}$ and $(v_n)y= \sum_{i=1}^n v_i$;
\item $(v_i)z=v_i$ for $i \in \{1, \ldots, n-1\}$ and $(v_n)z= \theta v_1 + \sum_{i=2}^n v_i$, 
\end{itemize}
where $\theta$ is a generator of $\mathbb{F}_q^*.$ Let $\varphi \in S \cap S^x \cap S^y \cap S^z.$ Since $\varphi \in S \cap S^x,$ it stabilises $\langle v_i \rangle $ for all $i \in \{1, \ldots, n\}.$ So $\varphi =\phi^j \diag(\alpha_1, \ldots, \alpha_n)$ for $j \in \{0,1, \ldots, f-1\}$ and $\alpha_i \in \mathbb{F}_q^*.$ Since $\varphi \in S^y,$ it stabilises $(V_1)y =\langle v_n \rangle y= \langle v_1 + \ldots +v_n \rangle.$ So $\alpha_1=\alpha_i$ for $i \in \{2, \ldots, n\}.$  Since $\varphi \in S^z,$ it stabilises $(V_1)z =\langle v_n \rangle z= \langle \theta v_1 + v_2 + \ldots +v_n \rangle.$ So $\theta^{p^j}=\theta$ and $j=0.$ Hence $\varphi \in Z(\GL_n(q)).$ 
\end{proof}

\begin{Prop}
\label{n3GL}
Fix a basis $\beta$ as in \eqref{basGLG}. Let $n \in \{2,3\}$ and $(n,q)$ is neither $(2,2)$ nor $(2,3)$. Let $S$ be a maximal solvable subgroup of $\GaL_n(q).$ If $S$ stabilises a  non-zero proper subspace $V_1$ of $V$, then there exist $x,y,z \in \SL_n(q)$ such that 
 $$S \cap S^x \cap S^y \cap S^z \le Z(\GL_n(q)).$$
\end{Prop}
\begin{proof}
We can assume that $V_1$ has no non-zero proper $S$-invariant subspaces. Let $V_i$ be as in \eqref{Gchain}.

If $n=2$, then $\dim V_1=1$ and $n_1=n_2=1$, so the statement follows by Proposition \ref{ni1}.


Assume $n=3$.  If $k=3,$ so $n_i=1$ for all $i$, then the statement follows by Proposition \ref{ni1}, so we assume $k=2.$

Let $n_1=1,$ so $n_2=2.$ Let $f>1.$ If $q=9$, the proposition is verified using computations, so we assume $q\ne 9.$ Assume that $\gamma_2(S)$ is a subgroup of the normaliser of a Singer cycle of $\GL_2(q)$ in $\GaL_2(q).$ Therefore, by \eqref{2sindiagodd} and \eqref{2sindiageven}, there exists $x_1 \in \SL_2(q)$ such that $\gamma_2(S) \cap (\gamma_2(S))^{x_1} \le \langle \varphi \rangle Z(\GL_2(q))$ where $\varphi = \diag(-1,1)$ if $q$ is odd and $\varphi = \begin{pmatrix}
1 & 0\\
1 & 1
\end{pmatrix}$ if $q$ is even. Hence, if $x=\diag[x_1,1],$ then $S \cap S^x \le \GL_3(q)$ and matrices in $S \cap S^x$ have shape 
$$
 \begin{pmatrix}
\alpha_1 \varphi & *\\
0 & \alpha_2
\end{pmatrix}
$$ with $\alpha_1, \alpha_2 \in \mathbb{F}_q^*.$ Let $y=\per((1,2,3))$ and 
$$
z=
\begin{cases}
\begin{pmatrix}
1&0&0\\
0&1&0\\
1&1&1
\end{pmatrix} \text{ if $q$ is odd;}\\
\begin{pmatrix}
0&1&0\\
0&0&1\\
1&0&0
\end{pmatrix} \text{ if $q$ is even.}\\
\end{cases}
$$
Calculations show that $S \cap S^x \cap S^y \cap S^z \le Z(\GL_n(q)).$

If $\gamma_2(S)$ does not normalise a Singer cycle of $\GL_2(q),$ then, by Theorem \ref{irred}, 
   there exists $x_1 \in \SL_2(q)$ such that $\gamma_2(S) \cap (\gamma_2(S))^{x_1} \cap \GL_2(q) \le Z(\GL_2(q)).$ Let $x=\diag[x_1,1]$ so $\varphi \in S \cap S^x$ has shape \eqref{phitri} with $g_2 \in Z(\GL_2(q))$ and $g_1 \in \mathbb{F}_q^*.$ In particular, let 
$$
\varphi =\phi^j
\begin{pmatrix}
\alpha_1 & 0 & \delta_1\\
 0&     \alpha_1 & \delta_2\\
0 & 0& \alpha_2  
\end{pmatrix}.
$$
 Let $y=
\begin{pmatrix}
0&1&0\\
1&0&0\\
\theta&1&1
\end{pmatrix}$ and consider $\varphi \in (S \cap S^x) \cap (S \cap S^x)^y.$ Since $\varphi \in S^y,$ it stabilises $$(V_1)y=\langle v_3 \rangle y= \langle \theta v_1 +v_2 +v_3 \rangle.$$ Therefore,
$$(\theta v_1 +v_2 +v_3)\varphi = \theta^{p^j} \alpha_1 v_1 + \alpha_1 v_2 + (\alpha_2 +\delta_1 +\delta_2)v_3 \in \langle \theta v_1 +v_2 +v_3 \rangle,$$
so $\theta^{p^j} \alpha_1= \theta \alpha_1$ and $j=0$. Thus, $(S \cap S^x) \cap (S \cap S^x)^y \le \GL_3(q).$ Now calculations  show that $(S \cap S^x) \cap (S \cap S^x)^y \le Z(\GL_3(q)).$

Now assume $f=1$ (we  continue to assume that $n_1=1$). Let $\sigma=(1,3) \in \Sym(3)$ and $x=\diag(\sgn(\sigma),1,1) \cdot {\rm perm}(\sigma) \in \SL_3(q)$. Matrices from $S \cap S^x$ have shape  
  \begin{equation*} 
  \begin{pmatrix}
*    & 0 & 0  \\
* &*    & *  \\
  0&    0    &   *      
\end{pmatrix}.
\end{equation*}
Let   
  \begin{equation*} 
y=  \begin{pmatrix}
0    & -1 & 0  \\
1 &0    & 0  \\
  1&    1    &   1      
\end{pmatrix} \in \SL_3(q).
\end{equation*}
It is easy to check that $$(S\cap S^x) \cap (S\cap S^x)^y \le Z(\GL_3(q)).$$

To complete the proof, we assume  $n_1=2,$ so $n_2=1.$ Let $\iota$ be the inverse-transpose map on $\GL_n(q)$. As mentioned in the introduction, $\iota$ induces automorphism on $\SL_3(q)$ and we can consider the group $\GaL_3(q) \rtimes \langle \iota \rangle$ (here $\iota$ and $\phi$ commute). It is easy to see that $S^{\iota}$ is a maximal solvable subgroup of $\GaL_3(q)$ stabilising $1$-dimensional subspace of $V$. Hence, by taking $x,y,z \in \SL_3(q)$ for $S^{\iota}$  as in the case $n_1=1,$ we obtain
$$(S\cap S^{x^{\iota}}) \cap (S\cap S^{x^{\iota}})^{y^{\iota}} \le Z(\GL_3(q)).$$
\end{proof} 

We now consider {\bf Cases 1} and {\bf 5} in the following two propositions. 



\begin{Prop}
\label{case1prop}
Theorem $\ref{lem41}$ holds in {\bf Case 1}.
\end{Prop} 
\begin{proof}
In this case, there are four conjugates of $M=S \cap \GL_n(q)$ whose intersection lies in $D(\GL_n(q)).$ Indeed, if $f>1,$ or $f=1$ and there is no $i \in \{1, \ldots, k\}$ (here $k$ is as in \eqref{Gchain}) such that $n_i=2$, then 
 there exist $x,y \in \SL_n(q)$ such that 
$$(M \cap M^x) \cap (M \cap M^x)^y \le D(\GL_n(q))$$
by Proposition \ref{clm44}  or Theorem \ref{irred}  and Lemma \ref{irrtog}.


Assume $f=1$ (so $S=M$) and there exists at least two $i \in \{1, \ldots, k\}$ such that $n_i=2.$ Let $t\ge 2$ be the number of such $i$-s and let $j_1, \ldots, j_t \in \{1, \ldots, n\}$ be such that    $(2 \times 2)$ blocks (corresponding to $V_i/V_{i-1}$ of dimension $2$) on the diagonal in matrices of $M$ occur in the rows $$(j_1, j_1+1), (j_2, j_2 +1), \ldots, (j_t, j_t+1).$$ 
Let $\tilde{x}=\diag(\sgn(\sigma), 1, \ldots, 1) \cdot {\rm perm}(\sigma) \in \SL_n(q),$ where  $$\sigma=(j_1, j_1+1, j_2, j_2 +1, \ldots, j_t, j_t+1)(j_1, j_1+1).$$  Let $x= \diag[x_k, \ldots, x_1]\tilde{x}$, where $x_i\in \SL_{n_i}(q)$ is such that $\gamma_i(M) \cap (\gamma_i(M))^{x_i} \le RT(\GL_{n_i}(q))$ if $n_i \ne 2$ and $x_i$ is the identity matrix if $n_i=2.$ Such $x_i$ exist by Theorem \ref{irred}. If $y$ is as in Lemma \ref{irrtog}, then calculations show that  
\begin{equation}\label{m4diag}
(M \cap M^x) \cap (M \cap M^x)^y \le D(\GL_n(q)).
\end{equation}








%Thus, by Lemma \ref{diag}*, there exists $z \in \SL_n(q)$ such that 
%$$(S,Sx,Sy,Sxy,Sz)$$ is a regular point in $\Omega^5$. 

Let $0<U <V$ be an $S$-invariant subspace and $\dim U=m.$ In general, we take $U=V_i$ for some $i<k$ (in most cases it is sufficient to take $U=V_1$). We fix $\beta$ to be as in \eqref{basGLG} such that Proposition \ref{clm44} holds for $f>1.$ So 
$$U = \langle v_{n-m+1}, v_{n-m+2}, \ldots, v_n \rangle.$$ Since in the proof we reorder $\beta,$ let us fix a second notation for  $$v_{n-m+1}, v_{n-m+2}, \ldots, v_n,$$ namely 
$w_1, \ldots, w_m$
respectively. So $U = \langle w_1, \ldots, w_m \rangle.$

  Our proof splits into the following four  subcases:
\begin{description}[before={\renewcommand\makelabel[1]{\bfseries ##1}}]
\item[{\bf Case (1.1)}] $n-m\ge 2 $ and $m \ge 2$;
\item[{\bf Case (1.2)}] $n-m=m=1$;
\item[{\bf Case (1.3)}] $n-m \ge 2$ and $m=1$;
\item[{\bf Case (1.4)}] $n-m=1$  and $m \ge 2$ (dual to {\bf Case (1.3)} via the inverse-transpose automorphism).
\end{description} 



{\bf Case (1.1).} Let $n-m\ge 2 $ and $m \ge 2$, so $n \ge 4$. We claim that there exist $z_i \in \SL_n(q)$ for $i=1, \ldots, 5$ such that the points $\omega_i=(S,Sx,Sy,Sxy,Sz_i)$ lie in distinct regular orbits. 

 Recall from \eqref{m4diag} and Proposition \ref{clm44} that, if either $f=1$ or  $f>1$ and  $\gamma_i(S)$ normalises a Singer cycle of $\GL_2(q)$ for some $i \in \{1, \ldots, k\}$, then $$(S \cap S^x) \cap (S \cap S^x)^y \le D(\GL_n(q)).$$ Also, we can assume that $n_l \ge 2$ for some $l \in \{1, \ldots, k\}$ in \eqref{Gchain}: otherwise Theorem~\ref{lem41} follows by Proposition \ref{ni1}. Let $s \in \{1, \ldots, n\}$ be such that $$V_l = \langle v_s, v_{s+1}, \ldots, v_{s+n_l-1}, V_{l-1} \rangle.$$

 If $(S \cap S^x) \cap (S \cap S^x)^y$ does not lie in $\GL_n(q)$, then   $\varphi \in (S \cap S^x) \cap (S \cap S^x)^y$ stabilises $W=\langle v_s, v_{s+1}, \ldots, v_{s+n_l-1} \rangle$ and the restriction $\varphi|_{_{W}}$ is $\gamma_l(\varphi)=\phi^j g_l$ where $g_l =\lambda I_{n_l}$ for some $\lambda \in \mathbb{F}_q^*$ by \eqref{phidiagG}. We now relabel vectors in $\beta$ as follows: $v_s, v_{s+1}, \ldots, v_{s+n_l-1}$ becomes $v_1, v_{2}, \ldots, v_{n_l}$ respectively; $v_1, v_{2}, \ldots, v_{n_l}$  becomes $v_s, v_{s+1}, \ldots, v_{s+n_l-1}$  respectively;  the remaining labels are unchanged.



  Let $z_1, \ldots, z_5 \in \SL_n(q)$ be such that $$(w_{j})z_i=u_{(i,j)}$$ for $i=1, \ldots, 5$ and $j=1, \ldots, m,$ where

\begingroup
\allowdisplaybreaks
\begin{align*}
& \left\{ 
\begin{aligned}
u_{(1,1)} & =\theta v_1+v_2;\\
u_{(1,2)} & =\theta v_1+v_3+v_4;\\
u_{(1,2+r)} & =\theta v_1+v_{4+r}\phantom{;} \text{ for } r \in \{1, \ldots, m-3\};\\
u_{(1,m)} & =\theta v_1 +\sum_{r=2+m}^{n}v_r;
\end{aligned}
\right. \\
& \left\{ 
\begin{aligned}
u_{(2,1)} & =\theta v_1+v_3;\\
u_{(2,2)} & =\theta v_1+v_2+v_4;\\
u_{(2,2+r)} & =\theta v_1+v_{4+r}\phantom{;} \text{ for } r \in \{1, \ldots, m-3\};\\
u_{(2,m)} & =\theta v_1 +\sum_{r=2+m}^{n}v_r;\\
\end{aligned}
\right. \\
& \left\{ 
\begin{aligned}
u_{(3,1)} & =\theta v_1+v_4;\\
u_{(3,2)} & =\theta v_1+v_2+v_3;\\
u_{(3,2+r)} & =\theta v_1+v_{4+r}\phantom{;} \text{ for } r \in \{1, \ldots, m-3\}; \\
u_{(3,m)} & =\theta v_1 +\sum_{r=2+m}^{n}v_r;
\end{aligned}
\right. \stepcounter{equation}\tag{\theequation} \label{orb}  \\
& \left\{ 
\begin{aligned}
u_{(4,1)} & =v_2+v_3;\\
u_{(4,2)} & =\theta v_1+v_2+v_4;\\
u_{(4,2+r)} & =v_2+v_{4+r}\phantom{;} \text{ for } r \in \{1, \ldots, m-3\};\\
u_{(4,m)} & =v_2 +\sum_{r=2+m}^{n}v_r;
\end{aligned}
\right. \\
& \left\{ 
\begin{aligned}
u_{(5,1)} & =v_2+v_4;\\
u_{(5,2)} & =\theta v_1+v_2+v_3;\\
u_{(5,2+r)} & =v_2+v_{4+r}\phantom{;} \text{ for } r \in \{1, \ldots, m-3\};\\
u_{(5,m)} & =v_2 +\sum_{r=2+m}^{n}v_r.
\end{aligned}
\right.  
\end{align*}
\endgroup

Recall that $\theta$ is a generator of $\mathbb{F}_q^*.$ Such $z_i$ always exist since $m<n$ and $u_{(i,1)}, \ldots, u_{(i,m)}$ are linearly independent  for every $i =1 ,\ldots, 5$. 
Consider $\varphi \in (S \cap S^x) \cap (S \cap S^x)^y \cap S^{z_1}$, so $\varphi=\phi^jg$ as in \eqref{phidiagG} with $g=\diag({\alpha_1, \ldots, \alpha_n})$ for $\alpha_i \in \mathbb{F}_q^*.$ Notice that $\varphi$ stabilises the subspace $Uz_1.$ 
Therefore, 
$$(u_{(1,1)})\varphi =
\begin{cases}
 \theta^{p^j} \alpha_1 v_1 + \alpha_2 v_2\\
 \mu_1 u_{(1,1)} + \ldots + \mu_m u_{(1,m)}
\end{cases}
$$ for $\mu_i \in \mathbb{F}_q^*.$ The first line does not contain $v_i$ for $i>2$, so $\mu_i=0$ for $i>1$ and $(u_{(1,1)})\varphi=\mu_1 u_{(1,1)}= \mu_1(\theta v_1 +v_2).$ Therefore, $\mu_1= \alpha_2.$  Assume  that $\varphi \notin \GL_n(q)$, so, by the arguments before \eqref{orb}, $\alpha_1 =\alpha_2$,  $\theta^{p^j-1}=1$ and $j=0$ which is a contradiction. Thus, we can assume $\varphi=g \in \GL_n(q).$
 The same arguments as in the proof of Lemma \ref{diag} show that $$(u_{(1,j)} )g = \delta_j u_{(1,j)} \text{ for all } j\in \{1,\ldots, m\}$$ for some $\delta_j \in \mathbb{F}_q$. Since all $u_{(1,j)}$ have $v_1$ in the decomposition \eqref{orb}, all $\delta_j$ are equal and $g$ is a scalar matrix, so $\omega_1$ is a regular point in $\Omega^5.$ It is routine to check that every $\omega_i \in \Omega^5$ is regular. 

Assume that $\omega_1$ and $\omega_2$ lie in the same orbit in $\Omega^5$, so there exists $\varphi \in \GaL_n(q)$ such that $\omega_1 \varphi =\omega_2$. This implies $$\varphi \in (S \cap S^x) \cap (S \cap S^x)^y \cap z_1^{-1}Sz_2.$$
In particular, $\varphi= \phi^j g$ with $g =\diag(\alpha_1, \ldots, \alpha_n) \in D(\GL_n(q))$ and $g=z_1^{-1}\psi z_2$, $\psi \in S$. Consider $(u_{(1,1)})\varphi$.  Firstly, 
\begin{equation}\label{orr}
(u_{(1,1)})\varphi=(\theta v_1+v_2)\varphi= \theta^{p^j} \alpha_1 v_1 + \alpha_2 v_2;
\end{equation}
on the other hand,
$$(u_{(1,1)})\varphi =(u_{(1,1)})z_1^{-1}\psi z_2=(w_{1})\psi z_2=(\eta_{1} w_{1} + \ldots +\eta_{m} w_{m})z_2 \in Uz_2,$$
for some $\eta_i \in \mathbb{F}_q$. So $(u_{(1,1)})\varphi= \delta u_{(2,2)}$ for some $\delta \in \mathbb{F}_q$ since there are no $v_i$ for $i>2$ in decomposition \eqref{orr}. However, 
$$\theta^{p^j} \alpha_1 v_1 + \alpha_2 v_2=\delta u_{(2,2)}$$
if and only if $(u_{(1,1)})\varphi=0$, which contradicts the fact that $\varphi$ is invertible. Hence  there is no such $\varphi$, so points $\omega_1$ and $\omega_2$ are in  distinct regular orbits on $\Omega^5.$ The same arguments show that all $\omega_i$ lie in  distinct regular orbits on $\Omega^5.$   

\medskip

{\bf Case (1.2).} If $n-m=m=1$, then $n=2$ and  Theorem \ref{lem41} follows by Proposition \ref{n3GL}.
% consists of upper triangular matrices and $S \cap S^x \le D(\GL_2(q))$, where $x=\diag(\sgn(\sigma),1) \cdot {\rm per}(\sigma) \in \SL_2(q)$ and  $\sigma=(1,2)$. Thus, $$b_S(G)\le 3$$ by Lemma \ref{diag}, so $\Reg_{S}(G,5) \ge 5$ by Lemma \ref{base4}.  

\medskip

{\bf Case (1.3).} Let $m=1$ and $n-m \ge 2$, so $n \ge 3.$ %so $k=2$ in \eqref{stup} and $S_1=Im(\gamma_1)$ is a maximal irreducible solvable subgroup of $GL(n-1,q).$
If $n=3$, then Theorem \ref{lem41} follows by Proposition \ref{n3GL}, so we assume that $m=1$ and $n \ge 4.$ If $k\ge 4$ in \eqref{stup}, then $U = \langle v_{n_1+n_2+1}, \ldots, v_n \rangle$ is an $S$-invariant subspace with $m=n_3 + \ldots + n_k \ge 2$ and $n-m=n_1 +n_2 \ge 2$, so the proof as in {\bf Case (1.1)} using \eqref{orb} works.

 Assume that $k=3.$ If $n_3 \ge 2$, then    $U = \langle v_{n_3+1}, \ldots, v_n \rangle$ is an $S$-invariant subspace with $m=n_1 + n_2 \ge 2$ and $n-m=n_3 \ge 2$, so the proof as in {\bf Case (1.1)} using \eqref{orb} also works.

Let $n_3=1.$ Since $m=1$, $n_1=1.$ If $n_2=2$ and $q\in \{2,3,5\}$, then Theorem~\ref{lem41} is verified by computation. Otherwise, by Theorem \ref{irred} and Lemma \ref{scfield}, there exist $\tilde{x}_1,\tilde{x}_2 \in \SL_{n_2} (q)$  such that $$\gamma_2(S) \cap (\gamma_2(S))^{\tilde{x}_1} \cap (\gamma_2(S))^{\tilde{x}_2} \le \langle \phi^i \rangle Z(\GL_{n_2}(q))$$ for some $i \in \{0, 1, \ldots, f-1\}$. Let $x_1, x_2 \in \GL_n(q)$ be the matrices $\diag[1,\tilde{x}_1,1]$ and $\diag[1,\tilde{x}_2,1]$ respectively.  Let $\tilde{x}= \diag(\sgn(\sigma),1 \ldots, 1) \cdot \per(\sigma) \in \SL_n(q)$ with $\sigma=(1,n)$. Calculations show that $(S \cap S^{x_1} \cap S^{x_2 \tilde{x}}) \cap \GL_n(q) \le D(\GL_n(q)).$ So, with respect to a basis $\beta$ as in \eqref{basGLG},  $\varphi \in S \cap S^{x_1} \cap S^{x_2 \tilde{x}}$ has shape
$$\phi^j \diag(\alpha_1, \alpha_2, \ldots, \alpha_2, \alpha_3)$$ with $j \in \{0,1, \ldots, f-1\}$ and $\alpha_1, \alpha_2, \alpha_3 \in \mathbb{F}_q^*.$ Let 
$$y=
\begin{pmatrix}
1& 0      & 0\ldots  & 0 &0 \\
0& 1      & 0 \ldots  & 0 &0\\
&       &\ddots &  &  &\\
&       & & \ddots &  &\\
0& 0      & 0& \ldots & 1 &0 \\
1& \theta & 1& \ldots & 1 &1
\end{pmatrix}.$$
Calculations show that $S \cap S^{x_1} \cap S^{x_2 \tilde{x}} \cap S^y \le Z(\GL_n(q)).$ 

Now let $k=2$, so $n_2=n-m=n-1$ and $n_2>2$. If $(n,q)$ is $(4,2)$ or $(5,3)$, then Theorem \ref{lem41} is verified by  computation.  Otherwise,  by Theorem \ref{irred} and Lemma \ref{scfield},
there exists $x \in \SL_n(q)$ such that $S \cap S^x$ consists of elements of  shape (with respect to a basis $\beta$ as in \eqref{basGLG})
\begin{equation*} 
\phi^j
  \begin{pmatrix}
\alpha_1 &0 & \dots    & 0 & \delta_1  \\
0 & \alpha_1 & \dots  & 0 & \delta_2 \\
 &  & \ddots  &  &  \\
 0& \dots & 0  & \alpha_{1} & \delta_{n-1} \\
0 & 0 & \dots  & 0 & \alpha_2 \\
\end{pmatrix}
\end{equation*}
where $\alpha_i, \delta_i \in \mathbb{F}_q$ and $j \in \{0,1, \ldots, f-1\}.$
Let $\sigma =(1, n)(2, n - 1) \ldots ([n/2], [n/2 + 3/2]) \in \Sym(n).$ Let $$y = \per(\sigma) \cdot \diag(\sgn(\sigma), 1, \ldots, 1, \theta, \theta^{-1}).$$ Calculations show that $(S \cap S^x) \cap (S \cap S^x)^y \le Z(\GL_n(q)).$
%Let $z \in \SL_n(q)$ be such that $(v_i)z=v_{i+1}$ for $1\le i <n-1$, $(v_{n-1})z=v_1$ and $(v_n)z= \sgn(\sigma)( v_1 + \ldots + v_n),$ where $\sigma:=(1,n-1)(2,n-2) \ldots ([n/2],[n/2]+1).$   
%Let $g=z^{-1}hz \in(S\cap S^x) \cap (S\cap S^x)^z$ with $h \in (S\cap S^x).$ For $i \ne n$
%$$(v_i)g = \alpha_i v_{i} + \beta_i v_{n}$$
%since $g \in (S\cap S^x)$ and 
%$$(v_i)g= (v_i)z^{-1}hz =(v_{i-1})hz= (\alpha_{i-1}' v_{i-1} + \beta_{i-1}' v_{n})z=\alpha_{i-1}'v_{i} + \beta_{i-1}'(v_1 +\ldots + v_n).$$ Here  $i-1$ is read modulo $(n-1),$ so if $i=1$, then $i-1=n-1.$ Therefore, $\beta_i=0$ and $g$ is diagonal. Since $g$ stabilises 
%$\langle(v_n)z \rangle= \langle v_1+ \ldots +v_n \rangle$, $g$ is scalar, so
%$$(S\cap S^x) \cap (S\cap S^x)^z \le Z(\GL_n(q)).$$
Thus, $b_S(G)\le 4$ and $\Reg_{S}(G,5) \ge 5.$
 

\medskip

{\bf Case (1.4).} Now let $n-m=1$  and $m \ge 2$,  so $n\ge 3.$   Let $\iota$ be the inverse-transpose map on $\GL_n(q)$ and  consider the group $\GaL_3(q) \rtimes \langle \iota \rangle$. It is easy to see that {\bf Case (1.3)} holds for $S^{\iota}$. Hence, by taking $x,y,z \in \SL_n(q)$ for $S^{\iota}$  as in {\bf Case (1.3)}, we obtain
$$(S\cap S^{x^{\iota}}) \cap (S\cap S^{x^{\iota}})^{y^{\iota}} \le Z(\GL_3(q)).$$ Thus, $b_S(G)\le 4$ and $\Reg_{S}(G,5) \ge 5.$ This concludes the proof of Proposition \ref{case1prop}.
\end{proof}

\bigskip

Recall that $f=1$ for {\bf Cases 2 -- 5}, so $\GaL_n(q)=\GL_n(q)$ and $S=M.$ Therefore, $\beta$ and $\gamma_i$ are as in Lemma \ref{supreduce} and matrices from $S$ have shape \eqref{stup}. Denote $\gamma_i(S)$ by $S_i$.






\begin{Prop}
\label{case5prop}
Theorem $\ref{lem41}$ holds in {\bf Cases 2} -- {\bf 5}.
\end{Prop}
\begin{proof}
As mentioned before Proposition \ref{ni1}, we only provide a detailed proof for {\bf Case 5}.

 Let $i \in \{1, \ldots, k\}$  be  such that $n_i=2$. Our proof  splits into three subcases:
\begin{description}[before={\renewcommand\makelabel[1]{\bfseries ##1}}]
\item[{\bf Case (5.1)}] $n\ge 9$ and $\gamma_i(g)$ appears in rows  $\{l,l+1\}$ of $g \in S$ where $l>n/2$ if $n$ is even and $l>(n+1)/2$ if $n$ is odd;
\item[{\bf Case (5.2)}] $n \ge 9$ and $\gamma_i(g)$ appears in rows $\{l-1,l\},$ where $l \le n/2$ if $n$ is even and $l \le (n-1)/2$ if $n$ is odd;
\item[{\bf Case (5.3)}] $n<9$.
\end{description} 

{\bf Case (5.1).} Consider the case when  the only $(2\times 2)$ block is in  rows $\{l,l+1\},$ where $l>n/2$ if $n$ is even and $l>(n+1)/2$ if $n$ is odd. Denote $s:=n-l$. Let $x= \diag[x_k, \ldots, x_1]$, where $x_i\in \SL_{n_i}(q)$ is such that $S_i \cap S_i^{x_i} \le RT(\GL_{n_i}(q))$ if $n_i \ne 2$ and $x_i$ is the identity matrix if $n_i=2.$ If $y$ is as in Lemma \ref{irrtog}, then  calculations show that  
$$(S \cap S^x) \cap (S \cap S^x)^y$$ consists of matrices of shape 
\begin{equation}\label{ugol} 
\diag[\alpha_1, \ldots, \alpha_{s-1}, \begin{pmatrix} \alpha_s & \beta_s \\ 0 & \alpha_{s+1} \end{pmatrix}, \alpha_{s+1}, \ldots, \alpha_{l-1},\begin{pmatrix} \alpha_l & 0 \\ \beta_{l+1} & \alpha_{l+1} \end{pmatrix}, \alpha_{l+2}, \ldots, \alpha_n ]
\end{equation}

Assume that $n \ge 9$. We take as the $S$-invariant subspace $U$ the subspace with  basis $\{v_{l}, \ldots, v_n\}$, so
$m\ge 2$ and $n-m\ge [n/2] \ge 4.$  Let us rename some basis vectors for convenience, so denote vectors 
$$v_1, \ldots, v_{s-1}, v_{s+2}, \ldots, v_{l-1}, v_{l+2}, \ldots, v_{n}$$ by $$w_1, w_2, \ldots, w_{n-4}$$ respectively.
Let $z_1, \ldots, z_5 \in \SL_n(q)$ be such that $$(v_{n-m+j})z_i=u_{(i,j)}$$ for $i=1, \ldots, 5$ and $j=1, \ldots, m,$ where
\begin{gather}\label{orb2t1}
\left\{ 
\begin{aligned}
u_{(1,1)} & = v_{s} + v_{l+1};\\
u_{(1,2)} & = v_{s}+ v_{s+1}+v_{l}+w_1;\\
u_{(1,2+r)} & =v_{s+1}+w_{1+r}\phantom{;} &\text{ for } r \in \{1, \ldots, m-3\};\\
u_{(1,m)} & =v_{s+1}+\sum_{r=m-1}^{n-4}w_r;
\end{aligned}
\right.
\end{gather}
\begin{gather}%\label{orb2t12}
\left\{ 
\begin{aligned}
u_{(2,1)} & = v_{s} + v_{l+1};\\
u_{(2,2)} & = v_{s}+v_{s+1}+v_{l}+w_2;\\
u_{(2,2+r)} & =v_{s+1}+w_{(2,r)}\phantom{;} &\text{ for } r \in \{1, \ldots, m-3\};\\
u_{(2,m)} & =v_{s+1}+\sum_{r=m-2}^{n-5}w_{(2,r)}.
\end{aligned}
\right. 
\end{gather}
Here $w_{(2,1)}, \ldots, w_{(2,n-5)}$ are equal to $w_1, w_3, \ldots, w_{n-4}$ respectively.  
\begin{gather}
\left\{ 
\begin{aligned}
u_{(3,1)} & = v_{s} + v_{l+1};\\
u_{(3,2)} & = v_{s}+ v_{s+1}+v_{l}+w_3;\\
u_{(3,2+r)} & =v_{s+1}+w_{(3,r)}\phantom{;} &\text{ for } r \in \{1, \ldots, m-3\};\\
u_{(3,m)} & =v_{s+1}+\sum_{r=m-2}^{n-5}w_{(3,r)}.
\end{aligned}
\right.
\end{gather} 
Here $w_{(3,1)}, \ldots, w_{(3,n-5)}$ are equal to $w_1, w_2, w_4 \ldots, w_{n-4}$ respectively.  
\begin{gather} 
\left\{ 
\begin{aligned}%\label{orb2t2}
u_{(4,1)} & = v_{s} + v_{l+1};\\
u_{(4,2)} & = v_{s}+ v_{s+1}+v_{l}+w_4;\\
u_{(4,2+r)} & =v_{s+1}+w_{(4,r)}\phantom{;} &\text{ for } r \in \{1, \ldots, m-3\};\\
u_{(4,m)} & =v_{s+1}+\sum_{r=m-2}^{n-5}w_{(4,r)}.
\end{aligned}
\right. 
\end{gather}
Here $w_{(4,1)}, \ldots, w_{(4,n-5)}$ are equal to $w_1, w_2, w_3, w_5 \ldots, w_{n-4}$ respectively.  
\begin{gather} 
\left\{ 
\begin{aligned}\label{orb2t3}
u_{(5,1)} & = v_{s} + v_{l+1};\\
u_{(5,2)} & = v_{s}+ v_{s+1}+v_{l}+w_5;\\
u_{(5,2+r)} & =v_{s+1}+w_{(5,r)}\phantom{;} &\text{ for } r \in \{1, \ldots, m-3\};\\
u_{(5,m)} & =v_{s+1}+\sum_{r=m-2}^{n-5}w_{(5,r)}.
\end{aligned}
\right. 
\end{gather}
Here $w_{(5,1)}, \ldots, w_{(5,n-5)}$ are equal to $w_1, w_2, w_3, w_4, w_6, \ldots, w_{n-4}$ respectively.

 Notice that $n \ge 9$ and $n-m \ge 3$ is necessary for such a definition of $u_{(i,j)}$. 

Let $$\omega_i=(S,Sx,Sy,Sxy,Sz_i).$$ We first show that the $\omega_i$ are regular points of $\Omega^5$. Consider $\omega_1$.  The regularity of the remaining points can be shown using the same arguments. Let $t \in (S \cap S^x) \cap (S \cap S^x)^y \cap S^{z_1}$, so it takes shape \eqref{ugol} for some $\alpha_i, \beta_i \in \mathbb{F}_q$ and stabilises the subspace $Uz_1.$ Thus
\begin{equation}\label{b2t1}
(u_{(1,1)})t=\alpha_{s}v_s +\beta_s v_{s+1} + \alpha_{l+1} v_{l+1} +\beta_{l+1}v_l,
\end{equation}
since $t \in (S \cap S^x) \cap (S \cap S^x)^y$ and 
\begin{equation}\label{b2t2}
(u_{(1,1)})t=\eta_1 u_{(1,1)} +\eta_2 u_{(1,2)} + \ldots + \eta_m u_{(1,m)},
\end{equation}
since $t$ stabilises  $Uz_1.$ There is no $w_i$ for $i \ge 1$ in decomposition \eqref{b2t1}, so in \eqref{b2t2} we obtain
$$(u_{(1,1)})t=\eta_1 u_{(1,1)}=\alpha_s u_{(1,1)},$$
in particular $\beta_{s}=\beta_{l+1}=0$, so $t$ is diagonal. Therefore,
\begin{equation}\label{2b2t2}
(u_{(1,2)})t=\alpha_{s}v_s +\alpha_{s+1}v_{s+1}+ \alpha_{l} v_{l} +\alpha_{1}w_1,
\end{equation}
must also lie in $\langle u_{(1,1)}, \ldots, u_{(1,m)} \rangle$. There are no $v_{l+1}$ and $w_i$ for $i \ge 2$ in \eqref{2b2t2}, so 
$$(u_{(1,2)})t= \alpha_s u_{(1,1)}.$$ If $i>2$ then
 the same arguments show that 
$$(u_{(1,i)})t=\alpha_{s+1} u_{(1,i)}=\alpha_s u_{(1,i)},$$
for the remaining $i \in \{1, \ldots , m\},$ so $t$ is scalar. Therefore, $\omega_1$ is regular.

Now we claim that  the $\omega_i$ lie in distinct orbits of $\Omega^5.$ Here we prove that $\omega_1$ and $\omega_2$  lie in distinct orbits;  the remaining cases are similar.  Assume that $\omega_1 g =\omega_2$ for $g \in \GL_n(q),$ so 
$$g \in (S \cap S^x) \cap (S \cap S^x)^y \cap z_1^{-1}Sz_2.$$
Therefore, $g$ has shape \eqref{ugol} and $g=z_1^{-1}hz_2$ for $h \in S,$ so 
\begin{equation}\label{b2t3}
(u_{(1,2)})g=\alpha_{s}v_s +\beta_s v_{s+1} + \alpha_{s+1} v_{s+1} +\alpha_{l}v_l+\alpha_1 w_1,
\end{equation}
and 
\begin{equation*}%\label{b2t4}
(u_{(1,2)})g=(u_{(1,2)})z_1^{-1}hz_2=(v_{n-m+2})hz_2 = \delta_1 u_{(2,1)} + \ldots + \delta_m u_{(2,m)}.
\end{equation*}
Since there is no $w_2$ in \eqref{b2t3}, $\delta_2$ is zero. Therefore, there must be no $v_l$ in \eqref{b2t3}, so $\alpha_l$ must be zero, which contradicts the existence of such an invertible $g.$ Thus, $\omega_1$ and $\omega_2$ lie in  distinct orbits. 

\medskip

{\bf Case (5.2).} Consider the case when $n \ge 9$ and the only $(2\times 2)$ block is in  rows $\{l-1,l\},$ where $l \le n/2$ if $n$ is even and $l \le (n-1)/2$ if $n$ is odd. We take as the $S$-invariant subspace $U$ the subspace with  basis $\{v_{l+1}, \ldots, v_n\}$, so
$m\ge [n/2]$ and $n-m \ge 2.$  If $n-m \ge 3$, then the proof is  as in {\bf Case (5.1).} 

 Consider the case  $n-m=2,$ so $n_k=2$. Let $x= \diag[x_k, \ldots, x_1]$, where $x_i\in \SL_{n_i}(q)$ is such that $S_i \cap S_i^{x_i} \le RT(\GL_{n_i}(q))$ if $n_i \ne 2$ and $x_i$ is the identity matrix if $n_i=2.$ Let $y$ be as in Lemma \ref{irrtog}. If $k=2$ in \eqref{stup}, then there exists $x_1 \in \SL_{n_1}(q)$ such that $S_1 \cap S_1^{x_1} \le D(\GL_{n_1}(q))$, by %Lemma \ref{sch}, Lemma \ref{prtoirr} and Section \ref{sec5}
Theorem \ref{irred}, since $n_1=n-2\ge 7.$ Thus, $(S\cap S^x) \cap (S\cap S^x)^y \le D(\GL_n(q))$ and the rest of the proof is as in {\bf Case 1}.

Let $k>2$ and $n_1 \ne 1$. So $n_1 >2$, since the only $(2\times 2)$ block is in  rows $\{l-1,l\},$ where $l \le n/2$. % there is only one $i$ such that $n_i=2,$ 
We can take $U$ to be the subspace with  basis $\{v_{n -n_1+1}, \ldots, v_n\}$. Therefore, $n-m\ge 3$  and the proof is the same as the one using \eqref{orb2t1} -- \eqref{orb2t3}. 

Let $k>2$ and $n_1=1.$ Denote $\diag(-1,1\ldots,1) \cdot \per((2,n))$ by $\tilde{x}$. Calculations show that  $S \cap S^{x \tilde{x}} \le RT(\GL_n(q))$, so $(S \cap S^{x \tilde{x}}) \cap (S \cap S^{x \tilde{x}})^y\le D(\GL_n(q))$. The rest of the proof is  as in {\bf Case 1}, since we can take $U=V_{k-1}=\langle v_3, \ldots, v_n \rangle,$ so $n-m=2,$ $m\ge 2.$

\medskip

{\bf Case (5.3).}  Now let $n<9.$ If $n=3,$ then Theorem~\ref{lem41} follows by Proposition \ref{n3GL}, so $n \in \{4,5,6,7,8\}$.

Assume that $n_i=1$ for some $i\le k$. So there is a $(1\times 1)$ block in the row $j \le n$ of $g \in S$. Let $\tilde{x}$ be $\diag(\sgn(\sigma), 1, \ldots, 1) \cdot \per(\sigma)$ with $\sigma=(s+1,j)$ if $j>s+1$ and  $\sigma=(j,s,s+1)$ if $j<s.$ It is easy to see that 
$$S \cap S^{x \tilde x}\le RT(\GL_n(q)),$$
so the rest of proof is  as in {\bf Case 1}.

If we exclude the cases previously resolved, then $k \in \{2,3\}$ and it remains to handle the following list of possibilities:
\begin{enumerate}
\item[$(a)$] $k=2$ and $(n,n_1,n_2)$ is one of the following:
$$(5,2,3), (5,3,2), (6,2,4), (6,4,2), (7,2,5), (7,5,2), (8,2,6), (8,6,2)$$
\item[$(b)$] $k=3$ and $(n, n_1, n_2, n_3)$ is $(8,2,3,3)$ or $(8,3,3,2).$
\end{enumerate}
If $(n,q) \ne (6,3),$ then, by Theorem \ref{irred}, for all $i$ such that $n_i\ne 2$  there exist $x_i \in \SL_{n_i}(q)$ such that $S_i \cap S_i^{x_i} \le D(\GL_{n_i}(q))$. Let $x=\diag[x_k, \ldots, x_1]$, where $x_i$ is the identity matrix  if $n_i=2$. It is easy to check directly that $$(S \cap S^x) \cap (S \cap S^x)^y \le D(\GL_n(q))$$
for all cases above, so the rest of the proof is  as in {\bf Case 1}.

  For $(n,q)=(6,3)$ the result is established by computation. This concludes the proof of Proposition \ref{case5prop}.
\end{proof}
  
 Theorem \ref{lem41}  now follows from     Propositions \ref{case1prop} and \ref{case5prop}.
\end{proof}


\begin{Th}
\label{lem42}
If $k>1$,  $q=2$ and  there exists  $i \in \{1, \ldots, k\}$ such that $\gamma_i(S)$ is the normaliser of a Singer cycle of $\GL_3(2),$ then Theorem {\rm \ref{theorem}} holds.
\end{Th} 
\begin{proof}
Notice that $\GL_n(2)=\SL_n(2),$ so $S \cdot \SL_n(2)=\GL_n(2).$  %If $S$ is irreducible, then $b_S(\GL_n(q)) \le 3$ by Lemma \ref{irred}, so $\Reg_S(G,5) \ge 5$ by Lemma \ref{base4}.


Since $k>1,$ $S$ stabilises a nontrivial invariant subspace $U$ of dimension $m<n$. Assume that matrices in $S$ take shape \eqref{stup} in the  basis $$\{v_1, v_2, \ldots, v_{n-m+1}, v_{n-m+2}, \ldots, v_n \}.$$

The main difference from the case $q>2$ is  that if $S_0$ is the normaliser of a Singer cycle of $\GL_3(2)$, then there is no $x, y \in \GL_3(2)$ such that $S_0^{x} \cap S_0^{y}$ is contained in $RT(\GL_3(2))$ (see Theorem \ref{irred}). Since all Singer cycles are conjugate in $\GL_3(2),$ we assume that $S_0$ is the normaliser of $$
\left\langle
\begin{pmatrix}
0&0&1\\
1&0&0\\
0&1&1
\end{pmatrix}
 \right\rangle.
$$ By Lemma \ref{lemRTdiagfield}, if $$x_0=\begin{pmatrix}
0&1&0\\
1&0&1\\
0&0&1
\end{pmatrix},
$$ then 
\begin{equation*}
S_0 \cap S_0^{x_0} = P =  \left\{
\begin{pmatrix}
1 & 1& 1\\
0& 1 & 0\\
1 & 0 & 0
\end{pmatrix},
\begin{pmatrix}
0 & 0& 1\\
0& 1 & 0\\
1 & 1 & 1
\end{pmatrix},
\begin{pmatrix}
1 & 0& 0\\
0& 1 & 0\\
0 & 0 & 1
\end{pmatrix} \right\}. 
\end{equation*} 


By Theorem \ref{irred}, for $n_i \ge 3$   there exist $x_i\in \GL_{n_i}(2)$ such that $S_i \cap S_i^{x_i} \le Z(\GL_{n_i}(2))$ if $S_i$ is not conjugate to $S_0$ and $S_i \cap S_i^{x_0} \le P$ if $S_i$ is conjugate to $S_0$. Notice that $Z(\GL_{n_i}(2))=1.$  Let $x=\diag[x_k, \ldots, x_1]$, where $x_i$ is an identity matrix if $n_i = 2.$  Therefore, matrices in $S \cap S^x$ are upper triangular except for $(2 \times 2)$ and  $(3 \times 3)$ blocks on the diagonal. Let $y$ be the permutation matrix corresponding to the permutation $$(1,n)(2, n-1) \ldots (n/2, n/2+1)$$ if $n$ is even, and $$(1,n)(2, n-1) \ldots ((n-1)/2, (n+1)/2)$$ if $n$ is odd.

 Fix some $(3\times 3)$ block on the diagonal of matrices in $S$ such that $S_i=S_0$ is in this block, so $P$ is in this block in $S \cap S^x$.
 This block intersects  one, two or three blocks in $(S \cap S^x)^y.$  
If it intersects  at least two blocks, then the  matrices in $(S \cap S^x) \cap (S \cap S^x)^y$ have the following restriction  to the chosen $(3 \times 3)$ block:
$$
\begin{pmatrix}
*&*&0\\
*&*&*\\
*&*&*
\end{pmatrix}. 
$$ 
Since  the only such matrix in $P$ is the identity, every matrix in $(S \cap S^x) \cap (S \cap S^x)^y$ has the identity submatrix in this block.  

If the chosen $(3 \times 3)$ block intersects  a bigger block in $(S \cap S^x)^y$, then it must lie in the block with  $S_i \cap S_i^{x_i}$ for some $i$ such that $n_i>3$ and  all such matrices are scalar.

Let the chosen $(3 \times 3)$ block intersect another $(3 \times 3)$ block in $(S \cap S^x)^y$ which consists of matrices in $P$. By Theorem \ref{irred}, $b_{S_0}(\GL_3(2))=3,$ so there exists $y_0$ such that $$S_0 \cap S_0^{x_0} \cap S_0^{y_0} =1.$$ Let $\tilde{y}=\diag[1, \ldots, 1, y_0, 1, \ldots, 1]$ where $y_0$ is in the chosen block. Therefore,
 $$(S \cap S^x) \cap (S \cap S^x)^{y\tilde{y}}$$
consists of matrices which are diagonal except, possibly, for $(2 \times 2)$ blocks.

The rest of the proof is  as in Theorem \ref{lem41}. 
\end{proof}

Theorem \ref{theorem} now follows by Theorems \ref{lem41} and \ref{lem42}.

\section{Solvable subgroups not contained in $\GaL_n(q)$}
\label{sec412}



Recall that $V=\mathbb{F}_q^n$ and let $\beta=\{v_1, \ldots, v_n\}$ be a basis of $V$. Let $\Gamma= \GaL _n(q)= \GL_n(q) \rtimes \langle  \phi_{\beta} \rangle$
and $A=A(V)=A(n,q):=\Gamma \rtimes \langle \iota_{\beta} \rangle$ where $\iota_{\beta}$ is the inverse-transpose  map of $\GL_n(q)$ with respect to $\beta.$  Our goal is to prove the following theorem.

\begin{T4}
Let $n\ge 3.$ If $S$ is a maximal solvable subgroup of $A(n,q)$ not contained in $\Gamma,$ then one of the following holds:
\begin{enumerate}
\item[$(1)$] $b_S(S \cdot \SL_n(q))\le 4$;
\item[$(2)$] $n=4,$ $q \in \{2,3\}$, $S$ is the normaliser in $A(n,q)$ of the stabiliser in $\GaL_n(q)$ of a $2$-dimensional subspace of $V$, $b_S(S \cdot \SL_4(q))=5$ and $\Reg_S(S \cdot \SL_4(q),5)\ge 5.$
\end{enumerate}
\end{T4} 

Before we start the proof, let us discuss the structure of a maximal solvable subgroup
$S$ of $A(n,q)$ and fix some notation. In this section, we assume  $S$ is not contained in $\Gamma.$

Consider the action of $\Gamma$ on the set $\Omega_1$ of subspaces of $V$ of dimension $m<n.$ This action is transitive and equivalent to the action of $\Gamma$ on the set
$$\Omega_1'=\{\mathrm{Stab}_{\Gamma}(U) \mid U<V, \dim U= m\}$$ by conjugation. Let $\Omega_2$ be the set of subspaces of $V$ of dimension $n-m$ and let  $$\Omega_2'=\{\mathrm{Stab}_{\Gamma}(W) \mid W<V, \dim W= n-m\}.$$ It is easy to see that $\Gamma$ acts on $\Omega=\Omega_1 \cup \Omega_2$ with orbits $\Omega_1$ and $\Omega_2$ (respectively, on $\Omega'=\Omega_1' \cup \Omega_2'$ with orbits $\Omega_1'$ and $\Omega_2'$). Notice that if $m=n/2$, then $\Omega_1=\Omega_2$ and there is only one orbit. We can extend this action to an action of $A$ on $\Omega'$ by conjugation which is equivalent to the following action of $A$ on $\Omega$:  for $U,W \in \Omega \text{ and } \varphi \in A$,
$$ U\varphi =W \text{ if and only if } 
(\mathrm{Stab}_{\Gamma}(U))^{\varphi}=\mathrm{Stab}_{\Gamma}(W).$$ 
In particular, if $U= \langle v_{n-m+1}, v_{n-m+2}, \ldots, v_n \rangle$, then $U\iota_{\beta}=\langle v_1, \ldots, v_{n-m}\rangle:=U'$ of dimension $n-m.$ Moreover, if $\varphi=\iota_{\beta}g$ with $g \in \Gamma$ and $h \in \Gamma,$ then 
$$(Uh)\varphi=U\iota_{\beta}h^{\iota_{\beta}}g=U'h^{\iota_{\beta}}g.$$  So elements of $A\backslash \Gamma$ interchange $\Omega_1$ and $\Omega_2.$

Let us now define the action of $A$ on the pairs of subspaces $(U,W)$ of $V$ with $\dim U=m \le n/2$ and $\dim W=n-m$ where either
\begin{itemize}
\item  $U<W$, or
\item $U\cap W= \{0\}.$
\end{itemize}
 Here we let $(U,W) \varphi =(U\varphi, W \varphi).$ The pair $(U, W)$ is not ordered, but for convenience we usually list first the subspace of smaller dimension.  Notice that this action is equivalent to the action of $A$ by conjugation on 
 \begin{itemize}
\item {$\{\mathrm{Stab}_{\Gamma}((U,W)) \mid U \le W < V, \dim U=m, \dim W=n-m\}$};
\item $\{\mathrm{Stab}_{\Gamma}((U,W)) \mid U , W < V, U \cap W = \{0\}, \dim U=m, \dim W=n-m\}$ 
 \end{itemize}
respectively
where $\mathrm{Stab}_{\Gamma}((U,W))=\mathrm{Stab}_{\Gamma}(U) \cap \mathrm{Stab}_{\Gamma}(W).$ 



\begin{Def}
\label{deftypepm}
 Let $M$ be the stabiliser in $A$ of a pair $(U,W)$ of subspaces of $V$ where $\dim U=m \le n/2$, $\dim W=n-m$. %Assume that $M$ is a maximal subgroup of $A.$
\begin{itemize}
\item If $U \le W$, then $M$ is a subgroup {\bf of type} $P_{m,n-m};$
\item If $U\cap W=\{0\}$, then $M$ is a subgroup {\bf of type} $\GL_m(q) \oplus \GL_{n-m}(q).$
\end{itemize}
We denote the class of subgroups of $A$ as above by $\mathcal{C}_1^*.$
\end{Def}
We follows \cite[\S 4.1]{kleidlieb} for the definition of ``type''. Let $S$  be a maximal solvable subgroup of $A$ not contained in $\Gamma$ such that $S \le M$ where $M$ is a maximal subgroup of $A$ contained in Aschbacher's class $\mathcal{C}_1.$ By \cite[\S 4.1]{kleidlieb}, $M$ is as in Definition \ref{deftypepm}.  Notice that  $\mathcal{C}_1$ does not contain the subgroups of type $\GL_{n/2} \oplus \GL_{n/2}$ when $n$ is even as such a subgroup lies in the maximal subgroup stabilising the decomposition $V=U \oplus W$ which is in the Aschbacher's class $\mathcal{C}_2.$ We define $\mathcal{C}_1^*$ to include such subgroups as it suits our purposes better.
In \cite[\S 4.1]{kleidlieb}, the type $P_{m,n-m}$ is used only when $m<n/2;$ when $m=n/2$ such $M$ are labelled $P_m$. We let $m \le n/2$ and use the label $P_{m,n-m}$ since it allows us to use more uniform statements.

%Let us discuss some structure of $S$. Let $m_1$ be the minimal such that there exist a maximal subgroup $M_1 \le A$ of type $P_{m_1, n-m_1}$ or $\GL_{m_1}(q) \oplus \GL_{n-m_1}(q).$ 
Let  $M$ be as in Definition \ref{deftypepm} and let $\beta=\{v_1, \ldots, v_n\}$ be a basis of $V$ such that $U=\langle v_{n-m+1}, \ldots, v_n \rangle$ and 
$$W= \begin{cases}
\langle v_{m+1}, \ldots, v_n\rangle &\text{ if } M \text{ is of type } P_{m, n-m};\\
\langle v_1, \ldots, v_{n-m}\rangle &\text{ if } M \text{ is of type } \GL_{m}(q) \oplus \GL_{n-m}(q).
\end{cases} 
$$ 
For such $\beta$ we say that it is {\bf associated} with $(U,W).$
 By Definition \ref{deftypepm}, $M$ is the normaliser of $\mathrm{Stab}_{\Gamma}(U,W)$ in $A$. Therefore, with respect to $\beta$,
\begin{equation}
\label{Mstr}
M =\begin{cases}
\mathrm{Stab}_{\Gamma}(U,W) \rtimes \langle \iota_{\beta}a(n,m) \rangle & \text{ if } M \text{ is of type } P_{m, n-m};\\
\mathrm{Stab}_{\Gamma}(U,W) \rtimes \langle \iota_{\beta} \rangle & \text{ if } M \text{ is of type } \GL_{m}(q) \oplus \GL_{n-m}(q)
\end{cases}
\end{equation}
where $a(n,m) \in \GL_n(q)$ is 
\begin{equation}
\label{adef}
\begin{pmatrix}
0 & 0 & I_{m} \\
0 & I_{n-2m} & 0 \\
I_{m} & 0 & 0
\end{pmatrix}.
\end{equation}

\begin{Lem}
\label{GRmmim2}
Let $n\ge 3$ and let $S \le A$ be a maximal solvable subgroup that is not contained in $\Gamma$. If $m$ is the least integer such that $S$ lies in $M$ as in \eqref{Mstr}, then 
 $\tilde{S}=S \cap \Gamma$ acts  on $U$ irreducibly.
\end{Lem} 
\begin{proof}
Let $\beta$ and $M$ be as in \eqref{Mstr} and let $P$ be $\mathrm{Stab}_{\Gamma}(U,W),$ so 
$$M=P \rtimes \langle \iota_{\beta} a\rangle$$
where $a$ is $a(n,m)$ if $M$ is of type $P_{m,n-m}$ and $a=I_n$ otherwise.


Assume that $\tilde{S}$ acts  reducibly on $U$, so  $\tilde{S}$ stabilises $U_1 <U$ of dimension $s<m.$  Assume that $U_1$ is a minimal such subspace, so $\tilde{S}$ stabilises no proper non-zero subspace of $U_1.$ Let $\varphi \in S \backslash \tilde{S}$ be such that
$S=\langle \tilde{S}, \varphi \rangle$, so $\varphi=\iota_{\beta}a \cdot g$ with $g \in P.$  Hence $\tilde{S}$ stabilises $W_1=U_1 \varphi$ of dimension  $n-s$. Notice that $U_1<W_1$ if $M$ is of type $P_{m,n-m}$ and $W_1 \cap U_1=0$ otherwise.
Since $\varphi^2 \in \Gamma,$ we obtain $\varphi^2 \in \tilde{S}$ and   $$W_1\varphi=U_1\varphi^2=U_1.$$
Therefore, $S$ normalises the stabiliser of $(U_1, W_1)$ in $\Gamma,$ so $S$ lies in a  subgroup of $A$ of type $P_{s, n-s}$ or $\GL_{s}(q) \oplus \GL_{n-s}(q)$  which contradicts the assumption of the lemma. 
\end{proof}

Notice that if we take $U=W=V$, then the proof of Lemma \ref{GRmmim2} implies the following statement.
\begin{Lem}
\label{GRirred2}
If $S \le A$ is not contained  in a  subgroup of $A$ from the class $\mathcal{C}_1^*$, then $\tilde{S}=S \cap \Gamma$ acts irreducibly on $V$.
\end{Lem}

\begin{Th}
Theorem {\rm \ref{theoremGR}} holds if $S$ is not contained  in a subgroup of $A$ from the class $\mathcal{C}_1^*$.
\end{Th}
\begin{proof}
Let $\hat{S}:=S \cap \GL_n(q).$ Lemmas \ref{GammairGL}  and \ref{GRirred2} and Theorem \ref{irred} imply that either there exists $x \in \SL_n(q)$ such that $\hat{S} \cap \hat{S}^x \le Z(\GL_n(q))$, or $\hat{S}$ lies in one of the groups in \ref{irred14}, \ref{irred15} of Theorem \ref{irred}. In the latter case the statement is verified by computation. 

Otherwise, $\overline{S} \cap \overline{S}^{\overline{x}}$ is an abelian subgroup of $G=(S \cdot \SL_n(q))/Z(\GL_n(q))$ where $\overline{\phantom{G}}:S \cdot \SL_n(q) \to G$ is the natural homomorphism.  By Theorem \ref{zenab}, there exist $\overline{y} \in G$ such that $\overline{S} \cap \overline{S}^{\overline{x}} \cap (\overline{S} \cap \overline{S}^{\overline{x}})^{\overline{y}}=1.$ So $b_S(S \cdot \SL_n(q)) \le 4$ and the statement follows. 
\end{proof}

For the rest of the section we assume that $S$ lies in a subgroup of $A$ from the class $\mathcal{C}_1^*$.

 Let $m_1$ be minimal such that $S$ lies in a  subgroup $M_1<A$ of type $P_{m_1, n-m_1}$ or $\GL_{m_1}(q) \oplus \GL_{n-m_1}(q)$ stabilising subspaces $(U_1,W_1)$. Let $V_1$ be $V$;  let $V_2=W_1/U_1$ and $n_2=\dim V_2$. Let $\beta$ be a basis associated with $(U_1,W_1).$ Notice that elements from $\mathrm{Stab}_{\Gamma}(U_1,W_1)$ induce semilinear transformations on $V_2$, and $\iota_{\beta}$ induces the inverse-transpose map on $\GL(V_2).$  Hence there exists a homomorphism 
$$\psi: S \to A(V_2) $$ mapping $x \in S$ to the element it induces on $V_2.$ Denote $\psi(S)$ by $S|_{_{V_2}}.$ Let $m_2$ be minimal such that $S|_{_{V_2}}$ lies in a  subgroup $M_2<A(V_2)$ of type $P_{m_2, n_2-m_2}$, or $\GL_{m_2}(q) \oplus \GL_{n_2-m_2}(q)$ stabilising subspaces $(U_2/U_1,W_2/U_1)$ of $V_2$.

 Repeating the arguments above we obtain a chain of subspaces 
\begin{equation}
\label{GRUs}
 0=U_0<U_1< \ldots < U_k<V \text { with } \dim U_i/U_{i-1} =m_i \text{ for } i \in\{1 \ldots, k\},
\end{equation} 
 subspaces $W_i$  for $i \in \{0, 1, \ldots, k\}$ where $W_0=V$, and groups $S|_{V_i}$ where $V_i=W_{i-1}/U_{i-1}$ for $i \in \{1, \ldots, k+1\}.$ Here $S|_{V_{i}}$ stabilises $(U_{i}/U_{i-1}, W_{i}/U_{i-1})$ for $i\in \{1, \ldots, k\}$ and $S|_{V_{k+1}}$ stabilises no subspace of $V_{k+1}.$ If $U_k=W_k,$ then $V_{k+1}$ is a zero space and  $S|_{V_{k+1}}$ is trivial.  Let $\beta=\{v_1, \ldots, v_n\}$ be a basis of $V$ such that  $\beta_i=\{v_1 +U_{i-1}, \ldots, v_n +U_{i-1}\} \cap V_i$ is a basis associated with $(U_{i}/U_{i-1}, W_{i}/U_{i-1})$ for $i\in \{1, \ldots, k\}.$

Let $\varphi \in S$. Hence $\varphi|_{_{V_i}}$ lies in $M_i$ for $i=1, \ldots, k$, so $$\varphi|_{_{V_i}}= (\iota_{\beta_i} a_i)^l \cdot g_i \text{ with } l\in \{0,1\} \text{ and } g_i \in \mathrm{Stab}_{\GaL (V_i)}((U_{i}/U_{i-1}, W_{i}/U_{i-1})),$$
where $a_i$ is $a(n_i,m_i)$  as in \eqref{adef} if $M_i$ is of type $P_{m_i, n_i-m_i}$ and $a_i$ is $I_{n_i}$ otherwise. Let $a \in \GL_n(q)$ be such that  $a|_{_{V_k}}=a_k$ and 
\begin{equation}
\label{adefGR}
a|_{_{V_i}}= \begin{cases}
\begin{pmatrix}
a|_{_{V_{i+1}}} & 0\\
0 & I_{m_{i}}
\end{pmatrix} &\text{ if } M_i \text{ is of type } \GL_{m_i}(q) \oplus \GL_{n_i-m_i}(q)\\
\begin{pmatrix}
0 & 0 & I_{m_i}\\
0 & a|_{_{V_{i+1}}} & 0\\
I_{m_i} & 0 & 0
\end{pmatrix} &\text{ if } M_i \text{ is of type } P_{m_i, n_i -m_i}
\end{cases}
\end{equation}
 for $i \in \{1, \ldots, k-1\}.$  Therefore, 
\begin{equation}
\label{GRvarphishape}
\varphi = (\iota_{\beta}a)^l \cdot (\phi_{\beta})^j \cdot g,
\end{equation}
where $g \in \GL_n(q)$ and $g|_{_{V_i}}$ stabilises $(U_{i}/U_{i-1}, W_{i}/U_{i-1}).$ More specifically, let $s \le k$ be the number of $i \in \{1, \ldots, k\}$ such that $M_i$ is of type $P_{m_i, n_i-m_i}$ and let $i_1 < \ldots < i_s$ be the corresponding $i$-s. Therefore,
\begin{equation}
\label{GRguppdiag}
g= \begin{pmatrix}
g_{i_1} '&* &\ldots &\ldots &\ldots & &* \\
 & \ddots &* &\ldots &\ldots& \ldots &* \\
 & &g_{i_s}' &* &\ldots & &* \\
 & & &g_{k+1} &* & &* \\
 & & & & g_k &* & *\\
 & & & & & \ddots &*\\
& & & & & & g_1
\end{pmatrix}
\end{equation} 
where $g_i, g_i' \in \GL_{m_i}(q)$ and $g_{k+1} \in \GL_{n_{k+1}}(q).$


\begin{example}
Let $k=3,$ and let $M_1$ and $M_3$ be of types $P_{m_i, n_i-m_i}$ for $i=1$ and $i=3$ respectively. Let $M_2$ be of type $\GL_{m_2}(q) \oplus \GL_{n_2-m_2}(q).$ Then 
$$
a=\begin{pmatrix}
 & & & & &  I_{m_1} \\
 &{\cellcolor{gray!50}} &{\cellcolor{gray!50}} &{\cellcolor{gray!50}} I_{m_3} &{\cellcolor{gray!20}} & \\
 & {\cellcolor{gray!50}} &{\cellcolor{gray!50}} I_{n_3-2m_3} &{\cellcolor{gray!50}} &{\cellcolor{gray!20}} & \\
 &{\cellcolor{gray!50}} I_{m_3} &{\cellcolor{gray!50}} & {\cellcolor{gray!50}} &{\cellcolor{gray!20}} & \\
 & {\cellcolor{gray!20}} & {\cellcolor{gray!20}} &{\cellcolor{gray!20}} &{\cellcolor{gray!20}} I_{m_2} & \\
I_{m_1} & & & & & 
\end{pmatrix}
$$
and $$\varphi = (\iota_{\beta}a)^l \cdot (\phi_{\beta})^j \cdot \begin{pmatrix}
 g_1 '&* &* &* &* &* \\
 & g_3' &* &* &0 &* \\
 & &g_4 &* &0 &* \\
 & & &g_3 &0 &* \\
 & & & & g_2 &* \\
 & & & & & g_1
\end{pmatrix}, $$
where $g_1, g_1' \in \GL_{m_1}(q),$ $g_2 \in \GL_{m_2}(q),$ $g_3, g_3' \in \GL_{m_3}(q),$ $g_4 \in \GL_{n_3-2m_3}(q).$
\end{example}

\begin{Lem}
\label{GRdiag}
Let $n \ge 3$. Let  $S$ be a maximal solvable subgroup of $A$. Assume that $S$ is contained in a subgroup of $A$ of type $P_{m,n-m}$ or $\GL_m(q) \oplus \GL_{n-m}(q)$ for some $m \le n/2.$ Let $\beta$ be as described after \eqref{GRUs}.
\begin{enumerate}
\item[$(1)$] If none of $(m_i, q)$ and $(n_i,q)$ lies in $\{(2,2), (2,3), (2,5)\},$ then there exist $x, y \in \SL_n(q)$ such that if $\varphi \in (S \cap S^x \cap S^y) \cap \Gamma,$ then 
\begin{equation}
\label{GRvarphdiagc1}
\varphi=(\phi_{\beta})^j \cdot \diag[g_{i_1}, \ldots,g_{i_s}, g_{k+1}, g_k, \ldots,  g_1]=(\phi_{\beta})^j \cdot \diag(\alpha_1, \ldots, \alpha_n)
\end{equation}
 where $\alpha_i \in \mathbb{F}_q^*$ and $j \in \{0,1, \ldots, f-1\}.$ Moreover, $g_i, g_i' \in Z(\GL_{m_i}(q))$ for $i \in \{1, \ldots, k\}$ and $g_{k+1} \in Z(\GL_{n_{k+1}}(q))$.
\item[$(2)$] If $q \in \{2,3,5\}$ and at least one of $m_i$ or $n_i$ is $2$, then there exist $x,y \in \SL_n(q)$ such that if $g \in S \cap S^x \cap S^y \cap \GL_n(q)$, then $$g=\diag[g_{i_1}, \ldots,g_{i_s}, g_{k+1}, g_k, \ldots,  g_1]$$ with $g_i, g_i' \in \GL_{m_i}(q)$ for $i \in \{1, \ldots, k\},$ $g_{k+1} \in \GL_{n_{k+1}}(q)$, and one of the following holds:
\begin{enumerate} 
\item[$(2a)$]  for each $i \in \{1, \ldots, k+1\}$, either  $g_i$, $g_i'$ are  scalar matrices over $\mathbb{F}_q$ or $m_i=2$ and  $g_i$, $g_i'$ are upper-triangular matrices in $\GL_2(q);$
\item[$(2b)$] there exist exactly one $j \in \{1, \ldots, k+1\} \backslash \{i_1, \ldots, i_s\}$ such that $g_j \in \GL_2(q),$ and $g_i$, $g_i'$ are scalar for $i \in \{1, \ldots, t\} \backslash \{j\}.$
\end{enumerate}
\end{enumerate}
\end{Lem}
\begin{proof}
We start with the proof of $(1),$ so neither $(m_1,q)$ nor $(n_2,q)$ lies in $\{(2,2), (2,3),(2,5)\}.$ 
Let us fix $q$ and assume that $n$ is minimal such that there exists $S\le A$ for which the statement of the lemma does not hold. Let $K$ be $S|_{_{U_1}} \le A(U_1)$ and let $R$ be $S|_{_{V_2}} \le A(V_2).$

We claim that there exist $x_2, y_2 \in \SL_{n_2}(q)$ such that $(1)$ holds for $R \le A(V_2)$. Indeed, if $k>1$, then $R$ stabilises $(U_2,W_2)$ and $R$ is not a counterexample to the lemma since $n_2<n.$ If $k=1$, then $\tilde{R}=R \cap \Gamma$ acts irreducibly on $V_2$ by Lemma \ref{GRirred2}.  Hence, by  Lemma \ref{GammairGL} and Theorem \ref{irred},   there exist $x_2,y_2 \in \SL(V_2)$ such that $\tilde{R} \cap \tilde{R}^{x_2} \cap \tilde{R}^{y_2} \cap \GL(V_2)= Z(\GL(V_2)).$ Now the claim follows by Lemma \ref{scfield}.

By the same argument, $\tilde{K}=K \cap \GaL_{m_1}(q)$ acts irreducibly on $U_1$, and, by Lemma \ref{GammairGL},   $K \cap \GL_{m_1}(q)$ lies in a subgroup $T$ of $\GL_{m_1}(q)$ such that  either $T$ is an  irreducible maximal solvable subgroup  or there exists $x_1 \in \SL_{m_1}(q)$ such that $T \cap T^{x_1} \le Z(\GL_{m_1}(q))$.

\medskip

 {\bf(1)} Our proof of $(1)$ splits into two cases: when $U_1 \le W_1$ and  $U_1\cap W_1=\{0\}$ respectively. We provide the proof for the first case, the proof for the second case is analogous. 
\medskip

 Let $U_1 \le W_1,$ so $n_2=n-2m_1$ and $S$ lies in a  subgroup $M_1$ of $A$ of type $P_{m_1, n-m_1}.$ Recall that, by \eqref{Mstr}, elements in $M_1$ have shape 
\begin{equation}
\label{GRPmnmdiag}
(\iota_{\beta} a)^l \cdot (\phi_{\beta})^j \cdot 
\begin{pmatrix}
h_1' & * & *\\
 0   & h_2 & *\\
 0   & 0 & h_1
\end{pmatrix}
\end{equation}
with $l \in \{0,1\}$, $j \in \{0, 1, \ldots, f-1\}$, $h_2 \in \GL_{n_2}(q)$ and $h_1, h_1' \in \GL_{m_1}(q).$ Let $K'$ be the restriction of $S$ on $V/W_1$ and let $\hat{K}$ and $\hat{K'}$ be $K \cap \GL_{m_1}(q)$ and $K' \cap \GL_{m_1}(q)$ respectively.

 Recall that $S=\langle \tilde{S}, \varphi \rangle,$ where $\varphi = \iota_{\beta} a \cdot g$ with $g \in \tilde{M_1}=M_1 \cap \GaL_n(q),$ so 
$$g= (\phi_{\beta})^{j_g} \cdot 
\begin{pmatrix}
g_3 & * & * \\
0  & g_2 & *\\
0 & 0 & g_1 \\
\end{pmatrix}
$$ with $g_1, g_3 \in \GL_{m_1}(q),$ $g_2 \in \GL_{n_2}(q)$ and $j_g \in \{0,1, \ldots, f-1\}.$ 

Since $(S \cap \GL_n(q))^{\varphi}=S \cap \GL_n(q),$  $$\hat{K}'=(\hat{K}^{\top})^{\phi^{j_g} g_3}.$$ Recall that $\hat{K}$ lies in a subgroup $T$ of $\GL_{m_1}(q)$ as described in the third paragraph of the proof.  Let $T'=(T^{\top})^{\phi^{j_g} g_3},$ so $\hat{K'}\le T'.$ 

By Lemma \ref{GammairGL} and Theorem \ref{irred}, there exist $x_1, y_1 \in \SL_{m_1}(q)$ such that $$T \cap T^{x_1} \cap T^{y_1}\le Z(\GL_{m_1}(q)).$$ Let $x_1'= ((x_1^{-1})^{\top})^{\phi^{j_g} g_3}$ and let $y_1'= ((y_1^{-1})^{\top})^{\phi^{j_g} g_3}$,  so $$T' \cap T'^{x_1'} \cap T'^{x_1'} \le Z(\GL_{m_1}(q)).$$  If $T$ is not as $S$ in  \ref{irred11} -- \ref{irred15} of Theorem \ref{irred}, then we assume $y_1=y_1'=I_{m_1}.$ So, by Lemma \ref{scfield}, we can assume that 
\begin{align*}
K \cap K^{x_1} \cap K^{y_1} \cap \GaL_{m_1}(q) \le \langle \phi_{\beta_1} \rangle Z(\GL_{m_1}(q));\\
K' \cap (K')^{x_1} \cap (K')^{y_1} \cap \GaL_{m_1}(q) \le \langle \phi_{\beta_3} \rangle Z(\GL_{m_1}(q)),
\end{align*}
 where $\beta_1$ is the basis of $U_1$ consisting of the last $m_1$ vectors of $\beta$ and $\beta_{3}=\{v_1 +W_1, \ldots, v_{m_1} + W_1\}$ is a basis of $V/W_1.$


If $T$ is as $S$ in  \ref{irred11} -- \ref{irred15} of Theorem \ref{irred}, then we claim that $T$ and $T'$ are conjugate by an element of $\GL_{m_1}(q).$ Indeed, if $T$ is as $S$ in \ref{irred11}, \ref{irred12}, \ref{irred14}  then both $T$ and $T'$ are normalisers of Singer cycles which are conjugate by Lemma \ref{singconj}. If $T$ is as $S$ in \ref{irred135} or \ref{irred13},  then $T$ and $T'$ are conjugate by \cite[\S 21, Theorem 6 $(1)$]{sup}.  If $T$ is as $S$ in  $(5),$ then $T=(T^{\top})^{\phi^{j_g}},$ so $T'=T^{g_3}$ with $g_3 \in \GL_{m_1}(q).$ Notice that $\Det(T)=\mathbb{F}_q^*$ if $T$ is not as $S$ in \ref{irred135}:  it is easy to check directly. If $T$ is as $S$ in \ref{irred135}, then $\Det(T)$ consists of all squares in $\mathbb{F}_q^*,$ in particular $\pm 1 \in \Det(T).$ Hence, for $\alpha = \pm 1$ and a given  $t_1 \in \GL_{m_1}(q)$ such that  $T'=T^{t_1}$ there exists $t_2 \in \GL_{m_1}(q)$ such that  $(T')^{t_2}=T$ and $\det(t_1)\det(t_2)=\alpha$: we can take $t_2=\alpha t_1^{-1}.$ 



Let $x= \diag[x_1', x_2, x_1]$, let 
$$y=
\begin{pmatrix}
0 & 0 & t_2 y_1 \\
0 & y_2 & 0 \\
t_1 y_1'& 0 & 0
\end{pmatrix}$$
where $t_1, t_2 \in \GL_{m_1}(q)$ are such that $\det(y)=1$, and 
$$T^{t_1}=T' \text{ and } (T')^{t_2}=T \text{ for } T \text{ as } S \text{ in \ref{irred11} -- \ref{irred15} of Theorem \ref{irred}.}$$
 So $x,y \in \SL_n(q).$ Such $t_1,t_2$ exist by the previous paragraph and since $\det(y_1)= \det(y_2)=1.$

Consider $h \in S \cap S^x \cap S^y \cap \Gamma.$ Using \eqref{GRPmnmdiag}, we obtain that elements in $S^y$ have shape 
\begin{equation}
\label{GRSy}
(\iota_{\beta} a)^l \cdot (\phi_{\beta})^j \cdot 
\begin{pmatrix}
h_1' & 0 & 0\\
 *   & h_2 & 0\\
 *   & * & h_1
\end{pmatrix}
\end{equation}
with $l \in \{0,1\}$, $j \in \{0, 1, \ldots, f-1\}$, $h_2 \in \GL_{n_2}(q)$ and $h_1, h_1' \in \GL_{m_1}(q).$ Therefore,  
by \eqref{GRPmnmdiag},
$$h= (\phi_{\beta})^j \cdot \diag[h_1',h_2,h_1],$$ so
\begin{align*}
h|_{_{V_2}} & =(\phi_{\beta_2})^j \cdot h_2 \le (R \cap R^{x_2} \cap R^{y_2}) \cap \GaL_{n_2}(q);\\  
h|_{_{U_1}} & =(\phi_{\beta_1})^j \cdot h_1 \le (K \cap K^{x_1} \cap K^{y_1}) \cap \GaL_{m_1}(q);\\
h|_{_{V/W_1}} & =(\phi_{\beta_3})^j \cdot h_1' \le (K' \cap (K')^{x_1'} \cap (K')^{y_1'}) \cap \GaL_{m_1}(q)
\end{align*} 
where $\beta_2=\{v_{m_1+1} +U_1, \ldots, v_{n-m_1}+U_1\}$ is a basis of $V_2.$ Therefore, $h_1$ and $h_1'$ are scalar matrices, and  $h_2$ is as $\varphi$ in \eqref{GRvarphdiagc1}, so $h$ has the shape claimed by the lemma.
 This concludes the proof of part $(1)$ of Lemma \ref{GRdiag}.

\bigskip

{\bf (2)} Now we start the proof of part $(2)$ of Lemma \ref{GRdiag}. So $q \in \{2,3,5\}$ and at least one of $m_i$ and $n_i$ is $2$ for $i \in \{1, \ldots, k\}$. Notice that $f=1$, so $\phi_{\beta}$ is trivial and $\Gamma =\GL_n(q)$. Therefore, matrices in $S \cap \Gamma$ are block-upper-triangular and have shape \eqref{GRguppdiag} where 
\begin{align*}
g_i \in S_i & :=S|_{_{U_i/U_{i-1}}} \cap \GL_{m_i} (q)  &&\text{ for } i\in \{1, \ldots, k\};\\
g_{k+1} \in S_k & := S|_{_{V_{n_{k+1}}}} \cap \GL_{n_{k+1}} (q); &&\\
g_i \in S_i' & :=S|_{_{V_i/W_{i}}} \cap \GL_{m_i} (q) && \text{ for } i\in \{i_1, \ldots, i_s\}.
 \end{align*}
Denote $\GL_{m_i}(q)$ by $G_i$ for $i \in \{1, \ldots, k\}$ and $\GL_{n_{k+1}}(q)$ by $G_{k+1}.$ By Lemma \ref{GRirred2}, $S_i$ and $S_i'$ are irreducible solvable subgroup of $G_i$ for $i \in \{1, \ldots, k+1\}$. By Theorem \ref{irred}, for $i \in \{1, \ldots, k+1\}$ one of the following holds:
\begin{enumerate}[label=(\roman*)]
\item $G_i=\GL_2(q);$ \label{RGlistA} 
\item there exists $x_i \in G_i$ of determinant $1$ such that $S_i \cap S_i^{x_i} \le Z(G_i);$ \label{RGlistB}
\item $(G_i, S_i) \in \{(\GL_4(3),\GL_2(3) \wr \Sym(2)), (\GL_3(2), N_{\GL_3(2)}(T))\}$ where $T$ is a Singer cycle of $\GL_3(2)$ and there exists $x_i, z_i \in G_i$ such that $S_i \cap S_i^{x_i} \cap S_i^{z_i} \le Z(G_i)$.\label{RGlistC}
\end{enumerate}
The same statement is true for $S_i'$ and we denote corresponding conjugating elements by $x_i'$ and $z_i'$. If $G_i=\GL_2(q)$, then we take $x_i=x_i'=I_2.$

 If condition \ref{RGlistC}  holds for $S_i,$ then, by Theorem \ref{irred}, as discussed in the proof of part $(1)$, $S_i'$ is conjugate to $S_i$ in $G_i$. % and $\Det(S_i)= \mathbb{F}_q$, so, for a given $\lambda \in \mathbb{F}_q^*$, there exists  $t \in G_i$ such that $\det(t)=1$ and $S_i'=S_i^t.$ 

Let us define $y_i, y_i' \in G_i$ for $i \in \{1, \ldots, k+1\}.$ % Assume first that $G_i=\GL_4(3)$ and $S_i=\GL_2(3) \wr \Sym(2).$ 
 If $M_i$ is of type $\GL_{m_i}(q) \oplus \GL_{n_i-m_i}(q)$ or $i=k+1$ then let $y_i \in G_i$ be such that $S_i \cap S_i^{x_i} \cap S_i^{y_i} \le Z(G_i).$ If $M_i$ is of type $P_{m_i, n_i-m_i},$ then let $y_i, y_i' \in G_i$ be such that 
\begin{align*}
S_i \cap S_i^{x_i} \cap (S_i')^{y_i'} & \le Z(G_i)\\
S_i' \cap (S_i')^{x_i'} \cap (S_i)^{y_i} & \le Z(G_i).
\end{align*} 

If conditions \ref{RGlistA} or \ref{RGlistB} hold for $S_i$ (respectively $S_i'$), then we take $y_i$ and $y_i'$ to be the identity matrix in $G_i.$ Let 
$$y=\diag[y_{i_1}', \ldots, y_{i_s}', y_{k+1}, \ldots, y_k] \cdot a.$$
If $\det (y) \ne 1$, then $q=3$ and $\det(y)=-1.$ If so, then we pick  $i \in \{1, \ldots, k+1\}$  such that $G_i=\GL_2(3)$ and change $y_i$ to be $\diag(-1,1)$ which is equivalent to multiplying a line of $y$ by $-1.$  Therefore, we can assume that $y \in \SL_n(q).$   

Let $r$ be the number of $G_i$ and $G_i'$ equal to $\GL_2(q)$. Our proof of $(2)$ splits into two cases: 
when $r \ge 2$
and $r=1$ respectively.


\medskip

{\bf Case (2.1).}  
  Let    $(2 \times 2)$ blocks (corresponding to the $S_i$ and the $S_i'$ lying in $\GL_2(q)$) on the diagonal in matrices of $\tilde{S}=S \cap \Gamma$ occur in the rows $$(j_1, j_1+1), (j_2, j_2 +1), \ldots, (j_r, j_r+1).$$ 
Let $\tilde{x}=\diag(\sgn(\sigma), 1, \ldots, 1) \cdot {\rm perm}(\sigma) \in \SL_n(q),$ where  $$\sigma=(j_1, j_1+1, j_2, j_2 +1, \ldots, j_r, j_r+1)(j_1, j_1+1).$$  
Let $x= \diag[x_{i_1}',\ldots, x_{i_s}', x_{k+1}  \ldots, x_1]\tilde{x}$. Notice that $\det x=1.$ Calculations show that if $h \in \tilde{S} \cap \tilde{S}^x \cap \tilde{S}^y$, then 
$$h=\diag[h'_{i_1}, \ldots, h'_{i_s}, h_{k+1}, \ldots, h_k]$$
where
\begin{itemize}
\item $h_i \in S_i \cap S_i^{x_i} \le Z(G_i)$ (respectively $h_i' \in S_i' \cap (S_i')^{x_i} \le Z(G_i)$) if condition \ref{RGlistB} holds for $S_i$ (respectively $S_i'$);
\item $h_i \in S_i \cap S_i^{x_i} \cap (S_i')^{y_i'} \le Z(G_i)$ (respectively $h_i' \in S_i' \cap (S_i')^{x'_i} \cap (S_i)^{y_i} \le Z(G_i)$) if condition \ref{RGlistC} holds for $S_i$ (respectively $S_i'$);
\item $h_i \in \GL_2(q)$ (respectively $h_i' \in \GL_2(q)$) is upper-triangular if condition \ref{RGlistA} holds
$S_i$ (respectively $S_i'$).
\end{itemize}
 Therefore, $h_i$ and $h_i'$ are either scalar or upper-triangular, so  $(2a)$   of Lemma \ref{GRdiag} holds.

\medskip

{\bf Case (2.2).} Assume that the number $r$ of $G_i$ and $G_i'$ equal to $\GL_2(q)$ is $1$. Notice that if $G_i'=\GL_2(q)$, then $G_i=\GL_2(q)$ and $r \ge 2.$ Hence there exists unique $j \in \{1, \ldots, k+1\} \backslash \{i_1, \ldots, i_s\}$ such that $G_j$ is $\GL_2(q)$.  Let $y_j=I_2$  and let $x_i, x'_i, y_i, y'_i$ for $i \in \{1, \ldots, k\} \backslash \{j\}$  be as defined before {\bf Case (2.1)}. Let $x=\diag[x_{i_1}',\ldots, x_{i_s}', x_{k+1}  \ldots, x_1]$ and $y=\diag[y_{i_1}', \ldots, y_{i_s}', y_{k+1}, \ldots, y_k] \cdot a$. It is easy to see that $(2b)$   of Lemma \ref{GRdiag} holds.  
\end{proof}

Now we prove Theorem \ref{theoremGR}.

\begin{proof}[Proof of Theorem {\rm \ref{theoremGR}}]
Let $U=U_1$, $W=W_1$ and $m=m_1.$ If  $Q \le V$ has dimension $r$, then we write 
\begin{equation} \label{Qnotation}
Q= \left \langle 
\begin{pmatrix}
u_1 \\
 \vdots \\
u_r
\end{pmatrix}
\right \rangle,
\end{equation}
where $u_1, \ldots, u_r \in V$ form a basis of $Q$.

The proof splits into two cases: when $(1)$ and $(2)$ of Lemma \ref{GRdiag} holds respectively.

\medskip

{\bf Case 1.} Assume that $(1)$ of Lemma \ref{GRdiag} holds. We study two subcases:
\begin{description}[before={\renewcommand\makelabel[1]{\bfseries ##1}}]
\item[{\bf Case (1.1)}] $m_i=1$ for $i \in \{1, \ldots, k\}$ and $n_{k+1} \in \{0,1\}$;
\item[{\bf Case (1.2)}] $m_i\ge 2$ for some $i \in \{1, \ldots, k\}$ or $n_{k+1}\ge 2.$
\end{description}

\medskip

{\bf Case (1.1).} Assume that $m_i=1$ for $i \in \{1, \ldots, k\}$ and $n_{k+1} \in \{0,1\}.$ Let $x=a$ where $a$ is as in \eqref{adefGR}, so $\det x = \pm 1$. Notice that if $\varphi \in S^x,$ then $\varphi$  has shape \eqref{GRvarphishape} with 
\begin{equation}
\label{GRginsy}
g= \begin{pmatrix}
g_{i_1}' & & & & & & \\
* & \ddots & & & & &\\
* & *&g_{i_s}' &  & & &\\
* & &* &g_{i_{k+1}} & & & \\
* &\ldots & &* & g_{k} & & \\
* &\ldots &\ldots& \ldots &*  & \ddots &\\
* &\ldots &\ldots &\ldots & &* &  g_{1}
\end{pmatrix}.
\end{equation}
So, if $\varphi \in S \cap S^x,$ then $g$ is diagonal.

 Let $y ,z \in \SL_n(q)$ be as in the proof of Proposition \ref{ni1}. Let us show that if $\varphi \in S \cap S^x \cap S^y$, then $\varphi \in S \cap \Gamma.$ Assume that $\varphi \notin \Gamma,$ so $l=1$ in  \eqref{GRvarphishape}.  Since $S$ stabilises $(U,W),$ $S^y$ stabilises $(U,W)y$. Therefore, $$(U,W)y \varphi =(U,W)y$$
and $Uy=Wy\varphi$ since $\dim Uy= \dim Wy\varphi.$ With respect to $\beta$,
$$Uy= \left \langle 
\begin{pmatrix}
1, \ldots, 1 \\
\end{pmatrix}
\right \rangle.$$ On the other hand, 
$$Uy=Wy \varphi =W y (\iota_{\beta} a) \cdot (\phi_{\beta})^j \cdot g =W \iota_{\beta} y^{\iota_{\beta}}a\cdot (\phi_{\beta})^j \cdot g= W'y^{\iota_{\beta}}a\cdot (\phi_{\beta})^j \cdot g,$$
where $W'$ is spanned by $\beta \backslash (W \cap \beta).$ 

 If $U \cap W = \{0\},$ then $W'=U.$ It is now easy to see that 
$$W'y^{\iota_{\beta}}a\cdot (\phi_{\beta})^j \cdot g = 
\begin{cases} \left \langle 
\begin{pmatrix}
1,0, \ldots, 0 \\
\end{pmatrix}
\right \rangle \text{ if } M_2 \text{ is of type } \GL_1(q) \oplus \GL_{n_2-1}(q); \\
\left \langle 
\begin{pmatrix}
0, \ldots,1, 0 \\
\end{pmatrix}
\right \rangle \text{ if } M_2 \text{ is of type } P_{1, n_2-1}
\end{cases} $$
since $g$ is diagonal. So, $Uy \ne Wy \varphi$ which is a contradiction. Hence $\varphi \in \Gamma.$

If $U\le W$, then $W'= \left \langle 
\begin{pmatrix}
1, 0, \ldots, 0 \\
\end{pmatrix}
\right \rangle$, $W'y^{\iota_{\beta}}= \left \langle 
\begin{pmatrix}
1, 0, \ldots, -1 \\
\end{pmatrix}
\right \rangle$ and  $$W'y^{\iota_{\beta}}a\cdot (\phi_{\beta})^j \cdot g=\left \langle 
\begin{pmatrix}
\alpha_1, 0, \ldots,0, \alpha_2 \\
\end{pmatrix}
\right \rangle$$ for some $\alpha_1, \alpha_2 \in \mathbb{F}_q^*.$ So, $Uy \ne Wy \varphi$ which is a contradiction. Hence $\varphi \in \Gamma.$

Therefore, $S\cap S^x \cap S^y \cap S^z = \tilde{S}\cap \tilde{S}^x \cap \tilde{S}^y \cap \tilde{S}^z \le Z(\GL_n(q))$ by Proposition \ref{ni1}. Notice that if $\det a=-1$, then we can take $x= \diag(-1, 1, \ldots, 1) \cdot a$ and the argument above still works, so we can assume $x,y,z \in \SL_n(q).$

\medskip

{\bf Case (1.2).}  Let $x,y$ be as in $(1)$ of Lemma \ref{GRdiag}. Assume that  $m_i\ge 2$ for some $i \in \{1, \ldots, k\}$ or $n_{k+1}\ge 2.$ So, if $\varphi \in (S \cap S^x \cap S^y) \cap  \Gamma$, then  there exists $r \in \{1, \ldots, n\}$ such that  $\alpha_r =\alpha_{r+1}$ in \eqref{GRvarphdiagc1}. We choose $r$ to be minimal such that $\alpha_r =\alpha_{r+1}$ for all such $\varphi.$

 By  \eqref{GRPmnmdiag}, \eqref{GRSy} and similar arguments for the situation when $U \cap W = \{0\}$ , if $\varphi \in S \cap S^y$, then 
\begin{equation}
\label{GRvarphisy}
\varphi = (\iota_{\beta} a )^l \cdot (\phi_{\beta}^j) \cdot g
\end{equation}
where $l \in \{0,1\},$ $j \in \{0, 1, \ldots, f-1\},$
$$a = \begin{cases}
I_n  &  \text{ if } M_1 \text{ is of type } \GL_m(q) \oplus \GL_{n-m(q)};\\
a(n,m) & \text{ if } M_1 \text{ is of type } P_{m, n-m}

\end{cases}
$$
and
$$g = 
\begin{cases}
\diag[g_2, g_1] &  \text{ if } M_1 \text{ is of type } \GL_m(q) \oplus \GL_{n-m(q)};\\
\diag[g_1', g_2, g_1] & \text{ if } M_1 \text{ is of type } P_{m, n-m}.
\end{cases}$$
Here $g_1 \in \GL_m(q)$, $g_2 \in \GL_{n-m}(q)$ in the first option and  $g_1, g_1' \in \GL_m(q)$, ${g_2 \in \GL_{n-2m}(q)}$ in the second. Our consideration of {\bf Case (1.2)} splits into two subcases:
when  $U \ne W$
and  $U=W$ respectively.

\medskip

{\bf Case (1.2.1).} Assume that $U \ne W$. Let $\theta$ be a generator of $\mathbb{F}_q^*$ and let $z \in \SL_n(q)$ be defined as follows:
\begin{equation}
\label{GRzdef1}
\begin{aligned}
(v_i)z & = v_i &&\text{ for } i \in \{1, \ldots, n-m\}; \\
(v_{n-m+1})z & =  \sum_{i=1}^{n-m} v_i -v_r + \theta v_{r} + v_{n-m+1}; && \text{ if } r \le n-m\\
(v_{n-m+1})z & =  \sum_{i=1}^{n-m-1} v_i  + \theta v_{n-m} + v_{n-m+1}; && \text{ if } r \ge n-m+1\\
(v_i)z & =  v_{n-m} + v_{i} && \text{ for } i \in \{n-m+2, \ldots, n\}.
\end{aligned} 
\end{equation}

Let $\varphi \in S \cap S^x \cap S^y \cap S^z$, so $\varphi$ has shape \eqref{GRvarphisy}.   Assume that $\varphi \notin \Gamma,$ so $l=1.$ Since $\varphi \in S^z$, it stabilises $(U,W)z,$ so $(U,W)z \varphi =(U,W)z$ and, therefore, 
$$Uz=Wz \varphi = W'z^{\iota_{\beta}}a\cdot (\phi_{\beta})^j \cdot g,$$
where $W'$ is spanned by $\beta \backslash (W \cap \beta).$  

 With respect to $\beta$, using the notation \eqref{Qnotation},
$$Uz= \left \langle 
\left(\begin{array}{ccccccc|ccccc}
1 & \ldots & 1 & \theta &1 & \ldots & 1     &             1 & & &  \\
1 & \ldots & 1 &  1    &1 & \ldots & 1       &           & 1 & &  \\
\vdots &   &   &       &  &        &\vdots    &          & &\ddots & \\
1 & \ldots & 1 &  1    &1 & \ldots & 1         &         & & & 1 
\end{array}\right)
\right \rangle$$
where $\theta$ in the first line is either in the $r$-th or $(m-n)$-th column, and the part after the vertical line forms $I_m.$

 If $U \cap W=\{0\},$ then $W'=U$ and it is easy to see that $Wz \varphi = W'z^{\iota_{\beta}}a\cdot (\phi_{\beta})^j \cdot g= U,$ since $g$ stabilises $U$ by \eqref{GRvarphisy}. So, $Uz \ne Wz \varphi$ which is a contradiction. Hence $\varphi \in \Gamma.$

If $U \le W$, then $W'= \langle (I_m \mid 0_{ m \times (n-m)}) \rangle$, so 
$$W'z^{\iota_{\beta}}a= \left \langle 
(A \mid 0_{n \times (n-2m)} \mid I_n)
\right \rangle$$
where $A$ is $m \times m$ matrix with entries $-1$ or $- \theta$, and $-\theta$ can occur at most once.
Therefore, $W'z^{\iota_{\beta}}a \cdot g = \langle (A g_1'\mid0_{m \times (n-2m)}\mid g_1) \rangle$. Notice that $n-2m\ge 1$ since $U \ne W.$ So, $Uz \ne Wz \varphi$ which is a contradiction. Hence $\varphi \in \Gamma.$

Therefore, $\varphi = (\phi_{\beta})^j \cdot \diag(\alpha_1, \ldots, \alpha_n)$ as in $(1)$ of Lemma
\ref{GRdiag}. Since $\varphi \in S \cap S^z \cap \Gamma,$ it stabilises $U$ and $Uz.$  

Consider $((v_{n-m+1})z) \varphi.$ First, let $r\le n-m,$ so %$m \ge 2$ and
\begin{equation*}((v_{n-m+1})z) \varphi=
\begin{cases}
\left(\sum_{\substack{i \in \{1, \ldots, n-m\}\\ i \ne r}} \alpha_i v_i \right) + \alpha_r \theta^{p^j} v_r + \alpha_{n-m+1} v_{n-m+1};\\
\sum_{i=1}^m \delta_i (v_{n-m+i})z
\end{cases}
\end{equation*}
for some $\delta_i \in \mathbb{F}_q.$ Since $((v_{n-m+1})z) \varphi$ has no $v_i$ for $i \in \{n-m+2, \ldots, n\}$ in the decomposition with respect to $\beta$, 
$$\delta_2 = \ldots = \delta_m =0$$
and $((v_{n-m+1})z) \varphi= \delta_1 ((v_{n-m+1})z).$ Hence 
$$\delta_1 = \alpha_r \theta^{p^j-1}=\alpha_{r+1} = \alpha_i \text{ for } i \in \{1, \ldots, n-m+1\} \backslash \{r,r+1\},$$
so $j=0$ and $\varphi \in Z(\GL_n(q)).$

Now let $r\ge n-m+1,$ so $m\ge 2$ and 
\begin{equation*}((v_{n-m+1})z) \varphi=
\begin{cases}
\left(\sum_{i=1}^{n-m-1} \alpha_i v_i \right) + \alpha_{n-m} \theta^{p^j} v_{n-m} + \alpha_{n-m+1} v_{n-m+1};\\
\sum_{i=1}^m \delta_i (v_{n-m+i})z
\end{cases}
\end{equation*}
for some $\delta_i \in \mathbb{F}_q.$ Since $((v_{n-m+1})z) \varphi$ has no $v_i$ for $i \in \{n-m+2, \ldots, n\}$ in the decomposition with respect to $\beta$, 
$$\delta_2 = \ldots = \delta_m =0$$
and $((v_{n-m+1})z) \varphi= \delta_1 ((v_{n-m+1})z).$ Hence 
$$\delta_1 = \alpha_{n-m} \theta^{p^j-1}=\alpha_{n-m+1} = \alpha_i \text{ for } i \in \{1, \ldots, n-m-1\}.$$
The same arguments for $((v_{n-m+2})z) \varphi$ show that $\alpha_{n-m+2}=\alpha_{n-m},$ so, since $\alpha_{n-m+1}=\alpha_{n-m+2}$ by $(1)$ of Lemma \ref{GRdiag}, we obtain $\theta^{p^j-1}=1.$ Hence $j=0$ and $\varphi \in Z(\GL_n(q)).$
 



\medskip
{\bf Case (1.2.2).} Assume that $U=W$, so $m=n/2$ and $M$ is of type $P_{n/2,n/2}.$ In particular, $n \ge 4,$ since $n \ge 3$ by the assumption of the theorem.  Let $\theta$ be a generator of $\mathbb{F}_q^*$ and let $z \in \SL_n(q)$ be defined as follows: 
\begin{equation}
\label{GRzdef2}
\begin{aligned}
(v_i)z= &v_i &&\text{ for } i \in \{1, \ldots, n-m\}; \\
(v_{m-n+1})z= & \theta v_{n-m} + v_{m-n+1}; &&\\
(v_i)z= & v_{n-m} + v_{i} && \text{ for } i \in \{n-m+2, \ldots, n\}.
\end{aligned} 
\end{equation}
Arguments, analogous to  ones in {\bf Case (1.2.1)},  show that  $S \cap S^x \cap S^y \cap S^z \le Z(\GL_n(q))$.
\bigskip

{\bf Case 2.}  Assume that $(2)$ of Lemma \ref{GRdiag} holds.
 For $n \in \{3,4\}$ the theorem follows by computation, so  we may assume that $n \ge 5$.


We adopt notation from the proof of Lemma \ref{GRdiag}, in particular $x,y$, $S_i$, $G_i$.
Let $(2 \times 2)$ blocks (corresponding to the $S_i$ and the $S_i'$ lying in $\GL_2(q)$) on the diagonal in matrices of $\tilde{S}=S \cap \Gamma$ occur in the rows $$(j_1, j_1+1), (j_2, j_2 +1), \ldots, (j_r, j_r+1).$$ Let $\Lambda=\Lambda_1 \cup \Lambda_2$ where 
$\Lambda_1=\{j_1, j_2, \ldots, j_r\}$ and $\Lambda_2=\{j_1+1, j_2+1, \ldots, j_r+1\}.$
 Let $U=U_1$, $W=W_1$ and $m=m_1.$  
 Notice, that if $\varphi \in S^y$, then it has shape \eqref{GRvarphishape} where $g$ has shape \eqref{GRginsy} with $g_i, g_i' \in \GL_{m_i}$ for $i \in \{1, \ldots, k\}$ and $g_{k+1} \in \GL_{n_{k+1}}(q).$ So, if $\varphi \in S \cap S^y$, then it has shape \eqref{GRvarphishape} with 
 \begin{equation}
 \label{GRgdiagcase2}
 g = \diag[g_{i_1}', \ldots, g_{i_s}', g_{k+1}, g_k, \ldots, g_1].
 \end{equation}
 
 Our consideration of {\bf Case 2} splits into two subcases: when $(2a)$ and $(2b)$ of Lemma \ref{GRdiag} holds respectively.

\medskip

{\bf Case (2.1).} Assume that $(2a)$ of Lemma \ref{GRdiag} holds. We consider two subcases: when $m \ge 2$ and $m=1$.

{\bf Case (2.1.1).} Assume that $m \ge 2.$ Let $z \in \SL_n(q)$ be defined as follows
\begin{equation}
\label{GRzdefq23}
\begin{aligned}
(v_i)z & =  v_i &&\text{ for } i \in \{1, \ldots, n-m\}; \\
%(v_{n-m+1})z= & \sum_{i =1}^{n-m} v_i + v_{n-m+1} && \text{ if } (a) \text{ holds};\\
%(v_{n-m+1})z= & \sum_{i \in \{1, \ldots, n-m\} \backslash \Lambda_1} v_i + v_{n-m+1} && \text{ if } (b) \text{ holds}; \\
(v_{n-m+1})z & =  \left(\underset{i \in \{1, \ldots, n-m\} \backslash \{j_1\}}{\sum} v_i \right) + v_{n-m+1};\\
(v_{n-m+2})z & =  \left(\underset{i \in \{1, \ldots, n-m\} \backslash \Lambda_2}{\sum} v_i \right)   + v_{n-m+2};\\
(v_i)z & =  \sum_{j=1}^{n-m} v_{j} + v_{i} && \text{ for } i \in \{n-m+2, \ldots, n\}.
\end{aligned} 
\end{equation}

Let $\varphi \in S \cap S^x \cap S^y \cap S^z$, so $\varphi$ has shape \eqref{GRvarphishape} where $g$ has shape \eqref{GRgdiagcase2}.   Assume that $\varphi \notin \Gamma,$ so $l=1.$ Since $\varphi \in S^z$, it stabilises $(U,W)z,$ so $(U,W)z \varphi =(U,W)z$ and, therefore, 
$$Uz=Wz \varphi = W'z^{\iota_{\beta}}a \cdot g,$$
where $W'$ is spanned by $\beta \backslash (W \cap \beta).$  

 With respect to $\beta$,
$$Uz= \left \langle 
\left(\begin{array}{ccc|ccccc}
\lambda_1 & \ldots & \lambda_{n-m}   &           1 & & & & \\
\mu_1     & \ldots & \mu_{n-m}       &           & 1 & & & \\
  1       & \ldots &  1              &           &   &1 & &  \\
\vdots    &        &\vdots           &           &  & &\ddots & \\
1         & \ldots &  1              &           &  & & & 1 
\end{array}\right)
\right \rangle$$
where $\lambda_i, \mu_i \in \{0,1\}$ according to \eqref{GRzdefq23}, so for each $i \in \{1, \ldots, n-m\}$ at least one of $\lambda_i$ and $\mu_i$ is $1$,  and the part after the vertical line forms $I_m.$

 If $U \cap W=\{0\},$ then $W'=U$ and it is easy to see that $Wz \varphi = W'z^{\iota_{\beta}}a\cdot (\phi_{\beta})^j \cdot g= U,$ since $g$ stabilises $U$ by \eqref{GRvarphishape}. So, $Uz \ne Wz \varphi$ which is a contradiction. Hence $\varphi \in \Gamma.$

If $U \le W$, then $W'= \langle (I_m  \mid  0_{(m \times n-m)}) \rangle$, so 
$$W'z^{\iota_{\beta}}a= \left \langle 
(A \mid 0_{m \times (n-2m)} \mid I_n)
\right \rangle$$
where $A$ is $m \times m$ matrix with entries $-1$ ,$-\lambda_i$ and $- \mu_i$.
Therefore, $$W'z^{\iota_{\beta}}a \cdot g = \langle (A g_1'\mid 0_{m \times (n-2m)}\mid g_1) \rangle.$$ Notice that $n-2m\ge 1,$ since otherwise  $U = W,$ so $m=2$ and $n=4.$ Thus, $Uz \ne Wz \varphi$ which is a contradiction. Hence $\varphi \in \Gamma.$

Therefore, $\varphi = g= \diag[g_{i_1}, \ldots,g_{i_s}, g_{k+1}, g_k, \ldots,  g_1]$ as in $(a)$ of $(2)$ of Lemma
\ref{GRdiag}. Specifically, let $(v_i)\varphi =\alpha_i v_i$ for $i \in \{1, \ldots, n\} \backslash \Lambda_1$ and let $(v_i)\varphi = \alpha_i v_i + \gamma_i v_{i+1}$ for $i \in \Lambda_1$ with $\alpha_i, \gamma_i \in \mathbb{F}_q.$

 Since $\varphi \in S \cap S^z \cap \Gamma,$ it stabilises $U$ and $Uz.$  
Therefore, $((v_{n-m+2})z) \varphi$ is
\begin{equation*}
\label{vnm2zph}
\left(\underset{i \in \{1, \ldots, n-m\} \backslash \Lambda_2}{\sum} \alpha_i v_i \right) +
\left(\underset{i \in \{1, \ldots, n-m\} \cap \Lambda_2}{\sum} \gamma_{i-1} v_i \right) + \alpha_{n-m+2} v_{n-m+2},
\end{equation*}
and
$$((v_{n-m+2})z) \varphi= \sum_{i=1}^m \delta_i (v_{n-m+i})z$$
for some $\delta_i \in \mathbb{F}_q.$ Since $((v_{n-m+2})z) \varphi$ does not contain $v_i$ for $i \in \{n-m+1,n-m+3, \ldots, n\}$ in the decomposition with respect to $\beta$, 
$$\delta_1= \delta_3 = \ldots = \delta_m =0$$
and $((v_{n-m+1})z) \varphi= \delta_1 ((v_{n-m+1})z).$ Hence 
$$\gamma_{j_1} = \ldots = \gamma_{j_r}=0,$$
so $\varphi= \diag(\alpha_1, \ldots, \alpha_n)$ where 
\begin{equation}
\label{GRalphasvnm2}
\alpha_i=\alpha_{n-m+2} \text{ for } i \in \{1, \ldots, n-m\} \backslash \Lambda_2.
\end{equation} 

Consider 
\begin{equation*}((v_{n-m+1})z) \varphi=
\begin{cases}
\underset{i \in \{1, \ldots, n-m\} \backslash \{j_1\}}{\sum} \alpha_i v_i + \alpha_{n-m+1} v_{n-m+1} +\underline{\gamma_{n-m+1} v_{n-m+1}}\\
\sum_{i=1}^m \delta_i (v_{n-m+i})z
\end{cases}
\end{equation*}
for some $\delta_i \in \mathbb{F}_q.$ The underlined part is present only if $m=2.$ Since $((v_{n-m+1})z) \varphi$ contains neither $v_{j_1}$ (notice that $j_1>n-m$ since if $(2a)$ of Lemma \ref{GRdiag} holds, then $r \ge 2$) nor $v_i$ for $i \in \{n-m+3, \ldots, n\}$ in the decomposition with respect to $\beta$, 
$$\delta_2 = \ldots = \delta_m =0.$$
Here $\delta_2=0$ since $((v_{n-m+2})z) $ contains $v_{j_1}$ in the decomposition with respect to $\beta$ and $((v_{n-m+1})z)\varphi$ does not.
Thus $((v_{n-m+1})z) \varphi= \delta_1 ((v_{n-m+1})z)$ and 
$$\alpha_i = \alpha_{n-m+1} \text{ for } i \in \{1, \ldots, n-m\} \backslash \{j_1\}.$$
Combined with \eqref{GRalphasvnm2}, it implies $\varphi \in Z(\GL_n(q)).$ 

\medskip

{\bf Case (2.1.2).} Assume that $m_1=1$ and let $t$ be the smallest $i \in \{2, \ldots, k\}$ such that $m_i\ge 2.$ Such $t$ exists since otherwise $m_i=1$ for $i \in \{1, \ldots, k\}$ and $n_{k+1}=2,$ so $(2b)$ of Lemma \ref{GRdiag} holds. Let $d_i=\sum_{j=1}^i m_j$ for $i \in \{1, \ldots, k\}$. Let $z \in \SL_n(q)$ be defined as follows:
\begin{equation}
\label{GRzdefq23mr}
\begin{aligned}
(v_i)z & =  v_i &&\text{ for } i \in \{1, \ldots, n-d_t\}; \\
(v_{n-d_t+1})z & =  \left(\underset{i \in \{1, \ldots, n-d_t\} \backslash \{j_1\}}{\sum} v_i \right) + v_{n-d_t+1};\\
(v_{n-d_t+2})z & =  \left(\underset{i \in \{1, \ldots, n-d_t\} \backslash \Lambda_2}{\sum} v_i \right)   + v_{n-d_t+2};\\
(v_i)z & =  v_i && \text{ for } i \in \{n-d_{t}+3, \ldots, n-1\};\\
(v_n)z & =  \sum_{i=1}^n v_i.
\end{aligned} 
\end{equation}
Arguments, analogous to  ones in {\bf Case (2.1.1)},  show that  $S \cap S^x \cap S^y \cap S^z \le Z(\GL_n(q))$.



\bigskip

{\bf Case (2.2).} Finally, to complete the proof, let us assume that $(2b)$ of Lemma \ref{GRdiag} holds. We consider two subcases: when $m=2$ and $m \ne 2$.

\medskip

{\bf Case (2.2.1).} Assume that $m=2,$ so if $g \in S \cap S^x \cap S^y \cap \GL_n(q),$ then $g_1 \in \GL_2(q)$ and $g_i$, $g_i'$ are scalar for $i \in \{2, \ldots, k+1\}.$ We may assume $U \cap W=0$ since otherwise the number of $G_i$ and $G_i'$ equal to $\GL_2(q)$ is at least 2 and $(2a)$  of Lemma \ref{GRdiag} holds. Let $z \in \SL_n(q)$ be defined as follows:
\begin{equation}
\label{GRzdefcase221}
\begin{aligned}
(v_i)z & =  v_i &&\text{ for } i \in \{1, \ldots, n-2\}; \\
(v_{n-1})z & =  \sum_{i=2}^{n-2} v_i + v_{n-1};\\
(v_{n})z & =  v_1+ \sum_{i=3}^{n-2} v_i   + v_{n}.\\
\end{aligned} 
\end{equation}
Let $\varphi \in S \cap S^x \cap S^y \cap S^z$, so $\varphi= (\iota_{\beta}a)^l \cdot g$ where $g=\diag[g_{i_1},\ldots, g_{i_s}, g_{k+1}, \ldots, g_1]$ where $g_i$ and $g_i'$ are as in \eqref{GRguppdiag}.  
 Assume that $\varphi \notin \Gamma,$ so $l=1.$ Since $\varphi \in S^z$, it stabilises $(U,W)z,$ so $(U,W)z \varphi =(U,W)z$ and, therefore, 
$$Uz=Wz \varphi = W'z^{\iota_{\beta}}a \cdot g,$$
where $W'$ is spanned by $\beta \backslash (W \cap \beta).$  

 With respect to $\beta$,
$$Uz= \left \langle 
\left(\begin{array}{ccccc|cc}
0  & 1& 1 & \ldots & 1   &           1 &0  \\
1  & 0& 1     & \ldots & 1       &    0       & 1 
\end{array}\right)
\right \rangle.$$
 Since $U \cap W=\{0\},$ we obtain $W'=U$ and it is easy to see that $Wz \varphi = W'z^{\iota_{\beta}}a\cdot g= U,$ since $g$ stabilises $U$. So, $Uz \ne Wz \varphi$ which is a contradiction. Hence $\varphi =g\in \Gamma.$ Therefore, $g=\diag[A,g_1]$ where $A=\diag(\alpha_1, \ldots, \alpha_{n-2})$ for some $\alpha_i \in \mathbb{F}_q^*$ and $$g_1 = \left( \begin{matrix} \delta_1 & \delta_2 \\ \delta_3 & \delta_4  \end{matrix} \right) \in \GL_2(q).$$ Since $g \in S^z$, it stabilises $Uz.$
 
 Consider 
\begin{equation}
\label{GRc221vn1z}
 ((v_{n-1})z)g =
 \begin{cases}
 \sum_{i=2}^{n-2} \alpha_i v_i + \delta_1 v_{n-1} + \delta_2 v_n;\\
 \lambda_1 (v_{n-1})z + \lambda_2 (v_{n})z
 \end{cases}
 \end{equation}
 for some $\lambda_{i} \in \mathbb{F}_q.$ Since there is no $v_1$ in the first line of \eqref{GRc221vn1z}, $\lambda_2=0,$ so
 $$\alpha_2 = \ldots = \alpha_{n-2}=\delta_1 \text{ and } \delta_2=0.$$
 The same arguments for $((v_{n})z)g$ show that $\delta_3=0$ and, since $n \ge 5$,
 $$\alpha_1=\alpha_{n-2}=\delta_4.$$
 Hence $g$ is scalar and $\varphi \in Z(\GL_n(q)).$
 
 \medskip
 
 {\bf Case (2.2.2).} Assume $m\ne 2.$  
  We may assume that $n_{k+1}\ne 2$. Indeed, if $n_{k+1}=2$, then  $n_{k}=3$ since $m_k\le n_k/2$ and there is only one $G_i$ equal to $\GL_2(q)$ for $i \in \{1, \ldots, k+1\}.$ Therefore, $$S_k=S|_{_{V_k}}\cap \GL(V_k)=\GL_2(q) \times \GL_1(q) \le \GL(V_k)=\GL_3(q)$$ and, using computation, we obtain that  there are $x_k, y_k \in \SL_n(q)$ such that $S_k \cap S_k^{x_k} \cap S_k^{y_k} \le Z(\GL_3(q))$, so the conclusion of  $(1)$ of Lemma \ref{GRdiag} holds and the theorem  holds by {\bf Case  1}. 
  
  Therefore, $j_1\ge 3$ where $j_1$ is as defined in the beginning of {\bf Case  2}, so the $(2 \times 2)$ block (corresponding to the $S_i$ lying in $\GL_2(q)$) on the diagonal in matrices of $\tilde{S}=S \cap \Gamma$ occurs in the rows $(j_1, j_1+1)$. Recall that if $h \in  S \cap S^x \cap S^y \cap  \GL_n(q)$, then $$h= \diag[\alpha_1, \ldots, \alpha_{j_1-1},A, \alpha_{j_1+2}, \ldots, \alpha_n]$$ where $\alpha_i \in \mathbb{F}_q^*$ and $A= \left( \begin{smallmatrix} \delta_1 & \delta_2 \\ \delta_3 & \delta_4  \end{smallmatrix} \right) \in \GL_2(q).$ Let $z \in \SL_n(q)$ be defined as follows:
\begin{equation}
\label{GRzdefcase222}
\begin{aligned}
(v_i)z & =  v_i &&\text{ for } i \in \{1, \ldots, n\} \backslash \{j_1, j_1+1, n\}; \\
(v_{j_1})z & =  v_1 + v_{j_1};\\
(v_{j_1+1})z & =  v_2 + v_{j_1+1};\\
(v_{n})z & =   \sum_{i=1}^{n-m} v_i   + v_{n}.\\
\end{aligned} 
\end{equation}
Let $\varphi \in S \cap S^x \cap S^y \cap S^z$, so $\varphi= (\iota_{\beta}a)^l \cdot g$ with $g=\diag[g_{i_1},\ldots, g_{i_s}, g_{k+1}, \ldots, g_1]$ where $g_i$ and $g_i'$ are as in \eqref{GRguppdiag}.  
 Assume that $\varphi \notin \Gamma,$ so $l=1.$ Since $\varphi \in S^z$, it stabilises $(U,W)z,$ so $(U,W)z \varphi =(U,W)z$ and, therefore, 
$$Uz=Wz \varphi = W'z^{\iota_{\beta}}a \cdot g,$$
where $W'$ is spanned by $\beta \backslash (W \cap \beta).$  

 With respect to $\beta$,
$$Uz= \left \langle 
\left(\begin{array}{ccc|ccccc}
 0        & \ldots & 0               &           1 & & &  \\
\vdots    &        &\vdots           &           & \ddots  & &  \\
0         & \ldots & 0               &            & &\ddots &  \\
1         & \ldots &  1              &           &  & & 1 
\end{array}\right)
\right \rangle$$
where the part after the vertical line forms $I_m.$

 If $U \cap W=\{0\},$ then $W'=U$ and it is easy to see that $Wz \varphi = W'z^{\iota_{\beta}}a \cdot g= U,$ since $g$ stabilises $U$ by \eqref{GRvarphisy}. So $Uz \ne Wz \varphi$ which is a contradiction. Hence $\varphi \in \Gamma.$
 
 If $U \le W,$ then 
 \begingroup
\allowdisplaybreaks
 \begin{align*}
 Wz\varphi & = W'z^{\iota_{\beta}}ag \\ & = 
 \left\langle I_m\mid 0_{m \times (n-m)} \right\rangle \, z^{\iota_{\beta}}ag \\
& =\Scale[0.95]{\left \langle \left(\begin{array}{ccccc|cccccc|cccc}
1 &  &   &        &          & 0 & \ldots 0 &-1 & 0 & 0 & \ldots 0 &       0  & 0 & \ldots &  0 \\         
  &1 &   &        &          & 0 & \ldots 0 & 0 &-1 & 0 & \ldots 0 &       0  & 0 & \ldots &  0 \\
  &  &1  &        &          & 0 & \ldots 0 & 0 &0  & 0 & \ldots 0 &       -1 & 0 & \ldots &  0 \\
  &  &   & \ddots &          & \vdots &     &   &   &   &          &   \vdots &   &        &    \\
  &  &   &        &1         & 0 & \ldots 0 & 0 &0  & 0 & \ldots 0 &       -1 & 0 & \ldots &  0 \\ 
\end{array}\right)
\right \rangle}\, a g\\
& = \Scale[0.95]{ \left \langle \left(\begin{array}{cccc|cccccc|ccccc}
 0  & 0 & \ldots &  0          & 0 & \ldots 0 &-1 & 0 & 0 & \ldots 0 &       1 &  &   &        & \\         
 0  & 0 & \ldots &  0        & 0 & \ldots 0 & 0 &-1 & 0 & \ldots 0 &         &1 &   &        &   \\
 -1 & 0 & \ldots &  0         & 0 & \ldots 0 & 0 &0  & 0 & \ldots 0 &         &  &1  &        &  \\
  &   \vdots &   &        &            & \vdots &     &   &   &   &           &  &   & \ddots &  \\
  -1 & 0 & \ldots &  0        & 0 & \ldots 0 & 0 &0  & 0 & \ldots 0 &        &  &   &        &1  \\ 
\end{array}\right)
\right \rangle}\, g
 \end{align*}
 \endgroup
 where $-1$ in the first row is in the $j_1$ entry, and $-1$ in the second row is in the $j_1+1$ entry. The result of the action of $g$ on the first two rows is 
 $$\langle 0_{2\times 1}, \ldots, 0_{2\times 1}, -g_t, 0_{2\times 1}, \ldots, 0_{2\times 1}, g_1^{(1,2)} \rangle $$ where $t \in \{1, \ldots, k+1\}$ is such that $G_t=\GL_2(q)$ and $g_1^{(1,2)}$ is the matrix formed by the first two rows of $g_1.$ It is easy to see that  two such vectors cannot lie in $Uz$, so  $Uz \ne Wz \varphi$ which is a contradiction. Hence $\varphi =g\in \Gamma.$ 
 
 Therefore $g=\diag[A,g_t, B]$ where $$A=\diag(\alpha_1, \ldots, \alpha_{j_1-1}), \text{ }B=\diag(\alpha_{j_1+2}, \ldots, \alpha_{n}),$$  for some $\alpha_i \in \mathbb{F}_q^*$ and $$g_t = \left( \begin{matrix} \delta_1 & \delta_2 \\ \delta_3 & \delta_4  \end{matrix} \right) \in \GL_2(q).$$ Since $g \in S^z$, it stabilises $Uz.$
 Notice that $S \cap \GL_n(q)$ also stabilises $\langle v_{j_1}, \ldots, v_n \rangle,$ so $g \in S^z \cap \GL_n(q)$ stabilises $\langle v_{j_1}, \ldots, v_n \rangle z$. 
 
 Consider 
\begin{equation}
\label{GR222vj1}
((v_{j_1})z)g= 
 \begin{cases}
 \alpha_1 v_1 + \delta_1 v_{j_1} + \delta_2 v_{j_1+1};\\
 \sum_{i=j_1}^n \lambda_i v_i
 \end{cases}
\end{equation}
 for some $\lambda_i \in \mathbb{F}_q.$ Since the first line of \eqref{GR222vj1} contains no terms with $v_2$ and $v_i$ for $i \ge j_1+2$, we obtain   
 $((v_{j_1})z)g= \alpha (v_{j_1})z $ for some $\alpha \in \mathbb{F}_q^*.$ Therefore, $\delta_2=0$ and $\alpha= \alpha_1 = \delta_1.$ The same arguments applied to $(v_{j_1+1})z$ show that
 $\delta_4=0$ and $\alpha_2=\delta_3.$ 
 
   The same arguments applied to $(v_{n})z$ show that
$\alpha= \alpha_1 = \ldots =\alpha_n,$ so $g$ is scalar and $\varphi \in Z(\GL_n(q)).$
\end{proof}

\section*{Acknowledgements}
Much of the work presented in this paper was done while I was a PhD student at the University of Auckland. I thank my supervisors Eamonn O'Brien and Jianbei An for all their guidance and inspiration. I also thank Professor Timothy  Burness and Professor Peter Cameron, who were  examiners of my PhD thesis, and the anonymous referees of this paper for their constructive comments and corrections.



\bibliographystyle{abbrv}
\bibliography{SolLin.bib}

\end{document}
\endinput
%%

\begin{thebibliography}{10}

\bibitem{asch}
M.~Aschbacher.
\newblock On the maximal subgroups of the finite classical groups.
\newblock {\em Invent. Math.}, 76(3):469--514, 1984.

\bibitem{bab}
L.~Babai, A.~J. Goodman, and L.~Pyber.
\newblock Groups without faithful transitive permutation representations of
  small degree.
\newblock {\em J. Algebra}, 195(1):1--29, 1997.

\bibitem{bailey}
R.~F. Bailey and P.~J. Cameron.
\newblock Base size, metric dimension and other invariants of groups and
  graphs.
\newblock {\em Bull. Lond. Math. Soc.}, 43(2):209--242, 2011.

\bibitem{bay}
A.~A. Baykalov.
\newblock Intersection of conjugate solvable subgroups in symmetric groups.
\newblock {\em Algebra Logika}, 56(2):135--149, 2017.

\bibitem{thesis}
A.~A. Baykalov.
\newblock {\em Intersection of conjugate solvable subgroups in finite classical
  groups}.
\newblock PhD thesis, The University of Auckland, 2021.
\newblock {\url{https://doi.org/10.48550/arXiv.1703.00124}}.

\bibitem{magma}
W.~Bosma, J.~Cannon, and C.~Playoust.
\newblock The {M}agma algebra system. {I}. {T}he user language.
\newblock {\em J. Symbolic Comput.}, 24(3-4):235--265, 1997.

\bibitem{maxlow}
J.~N. Bray, D.~F. Holt, and C.~M. Roney-Dougal.
\newblock {\em The maximal subgroups of the low-dimensional finite classical
  groups}, volume 407 of {\em London Mathematical Society Lecture Note Series}.
\newblock Cambridge University Press, Cambridge, 2013.

\bibitem{fpr}
T.~C. Burness.
\newblock Fixed point ratios in actions of finite classical groups, {I}.
\newblock {\em J. Algebra}, 309(1):69--79, 2007.

\bibitem{fpr2}
T.~C. Burness.
\newblock Fixed point ratios in actions of finite classical groups, {II}.
\newblock {\em J. Algebra}, 309(1):80--138, 2007.

\bibitem{fpr3}
T.~C. Burness.
\newblock Fixed point ratios in actions of finite classical groups. {III}.
\newblock {\em J. Algebra}, 314(2):693--748, 2007.

\bibitem{fpr4}
T.~C. Burness.
\newblock Fixed point ratios in actions of finite classical groups. {IV}.
\newblock {\em J. Algebra}, 314(2):749--788, 2007.

\bibitem{burness}
T.~C. Burness.
\newblock On base sizes for actions of finite classical groups.
\newblock {\em J. Lond. Math. Soc. (2)}, 75(3):545--562, 2007.

\bibitem{sfa}
T.~C. Burness.
\newblock Simple groups, fixed point ratios and applications.
\newblock In {\em Local representation theory and simple groups}, EMS Ser.
  Lect. Math., pages 267--322. Eur. Math. Soc., Z\"{u}rich, 2018.

\bibitem{burPS}
T.~C. Burness.
\newblock Base sizes for primitive groups with soluble stabilisers.
\newblock {\em Algebra Number Theory}, 15(7):1755--1807, 2021.

\bibitem{burspor}
T.~C. Burness.
\newblock On soluble subgroups of sporadic groups.
\newblock {\em Israel J. Math.}, 254(1):313--340, 2023.

\bibitem{bls}
T.~C. Burness, M.~W. Liebeck, and A.~Shalev.
\newblock Base sizes for simple groups and a conjecture of {C}ameron.
\newblock {\em Proc. Lond. Math. Soc. (3)}, 98(1):116--162, 2009.

\bibitem{spor}
T.~C. Burness, E.~A. O'Brien, and R.~A. Wilson.
\newblock Base sizes for sporadic simple groups.
\newblock {\em Israel J. Math.}, 177:307--333, 2010.

\bibitem{buturl}
A.~A. Buturlakin and M.~A. Grechkoseeva.
\newblock The cyclic structure of maximal tori in finite classical groups.
\newblock {\em Algebra Logika}, 46(2):129--156, 2007.

\bibitem{Camer}
P.~J. Cameron and W.~M. Kantor.
\newblock Random permutations: some group-theoretic aspects.
\newblock {\em Combin. Probab. Comput.}, 2(3):257--262, 1993.

\bibitem{carter}
R.~W. Carter.
\newblock {\em Finite groups of {L}ie type: Conjugacy classes and complex
  characters}.
\newblock Pure and Applied Mathematics (New York). John Wiley \& Sons, Inc.,
  New York, 1985.

\bibitem{cher}
A.~Chermak and A.~Delgado.
\newblock A measuring argument for finite groups.
\newblock {\em Proc. Amer. Math. Soc.}, 107(4):907--914, 1989.

\bibitem{Dieu}
J.~Dieudonn\'{e}.
\newblock {\em On the automorphisms of the classical groups}, volume No. 2 of
  {\em Memoirs of the American Mathematical Society}.
\newblock American Mathematical Society, Providence, RI, 1980.
\newblock With a supplement by Loo Keng Hua [Luo Geng Hua], Reprint of the 1951
  original.

\bibitem{dolfi}
S.~Dolfi.
\newblock Large orbits in coprime actions of solvable groups.
\newblock {\em Trans. Amer. Math. Soc.}, 360(1):135--152, 2008.

\bibitem{GAP4}
The GAP~Group.
\newblock {\em {GAP -- Groups, Algorithms, and Programming, Version 4.11.1}},
  2021.

\bibitem{gluck}
D.~Gluck and O.~Manz.
\newblock Prime factors of character degrees of solvable groups.
\newblock {\em Bull. London Math. Soc.}, 19(5):431--437, 1987.

\bibitem{conjaut}
D.~Gorenstein and R.~Lyons.
\newblock The local structure of finite groups of characteristic {$2$} type.
\newblock {\em Mem. Amer. Math. Soc.}, 42(276):vii+731, 1983.

\bibitem{gorlyo}
D.~Gorenstein, R.~Lyons, and R.~Solomon.
\newblock {\em The classification of the finite simple groups. {N}umber 3.
  {P}art {I}. {C}hapter {A}}, volume~40 of {\em Mathematical Surveys and
  Monographs}.
\newblock American Mathematical Society, Providence, RI, 1998.

\bibitem{PrimGrp}
A.~Hulpke, C.~Roney-Dougal, and C.~Russell.
\newblock {\em PrimGrp -- GAP Primitive Permutation Groups Library}, Version
  3.4.0, 2019 (https://gap-packages.github.io/primgrp/).

\bibitem{hump}
J.~E. Humphreys.
\newblock {\em Linear algebraic groups}.
\newblock Graduate Texts in Mathematics, No. 21. Springer-Verlag, New
  York-Heidelberg, 1975.

\bibitem{hupp}
B.~Huppert.
\newblock {\em Endliche {G}ruppen. {I}}.
\newblock Die Grundlehren der Mathematischen Wissenschaften, Band 134.
  Springer-Verlag, Berlin-New York, 1967.

\bibitem{james}
J.~P. James.
\newblock Two point stabilisers of partition actions of linear groups.
\newblock {\em J. Algebra}, 297(2):453--469, 2006.

\bibitem{kaz}
L.~S. Kazarin and S.~P. Strunkov.
\newblock On intersections of subgroups in groups.
\newblock {\em J. Math. Sci. (New York)}, 100(1):1888--1892, 2000.

\bibitem{kt}
E.~I. Khukhro and V.~D. Mazurov, editors.
\newblock {\em The {K}ourovka notebook}.
\newblock Sobolev Institute of Mathematics. Russian Academy of Sciences.
  Siberian Branch, Novosibirsk, 2018.
\newblock Unsolved problems in group theory, Nineteenth edition, March 2019
  update.

\bibitem{kleidlieb}
P.~Kleidman and M.~Liebeck.
\newblock {\em The subgroup structure of the finite classical groups}, volume
  129 of {\em London Mathematical Society Lecture Note Series}.
\newblock Cambridge University Press, Cambridge, 1990.

\bibitem{lieb}
M.~W. Liebeck and A.~Shalev.
\newblock Simple groups, permutation groups, and probability.
\newblock {\em J. Amer. Math. Soc.}, 12(2):497--520, 1999.

\bibitem{luccini}
A.~Lucchini.
\newblock On the order of transitive permutation groups with cyclic
  point-stabilizer.
\newblock {\em Atti Accad. Naz. Lincei Cl. Sci. Fis. Mat. Natur. Rend. Lincei
  (9) Mat. Appl.}, 9(4):241--243 (1999), 1998.

\bibitem{manz}
O.~Manz and T.~R. Wolf.
\newblock {\em Representations of solvable groups}, volume 185 of {\em London
  Mathematical Society Lecture Note Series}.
\newblock Cambridge University Press, Cambridge, 1993.

\bibitem{mon}
M.~Neunh\"{o}ffer, F.~Noeske, E.~A. O'Brien, and R.~A. Wilson.
\newblock Orbit invariants and an application to the {B}aby {M}onster.
\newblock {\em J. Algebra}, 341:297--305, 2011.

\bibitem{passman}
D.~S. Passman.
\newblock Groups with normal solvable {H}all {$p^{\prime}$}-subgroups.
\newblock {\em Trans. Amer. Math. Soc.}, 123:99--111, 1966.

\bibitem{seress}
A.~Seress.
\newblock The minimal base size of primitive solvable permutation groups.
\newblock {\em J. London Math. Soc. (2)}, 53(2):243--255, 1996.

\bibitem{short}
M.~W. Short.
\newblock {\em The primitive soluble permutation groups of degree less than
  {$256$}}, volume 1519 of {\em Lecture Notes in Mathematics}.
\newblock Springer-Verlag, Berlin, 1992.

\bibitem{singer38}
J.~Singer.
\newblock A theorem in finite projective geometry and some applications to
  number theory.
\newblock {\em Trans. Amer. Math. Soc.}, 43(3):377--385, 1938.

\bibitem{sup}
D.~A. Suprunenko.
\newblock {\em Matrix groups}.
\newblock American Mathematical Society, Providence, R.I., 1976.
\newblock Translated from the Russian, Translation edited by K. A. Hirsch,
  Translations of Mathematical Monographs, Vol. 45.

\bibitem{vdovinReg}
E.~P. Vdovin.
\newblock Regular orbits of solvable linear {$p'$}-groups.
\newblock {\em Sib. \`Elektron. Mat. Izv.}, 4:345--360, 2007.

\bibitem{vd}
E.~P. Vdovin.
\newblock On the base size of a transitive group with solvable point
  stabilizer.
\newblock {\em J. Algebra Appl.}, 11(1):1250015, 14, 2012.
\newblock \url{https://arxiv.org/abs/1011.4341}.

\bibitem{Zen93}
V.~I. Zenkov.
\newblock The structure of intersections of nilpotent {$\pi$}-subgroups in
  {$\pi$}-solvable finite groups.
\newblock {\em Sibirsk. Mat. Zh.}, 34(4):103--107, iii, ix, 1993.

\bibitem{zen}
V.~I. Zenkov.
\newblock Intersections of abelian subgroups in finite groups.
\newblock {\em Mat. Zametki}, 56(2):150--152, 1994.

\bibitem{ZenP}
V.~I. Zenkov.
\newblock Intersections of nilpotent subgroups in finite groups.
\newblock {\em Fundam. Prikl. Mat.}, 2(1):1--92, 1996.

\bibitem{zen21}
V.~I. Zenkov.
\newblock Intersections of three nilpotent subgroups in a finite group.
\newblock {\em Sib Math J}, 62(4):621--637, 2021.

\end{thebibliography}
